\documentclass[11pt,a4paper]{article}
% !TEX root = main.tex
\usepackage{fontspec}
\setmainfont{TeX Gyre Pagella}
\setsansfont{TeX Gyre Heros}
\setmonofont{Latin Modern Mono}
\usepackage[spanish,es-tabla,es-nodecimaldot]{babel}
\usepackage[a4paper,margin=24mm,headheight=15pt]{geometry}
\usepackage{microtype}
\usepackage{array,booktabs,longtable,tabularx}
\usepackage{pdflscape}
\usepackage{xcolor}
\definecolor{accent}{HTML}{164E63}
\usepackage{enumitem}
\setlist{nosep,leftmargin=*}
\usepackage{xurl}
\usepackage{csquotes}
\usepackage[backend=biber,style=numeric,sorting=none,defernumbers=true,maxbibnames=3,giveninits=true,doi=true,url=true,isbn=false]{biblatex}
\addbibresource{bibliography/references.bib}
\usepackage[colorlinks=true,linkcolor=accent,citecolor=accent,urlcolor=accent,unicode=true]{hyperref}
\usepackage{fancyhdr}
\pagestyle{fancy}
\fancyhf{}
\fancyhead[L]{\small Escribas digitales · borrador}
\fancyfoot[C]{\thepage}
\setlength{\parindent}{0pt}
\setlength{\parskip}{5pt}
\setcounter{tocdepth}{1}
\setlength{\emergencystretch}{2em}
\raggedbottom
\AtBeginBibliography{\small\raggedright}
\AtEveryBibitem{\clearfield{addendum}}
\setcounter{biburlnumpenalty}{100}
\setcounter{biburlucpenalty}{100}
\setcounter{biburllcpenalty}{100}
\renewcommand{\arraystretch}{1.15}
\newcolumntype{L}[1]{>{\raggedright\arraybackslash}p{#1}}
\newcommand{\pendiente}[1]{\par\begingroup\small\color{accent}\textbf{Aspectos pendientes de definición:} #1\par\endgroup}
\newcommand{\fuente}[2]{\href{#1}{#2}}

\title{Escribas digitales para la documentación clínica\\\large Evaluación de viabilidad para un hospital con Dedalus HCIS}
\author{Documento de trabajo para Dirección de Innovación}
\date{18 de septiembre de 2026}

\begin{document}
\maketitle
\begin{center}
\small Versión preliminar para evaluación institucional.\\
Mercado: 9 de septiembre de 2026; ampliación AWS: 17 de septiembre de 2026.\\
Bibliografía: revisión inicial del 11 de septiembre; ampliaciones del 12 y 18 de septiembre de 2026.
\end{center}

\begin{abstract}
Los escribas digitales ofrecen una vía para reducir la carga documental mediante la generación de notas a partir de conversaciones clínicas. Este informe examina su posible incorporación a un hospital con Dedalus HCIS y atención en español, a partir de literatura y documentación pública de 43 proveedores. La evidencia muestra beneficios variables y requisitos de seguridad e interoperabilidad pendientes de acreditación. Se presenta una selección corta y una propuesta progresiva de evaluación, cuya arquitectura, servicios participantes, presupuesto y protocolo requieren definición institucional.
\end{abstract}

\begingroup
\small
\setlength{\parskip}{0pt}
\tableofcontents
\endgroup
\clearpage

% !TEX root = ../main.tex
\section{Introducción}
\label{sec:introduccion}
\label{sec:resumen}
Los escribas digitales generan un borrador de la nota clínica a partir de la conversación entre el profesional y el paciente. Su promesa es reducir el trabajo de documentación y permitir que el profesional dedique más atención a la consulta. La nota requiere revisión y aprobación antes de incorporarse a la historia clínica~\cite{7MXFHMBA}.

La rápida expansión del mercado refleja el interés por aliviar esa carga. Los estudios describen mejoras de bienestar y, en algunos contextos, de tiempo de documentación, con resultados variables~\cite{FKHM4UFL,4PZBAAZ7}. El beneficio depende también del trabajo de revisión y del efecto sobre la relación asistencial. Una encuesta de pacientes de UC Davis recoge expectativas de mayor atención y preocupaciones sobre privacidad y exactitud. Al ser previa a la implantación, informa sobre aceptación prevista, no sobre cambios observados en la comunicación~\cite{54NFKG36}.

Este informe compara \textbf{43 proveedores} para un hospital con \textbf{Dedalus HCIS y atención en español}. Las prioridades son protección de datos, ciberseguridad, calidad, integración y coste total. ENS de categoría ALTA se mantiene como criterio prioritario, sujeto a definición institucional. La documentación aún no permite confirmar un candidato con integración ambiental operativa con HCIS y alcance ENS íntegramente verificado (sección~\ref{sec:candidatos}).

Se propone avanzar de forma progresiva: definir condiciones de uso, valorar un conjunto local de consultas de prueba y comparar entre uno y tres proveedores antes de ampliar el uso asistencial. La evaluación deberá comprobar el beneficio en español y en el flujo completo hasta HCIS, incluyendo la revisión de las notas. Servicios participantes, política de datos, presupuesto, protocolo y umbrales de aceptación permanecen abiertos.

% !TEX root = ../main.tex
\section{Motivación y beneficios esperados}
\label{sec:motivacion}
El escriba clínico humano registra información del encuentro asistencial para facilitar su incorporación a la historia clínica electrónica. Los escribas digitales buscan automatizar parte de esa tarea mediante la captura de la conversación y la generación de un borrador. Su propuesta responde a una carga documental relevante: un estudio de observación de 57 médicos de cuatro especialidades en Estados Unidos atribuyó el 49,2\% de la jornada en consulta al trabajo con la historia electrónica y de escritorio~\cite{7MXFHMBA,sinsky2016}.

\textbf{Adopción y mercado.} PHTI describió en 2025 una expansión excepcionalmente rápida, basada en la experiencia de directivos y empresas del sector. Esta valoración respalda el interés por la tecnología, pero no establece un récord histórico mundial de adopción. La revisión de mercado del presente informe reúne 43 proveedores con diferentes modelos comerciales y grados de integración; su amplitud justifica una selección orientada a los requisitos locales~\cite{FKHM4UFL}.

\textbf{Beneficio esperado.} Una revisión de 35 estudios identifica resultados generalmente favorables o mixtos sobre bienestar, con predominio de investigaciones estadounidenses y medidas percibidas~\cite{4PZBAAZ7}. En el ensayo aleatorizado de UCLA, Nabla redujo el tiempo de documentación un 9,5\% frente al control; DAX no mostró una diferencia significativa~\cite{PR93N8FN}. Por tanto, el ahorro debe estimarse localmente, sin adoptar un porcentaje único como expectativa de todos los productos.

\textbf{Incertidumbres.} Persisten preguntas sobre errores, variación entre especialidades, efectos sobre el razonamiento profesional y cambios acumulados en la historia clínica. La posible dependencia cognitiva es una hipótesis de investigación, no un perjuicio demostrado de estos sistemas~\cite{ZI87XIWC}. El objetivo del proyecto es comprobar si el alivio documental se traduce en una mejora útil y sostenible para el hospital.

% !TEX root = ../main.tex
\section{Alcance de la implantación}
\label{sec:flujo}
El alcance inicial propuesto es \textbf{conversación en español, nota clínica editable e incorporación a HCIS tras revisión profesional}. La documentación de referencia distingue la captura, la generación y la validación del resultado; las funciones concretas dependen del producto y de su integración~\cite{7MXFHMBA,nhsAmbient2025}.

\begin{center}
\small
\fbox{Captura de audio} $\longrightarrow$ \fbox{Borrador clínico}
$\longrightarrow$ \fbox{Revisión y aprobación}\\[6pt]
$\longrightarrow$ \fbox{Incorporación al paciente y episodio correctos en HCIS}
\end{center}

\textbf{Captura y procesamiento.} La oferta revisada incluye micrófonos conectados a un ordenador, dispositivos dedicados y aplicaciones móviles, web o extensiones de navegador. Predominan los servicios con procesamiento remoto, aunque las condiciones regionales y la retención varían (apéndice~\ref{app:proveedores}). Un dispositivo autónomo de captura no implica procesamiento local: algunos conservan audio hasta recuperar la conexión. La compatibilidad con equipos, navegadores y red debe comprobarse en los servicios participantes.

\textbf{Revisión e incorporación.} La nota puede editarse en la plataforma del proveedor o en una interfaz integrada. El flujo exigido deberá identificar paciente y episodio, permitir correcciones y registrar su incorporación a HCIS. Aprobar el borrador, transferirlo y firmarlo en la historia clínica pueden ser operaciones diferentes; no se presupone una función común de firma ni un conector disponible en todos los productos. El apéndice~\ref{app:hcis} distingue integración publicada, dictado, transferencia manual y desarrollo pendiente.

Las plantillas se incluyen cuando determinan el contenido y la estructura de la nota. La codificación automática, las recomendaciones clínicas y otros documentos quedan fuera del alcance inicial y requerirían evaluación específica. Para sistemas remotos, una conexión insuficiente puede retrasar la generación; el piloto deberá contemplar la documentación habitual como alternativa ante fallos.

\textbf{Funciones y alcance regulatorio.} Transcribir lo expresado por el médico y generar una recomendación de tratamiento son funciones distintas. La TGA australiana ilustra esta diferencia: la transcripción sin análisis queda fuera de su definición de producto sanitario, mientras que generar diagnósticos o recomendaciones no expresados por el profesional entra en ese ámbito~\cite{tgaScribes2026}. Para España, la finalidad y las funciones deben examinarse conforme al marco europeo; la guía australiana aporta contexto, sin sustituirlo~\cite{mdcg2019rev2025}.

% !TEX root = ../main.tex
\clearpage
\section{Criterios de adecuación al hospital}
\label{sec:evaluacion}
La revisión documental delimita candidatos y preguntas pendientes. Los criterios siguientes orientan la selección; \textbf{el diseño del piloto, las métricas y los umbrales de aceptación siguen pendientes de definición}.

\subsection{Protección de datos, seguridad e integración}
Se propone priorizar \textbf{almacenamiento y procesamiento en la UE/EEE, ENS de categoría ALTA e integración con Dedalus HCIS}. La ubicación europea no acredita por sí sola cumplimiento del RGPD: deben verificarse finalidades, responsabilidades, contratos, subencargados y transferencias~\cite{rgpd2016}. Para ENS deben comprobarse titular, sistema, alcance, emisor y vigencia, así como la categoría exigible al sistema concreto~\cite{ens2022}.

La arquitectura dependerá del tratamiento de datos admisible y la capacidad operativa del hospital. Las certificaciones no sustituyen la evaluación clínica; una alianza con Dedalus o una API tampoco acredita integración con la instalación local de HCIS (apéndices~\ref{app:certificaciones} y~\ref{app:hcis}).

\subsection{Factores de comparación}
\begin{small}
\begin{longtable}{L{.22\linewidth}L{.70\linewidth}}
\toprule
Factor & Evidencia relevante para la decisión \\
\midrule\endhead
Seguridad y datos & Certificados e informes completos; alojamiento y procesamiento; acceso de soporte; retención de audio y notas; subencargados y reutilización para entrenamiento. \\
Calidad en español & Evaluaciones en el idioma y contexto previstos: omisiones, adiciones, negaciones, fármacos, dosis y atribución de hablantes. Resultados por especialidad y condiciones de captura. \\
Utilidad clínica & Tiempo total de documentación, carga de revisión, experiencia del profesional y del paciente, utilización y abandono. \\
Integración y operación & Identificación de paciente y episodio, edición, transferencia, firma, registro de errores, dispositivos compatibles y comportamiento ante pérdida de conexión. \\
Coste total & Licencias, consumo, equipos, integración, mantenimiento, formación, soporte, mínimos contractuales y condiciones de salida. \\
\bottomrule
\end{longtable}
\end{small}
La matriz combina las necesidades del proyecto con dominios de estudios comparativos~\cite{NNWFT2JX,PRTISMYF}. Declarar capacidad multilingüe no acredita calidad equivalente: interesa evaluar español, acentos, expresiones coloquiales y varios hablantes. Las simulaciones disponibles orientan las pruebas, sin sustituir la validación local.

\subsection{Interpretación de la evidencia}
Se distinguen documentos examinados, declaraciones comerciales, evidencia de uso y aspectos pendientes. Revisión por pares, independencia del proveedor y aplicabilidad local se valoran por separado. Los estudios de implantación, los ensayos comparativos y las evaluaciones de errores responden a preguntas diferentes; sus diseños y medidas desaconsejan una puntuación global de superioridad clínica~\cite{956B6PZT}.

La similitud textual no equivale a exactitud clínica ni la nota del médico constituye una referencia infalible. La evaluación distinguirá calidad del borrador, correcciones y calidad de la nota final~\cite{4IHZ7QFQ,ZI87XIWC}.

% !TEX root = ../main.tex
\clearpage
\section{Opciones y selección corta}
\label{sec:candidatos}
La muestra comprende \textbf{43 proveedores}: 20 del conjunto inicial, 22 adicionales del registro NHS y AWS HealthScribe, incorporado el 17 de septiembre de 2026. La inclusión en dicho registro responde a un procedimiento de autocertificación y no constituye una aprobación clínica general~\cite{EK6X6DWK}. Los apéndices presentan el inventario completo y sus fuentes. La preselección responde a los requisitos del hospital: procesamiento remoto admisible, ENS ALTA, atención en español e integración ambiental con HCIS. No establece una clasificación de calidad clínica del mercado.

\subsection{Certificación ENS, integración con HCIS y precio}
\textbf{Tandem, Invox Genesis y Solventum Align} se priorizan para solicitar documentación y demostraciones por la evidencia o declaraciones ENS identificadas y su presencia comercial o asistencial en España. Ninguno reúne todavía todos los requisitos acreditados. La tabla resume ENS, integración y precio; los párrafos siguientes completan los factores de la sección~\ref{sec:evaluacion} con datos, evidencia clínica y español. La preselección documental no equivale a una recomendación de compra.
\begin{small}
\begin{longtable}{L{.14\linewidth}L{.25\linewidth}L{.23\linewidth}L{.26\linewidth}}
\toprule
Candidato & ENS ALTA & Integración ambiental con HCIS & Precio y forma de cobro \\
\midrule\endhead
Tandem & Certificado Bureau Veritas examinado; copia pública parcial (página 1 de 2). & El proveedor indica «no disponible». & Pro: 125 €/usuario/mes, facturado anualmente. Enterprise: por usuario, importe a medida. \\
Invox Genesis & Obtención anunciada; certificado y alcance Genesis no verificados. & Dictation anuncia HCIS; Genesis no confirmado. & Presupuesto comercial. Importe, unidad facturable y mínimos no publicados en las fuentes consultadas. \\
Solventum Align & Cumplimiento declarado; verificación del certificado pendiente. & No localizado para Align. & Presupuesto comercial. Importe, volumen incluido y coste de implantación no publicados en las fuentes consultadas. \\
\bottomrule
\end{longtable}
\end{small}
Fuentes: \fuente{https://trust.tandemhealth.ai/}{expediente Tandem}, \fuente{https://tandemhealth.ai/es/producto/integraciones/hp-his-hp-hcis}{estado HCIS Tandem}, \fuente{https://www.invoxmedical.com/blog/invox-medical-refuerza-su-compromiso-con-la-seguridad-de-la-informacion-con-la-obtencion-de-la-certificacion-iso-27001-y-el-ens-nivel-alto}{ENS Invox}, \fuente{https://dictation.invoxmedical.com/software-medico/dedalus-hcis/}{Dictation--HCIS} y \fuente{https://www.solventum.com/es-es/home/health-information-technology/solutions/fluency-align/}{Align}. Documentación adicional en apéndices~\ref{app:proveedores} y~\ref{app:hcis}. «No localizado» no significa imposibilidad técnica.

\textbf{Lectura económica.} Los importes proceden del corte de mercado del 9 de septiembre de 2026; las fuentes y condiciones están en el apéndice~\ref{app:precios}. La tarifa Pro de Tandem permite cuantificar una licencia, pero no el despliegue hospitalario con HCIS. En los tres casos falta un presupuesto comparable de licencias, consumo, conector, equipos, formación, soporte e impuestos. Una tarifa más transparente facilita presupuestar; no demuestra menor coste total.

\subsection{Aspectos que condicionan la preselección}
\textbf{Tandem.} En seguridad y datos, aporta documentos externos examinables, declara tratamiento UE y publica subencargados; el dictamen C5 tiene salvedades y el alcance ENS público debe completarse (apéndice~\ref{app:proveedores}). En utilidad clínica, su estudio sueco es observacional y selecciona usuarios con utilización sostenida. El ahorro de tiempo es autodeclarado y no acredita rendimiento en español~\cite{C7S8N9NA}.

\textbf{Invox Genesis.} En seguridad y datos, quedan por verificar el certificado ENS, el país de procesamiento y la cadena contractual. En experiencia asistencial, la Comunidad de Madrid confirma su uso en Psiquiatría del Hospital Fundación Alcorcón. Es una referencia de implantación en España, sin resultados cuantificados que establezcan calidad clínica en español o ahorro de trabajo. Tampoco acredita un despliegue general del SERMAS ni una integración Genesis--HCIS. \fuente{https://www.comunidad.madrid/noticias/2025/06/10/hospital-publico-fundacion-alcorcon-incorpora-aplicacion-ia-optimizar-generacion-documentacion-clinica-partir-entrevistas}{Fuente hospitalaria}; documentación en el apéndice~\ref{app:proveedores}.

\textbf{Solventum Align.} En seguridad y datos, declara almacenamiento y procesamiento UE, con certificado ENS íntegro y condiciones completas de tratamiento pendientes de verificación. En calidad y utilidad clínica, no se identificó un estudio revisado por pares específico de Align, por lo que su oferta en España no basta para acreditar rendimiento en español. Los casos de Fluency Direct corresponden a otro producto (apéndice~\ref{app:proveedores}).

La exigencia de infraestructura propia, la aceptación de dictado o el desarrollo de un conector obligarían a revisar esta selección.

\subsection{Otros proveedores con evidencia relevante para la evaluación}
Heidi, Nabla, Abridge y Ambience permiten examinar preguntas complementarias. Su literatura justifica incluirlos en el análisis, pero no acredita los requisitos ENS ALTA y HCIS del proyecto. Tampoco permite ordenar estos productos por eficacia a partir de estudios con diseños y contextos diferentes (apéndices~\ref{app:proveedores} y~\ref{app:hcis}).

\textbf{Heidi: evidencia preliminar en español.} Un preprint con simulaciones en cinco países incluye Barcelona y consultas en español. Ofrece un motivo concreto para examinar el producto en pruebas locales, pero no lo convierte en un patrón de calidad lingüística ni demuestra su rendimiento asistencial. El estudio está financiado por Heidi y carece de revisión por pares~\cite{heidiSpanish2026}.

\textbf{Nabla: comparación aleatorizada del trabajo documental.} El ensayo de UCLA compara Nabla, DAX y documentación habitual. Informa una reducción modesta del tiempo con Nabla y aporta un diseño útil para estimar el efecto frente al flujo habitual. Se realizó en un centro y en inglés; no compara Nabla con Abridge o Ambience ni demuestra superioridad sobre ellos~\cite{PR93N8FN}.

\textbf{Abridge: bienestar en varios sistemas sanitarios.} Un estudio de 263 profesionales en seis sistemas estadounidenses compara encuestas antes y después de 30 días de uso. Describe mejoras percibidas de carga documental y bienestar, con participación voluntaria, sin aleatorización y con coautoría del proveedor~\cite{olson2025}. Complementa el ensayo de UCLA con una experiencia multicéntrica, aunque responde a otra pregunta. Su catálogo incluye también una \fuente{https://doi.org/10.1056/AIcs2400977}{publicación sobre aseguramiento de calidad}, recogida en el apéndice~\ref{app:proveedores}.

\textbf{Ambience: organización de la implantación.} La experiencia de Cleveland Clinic describe la incorporación de más de 4.000 profesionales en cuatro meses y las medidas de gobernanza, formación y soporte. Su principal aportación al proyecto es organizativa; el informe descriptivo, con coautoría del proveedor, no establece una ventaja clínica frente a otros productos~\cite{F5FVUMMX}.

% !TEX root = ../main.tex
\clearpage
\section{Lecciones de despliegues y propuesta de piloto}
\label{sec:piloto}
\subsection{Resultados y límites de las experiencias previas}
\textbf{El alivio documental no implica necesariamente consultas más breves.} En Suecia se comunica ahorro de tiempo percibido, mientras que en Quirónsalud las mejoras de experiencia coexisten con una duración media de consulta ligeramente mayor entre los usos de Scribe. El ensayo de UCLA muestra además que el resultado puede diferir entre productos~\cite{C7S8N9NA,U5GHQW86,PR93N8FN}. Para el hospital, la consecuencia es evaluar por separado carga percibida, tiempo total de documentación y calidad clínica. La experiencia de Cleveland Clinic subraya también el papel de la formación y el soporte en la adopción~\cite{F5FVUMMX}. Permanente añade una advertencia: el uso se concentra en una parte de los profesionales~\cite{tierney2025}. Los casos siguientes orientan esas decisiones; no constituyen una comparación directa entre proveedores.

\begin{small}
\begin{longtable}{L{.23\linewidth}L{.69\linewidth}}
\toprule
Experiencia & Resultado y consecuencia para el proyecto \\
\midrule\endhead
UCLA, Estados Unidos & Ensayo aleatorizado con 238 médicos: reducción del tiempo de documentación con Nabla, sin diferencia significativa con DAX. Cerca del 15\% de los médicos asignados a escribas no los utilizó. Conviene incluir a todos los participantes y registrar el abandono~\cite{PR93N8FN}. \\
The Permanente Medical Group, Estados Unidos & Seguimiento de 63 semanas: 7.260 médicos y 2,58 millones de encuentros; ahorro agregado anual estimado de 15.791 horas, comparando cambios frente a no usuarios. Diseño observacional, sin asignación aleatoria. La escala permite un ahorro acumulado importante, pero no demuestra un beneficio uniforme. Conviene analizar utilización y carga inicial por profesional~\cite{tierney2025}. \\
Abridge, seis sistemas de Estados Unidos & Estudio antes/después de 263 profesionales: mejoras autodeclaradas de bienestar y carga documental tras 30 días. La participación voluntaria y la ausencia de un control contemporáneo limitan la atribución causal. Conviene medir bienestar sin equipararlo a ahorro objetivo~\cite{olson2025}. \\
Cleveland Clinic, Estados Unidos & Informe de implantación de Ambience: más de 4.000 profesionales incorporados en cuatro meses, con gobernanza, formación y soporte. La coautoría del proveedor y el diseño descriptivo limitan la interpretación clínica. El caso aporta un modelo organizativo~\cite{F5FVUMMX}. \\
Tandem, Suecia & Estudio observacional de 236.153 notas; encuesta de 177 usuarios con utilización sostenida. Reducción autodeclarada del 29\% en tiempo por nota. La selección de usuarios y la edición fuera de la aplicación pueden sesgar la estimación del beneficio~\cite{C7S8N9NA}. \\
Quirónsalud, España & Más de 2,33 millones de usos ambulatorios, con mejoras percibidas; duración media de 15,01 minutos con Scribe y 14,65 sin él. El acuerdo semántico de 87,4--89,2\% compara borrador y texto editado, sin establecer exactitud clínica independiente. Satisfacción y duración requieren análisis separado~\cite{U5GHQW86}. \\
GOSH y otros centros NHS, Londres & Evaluación de Tortus en nueve centros y más de 17.000 encuentros. El hospital comunica menor carga documental y mayor tiempo de interacción. Se trata de un estudio antes/después no aleatorizado, difundido como informe y preprint; los efectos no se extrapolan automáticamente al español~\cite{gosh2025,wray2025}. \\
\bottomrule
\end{longtable}
\end{small}
En España, el Hospital Fundación Alcorcón aporta además una referencia institucional de uso de Invox Genesis en Psiquiatría, descrita en la sección~\ref{sec:candidatos}. Esa adopción demuestra experiencia de implantación, pero no sustituye una evaluación comparativa del rendimiento.

\subsection{Secuencia de trabajo propuesta}
Se propone un proceso progresivo, sujeto a disponibilidad presupuestaria y a la definición del hospital. El apéndice~\ref{app:evaluacion-piloto} desarrolla opciones de diseño y medidas de resultado para su revisión. Es una base de discusión; no fija todavía el protocolo, la muestra ni los umbrales de aceptación.
\begin{enumerate}
\item \textbf{Preparación local.} Seleccionar servicios, profesionales interesados y condiciones técnicas; registrar la situación inicial de documentación y resolver el tratamiento de datos permitido.
\item \textbf{Conjunto común de consultas de prueba.} Valorar la recogida de audio y notas con finalidad y condiciones específicas, sin que ello autorice su envío a proveedores. Si no es viable utilizar consultas reales, pueden emplearse escenarios simulados~\cite{rgpd2016,NNWFT2JX}.
\item \textbf{Comparación previa al uso asistencial.} Probar entre uno y tres proveedores sobre casos comunes, según coste y requisitos, con revisión clínica de los borradores y sin incorporarlos a HCIS como documentación asistencial.
\item \textbf{Piloto asistencial limitado.} Evaluar los candidatos que superen la fase previa, con revisión profesional y alternativa ante fallos. La ampliación dependerá de criterios acordados antes del inicio.
\end{enumerate}

La asignación debería equilibrar especialidades, complejidad y experiencia digital; asignar un producto a cada servicio confundiría ambos efectos. Los profesionales interesados pueden facilitar el inicio, pero deberán registrarse también la falta de uso, los rechazos y el abandono~\cite{PR93N8FN,FKHM4UFL}.

El banco de consultas permite comparar notas; para estimar ahorro asistencial será necesario comparar flujos de trabajo. Las alternativas de diseño, la referencia clínica y las medidas propuestas se desarrollan en el apéndice~\ref{app:evaluacion-piloto}.

% !TEX root = ../main.tex
\section{Decisiones solicitadas}
\label{sec:decision}
La fase preparatoria requiere definir cinco aspectos antes de seleccionar proveedor y concretar el piloto:
\begin{enumerate}
\item \textbf{Política de datos y nivel de aseguramiento.} Determinar la arquitectura admisible, el alcance ENS y la documentación exigida. El alojamiento local también requiere seguridad, mantenimiento y evaluación clínica.
\item \textbf{Servicios y participantes.} Identificar áreas prioritarias, tipos de consulta, profesionales y responsables clínico y técnico. Esta selección orientará los requisitos lingüísticos y de captura.
\item \textbf{Preparación del conjunto de evaluación.} Resolver si procede grabar consultas reales. Protección de Datos y el comité competente, cuando corresponda, deberán definir finalidad, base jurídica, información, consentimiento exigible, acceso, conservación y reutilización~\cite{rgpd2016}.
\item \textbf{Presupuesto y alcance comercial.} Fijar un límite para proveedores, equipos, integración, soporte y evaluación clínica. Las tarifas individuales no equivalen a un presupuesto hospitalario (apéndice~\ref{app:precios}).
\item \textbf{Infraestructura y trazabilidad.} Precisar equipos, sistemas, navegadores, periféricos y calidad de red. Decidir si basta con audio y nota o se requieren transcripción, versiones, correcciones y registros de transferencia y error, con finalidad y conservación definidas.
\end{enumerate}

Estas definiciones permitirán solicitar propuestas comparables y elaborar el protocolo. El siguiente hito será decidir si procede el piloto y en qué condiciones; la preparación no implica adjudicación ni autorización de uso asistencial.

% !TEX root = ../main.tex
\section*{Declaración sobre el uso de inteligencia artificial}
\label{sec:uso_ia}
Se utilizaron Codex y ChatGPT (OpenAI) como apoyo parcial en la búsqueda de referencias, la estructuración y revisión de la redacción a partir de notas y textos del autor y la organización bibliográfica. La valoración crítica de las fuentes, las decisiones de contenido y la responsabilidad sobre la versión final corresponden al autor.


\clearpage
\printbibliography[heading=bibintoc,notkeyword=tecnica,title={Referencias bibliográficas}]

\clearpage
\appendix
% !TEX root = ../main.tex
% GENERADO por scripts/sync_market_appendices.py; editar el Markdown de origen.
\clearpage
\section{Certificaciones y marcos de aseguramiento: alcance y evidencia}
\label{app:certificaciones}
\begingroup\footnotesize

Revisión pública inicial: \textbf{9 de septiembre de 2026}. \textbf{43 proveedores}: 20 del conjunto inicial, 22 adicionales del registro NHS y AWS HealthScribe, incorporado el \textbf{17/\allowbreak{}09/\allowbreak{}2026}. Esta ampliación no supone una nueva revisión de las otras 42 fichas. La comparación aplica un formato y unos criterios de evidencia uniformes; el número de credenciales no constituye una medida de calidad clínica. \hyperref[app:hcis]{Integraciones con HCIS} · \hyperref[app:precios]{Precios} · \hyperref[sec:candidatos]{Evaluación comparativa}.

\textbf{Criterios de clasificación de la evidencia:} «declara» significa que el proveedor anuncia la credencial; la categoría «documentación primaria examinada» identifica certificados, informes o registros cuyo contenido ha sido consultado directamente. La verificación íntegra del emisor, la vigencia y el alcance requiere el certificado o informe completo. «No localizado» indica ausencia de evidencia identificada en las fuentes públicas consultadas; no permite concluir que la credencial no exista. Las condiciones de datos son declaraciones/\allowbreak{}documentos del proveedor, salvo indicación contraria.

\subsection*{Definición y alcance de las credenciales}

\begin{longtable}{L{\dimexpr 0.23\linewidth-2\tabcolsep\relax}L{\dimexpr 0.77\linewidth-2\tabcolsep\relax}}
\toprule
\textbf{Credencial} & \textbf{Objeto y límites de acreditación} \\
\midrule
\endfirsthead
\textbf{Credencial} & \textbf{Objeto y límites de acreditación} \\
\midrule
\endhead
\midrule
\multicolumn{2}{r}{\footnotesize Continúa en la página siguiente}
\endfoot
\bottomrule
\endlastfoot
\textbf{ENS, categoría ALTA: España} & Esquema Nacional de Seguridad: requisitos organizativos y técnicos para sistemas del sector público y proveedores incluidos en su ámbito. \textbf{ALTA es la categoría más exigente}, sujeta a auditoría; importa el sistema incluido en el certificado. No acredita por sí sola exactitud clínica ni residencia UE. \href{https://www.boe.es/eli/es/rd/2022/05/03/311}{BOE, RD 311/\allowbreak{}2022}. \\
\textbf{ISO/\allowbreak{}IEC 27001} & Certifica el sistema de gestión de seguridad de la información: riesgos, accesos, incidentes, continuidad y protección de confidencialidad, integridad y disponibilidad. El alcance puede cubrir solo parte de la empresa. \href{https://www.iso.org/standard/27001}{ISO}. \\
\textbf{ISO/\allowbreak{}IEC 27701} & Sistema de gestión de privacidad: responsabilidades y controles sobre el tratamiento de datos personales. No sustituye el análisis de cumplimiento del RGPD. \href{https://www.iso.org/standard/27701}{ISO}. \\
\textbf{ISO/\allowbreak{}IEC 27017 /\allowbreak{} 27018} & \textbf{27017:} controles de seguridad específicos de servicios cloud. \textbf{27018:} protección de datos personales cuando la nube pública actúa como encargada del tratamiento. \href{https://www.iso.org/standard/43757.html}{27017}, \href{https://www.iso.org/standard/76559.html}{27018}. \\
\textbf{ISO/\allowbreak{}IEC 20000-1} & Sistema de gestión de servicios de TI: planificación, prestación y mejora del servicio. No evalúa exactitud clínica. \href{https://www.iso.org/standard/70636.html}{ISO}. \\
\textbf{ISO/\allowbreak{}IEC 42001} & Certifica la \textbf{gestión de IA}: gobernanza, riesgos, responsabilidades y mejora continua. No acredita la ausencia de contenido no sustentado ni certifica la precisión clínica del escriba. \href{https://www.iso.org/standard/42001}{ISO}. \\
\textbf{Ethik-IA: Francia} & Sello privado de supervisión humana y uso ético de IA sanitaria. Examina, entre otros aspectos, equidad, trazabilidad, datos y robustez; no es el marcado CE ni una aprobación del AI Act. \href{https://esante.gouv.fr/sites/default/files/media\_entity/documents/guide-ia\_vf.pdf}{Guía oficial francesa, §5.6}. \\
\textbf{ISO 9001 /\allowbreak{} 13485 /\allowbreak{} 14001} & \textbf{9001:} gestión general de calidad. \textbf{13485:} gestión de calidad para productos sanitarios y requisitos regulatorios. \textbf{14001:} gestión ambiental. Ninguna sustituye la evaluación clínica del producto. \href{https://www.iso.org/standard/62085.html}{9001}, \href{https://www.iso.org/standard/59752.html}{13485}, \href{https://www.iso.org/standard/14001}{14001}. \\
\textbf{SOC 1 /\allowbreak{} SOC 2 /\allowbreak{} SOC 3; tipos I y II} & Informes de auditoría: \textbf{SOC 1}, controles relevantes para información financiera; \textbf{SOC 2}, controles de seguridad y, según alcance, disponibilidad, confidencialidad, privacidad e integridad del procesamiento. \textbf{I:} diseño en una fecha; \textbf{II:} también funcionamiento durante un periodo. \textbf{SOC 3:} informe de uso general sobre controles, con menos detalle que SOC 2. \href{https://www.aicpa-cima.com/resources/landing/system-and-organization-controls-soc-suite-of-services}{AICPA}. \\
\textbf{C5 tipo II} & Auditoría de controles de seguridad cloud según el catálogo del BSI alemán; tipo II incluye eficacia durante un periodo. El auditor emite el informe, no un certificado del BSI. \href{https://www.bsi.bund.de/SharedDocs/Downloads/EN/BSI/Publications/CloudComputing/ComplianceControlsCatalogue-Cloud\_Computing-C5.pdf?\_\_blob=publicationFile\&v=3}{BSI}. \\
\textbf{ISAE 3000} & Norma para encargos de aseguramiento no financiero. Puede respaldar una auditoría de privacidad, pero su interpretación requiere examinar criterios, alcance y periodo: el nombre solo no acredita cumplimiento íntegro del RGPD. \href{https://www.iaasb.org/publications/international-standard-assurance-engagements-isae-3000-revised-assurance-engagements-other-audits-or}{IAASB}. \\
\textbf{NEN 7510 /\allowbreak{} 7512 /\allowbreak{} 7513: Países Bajos} & \textbf{7510:} seguridad de información sanitaria; \textbf{7512:} intercambio seguro; \textbf{7513:} registro de accesos y operaciones. Aplicar las normas no equivale automáticamente a poseer certificados externos de las tres. \href{https://www.nen.nl/en/zorg-welzijn/ict-in-de-zorg/informatiebeveiliging-in-de-zorg}{NEN}. \\
\textbf{Cyber Essentials /\allowbreak{} Plus: UK} & Protección frente a ciberataques comunes. La modalidad básica usa autoevaluación verificada; \textbf{Plus añade comprobación técnica independiente}. No evalúa calidad clínica. \href{https://www.ncsc.gov.uk/cyberessentials/overview}{NCSC}. \\
\textbf{HDS: Francia} & Certificación del alojamiento de datos de salud por un organismo acreditado. Su interpretación requiere comprobar el titular y las actividades cubiertas: usar un proveedor HDS no convierte automáticamente al escriba en certificado HDS. \href{https://esante.gouv.fr/produits-services/hds}{Agencia de Salud Digital}. \\
\textbf{HITRUST} & Certificaciones de controles de ciberseguridad con distintas modalidades y profundidad. Deben identificarse certificado, sistema y vigencia; mencionar el marco no demuestra estar certificado. \href{https://hitrustalliance.net/assessments-and-certifications}{HITRUST}. \\
\textbf{TX-RAMP: Texas} & Programa público de evaluación de seguridad cloud. \textbf{Nivel 2:} datos confidenciales en sistemas de impacto moderado/\allowbreak{}alto; existe también nivel 1 y estatus provisional. No es una aprobación clínica. \href{https://dir.texas.gov/information-security/texas-risk-and-authorization-management-program-tx-ramp/tx-ramp-frequently}{Texas DIR}. \\
\textbf{CE bajo MDR /\allowbreak{} UKCA médico; clases I y IIa} & Marcados de conformidad para productos sanitarios de la UE/\allowbreak{}Gran Bretaña. En software de \textbf{clase I} suele bastar la declaración del fabricante; \textbf{IIa requiere evaluación de un organismo externo}. Alcance e indicación de uso son esenciales. \href{https://eur-lex.europa.eu/eli/reg/2017/745/oj}{MDR UE}, \href{https://www.gov.uk/guidance/regulating-medical-devices-in-the-uk}{MHRA}. \\
\textbf{Registro MHRA} & Inscripción para comercializar en Gran Bretaña; no equivale por sí misma a una evaluación clínica independiente o certificación de clase IIa. \href{https://www.gov.uk/guidance/regulating-medical-devices-in-the-uk}{MHRA}. \\
\textbf{Certificado de libre venta} & Acredita la comercialización de los productos identificados y facilita su exportación. Es distinto del certificado de conformidad MDR y de un estudio clínico. \href{https://www.lakemedelsverket.se/en/medical-devices/sale}{Agencia sueca del medicamento}. \\
\textbf{LVV/\allowbreak{}Valvira A1: Finlandia} & Clasificación de sistemas de información sociosanitaria, con evaluación de seguridad recogida en Astori. \textbf{A1 no es una clase de producto sanitario MDR}; son marcos diferentes. \href{https://astori.lvv.fi/}{Registro LVV}. \\
\textbf{ISO 14971 /\allowbreak{} IEC 62366-1} & \textbf{14971:} gestión de riesgos de productos sanitarios. \textbf{62366-1:} usabilidad y prevención de errores de uso que afectan a seguridad. En las fichas son normas aplicadas, salvo certificado externo identificado. \href{https://www.iso.org/standard/72704.html}{ISO}, \href{https://webstore.iec.ch/en/publication/21863}{IEC}. \\
\textbf{DTAC /\allowbreak{} DCB0129 /\allowbreak{} DCB0160 /\allowbreak{} DSPT: NHS} & \textbf{DTAC:} criterios para evaluar tecnología digital al adquirirla. \textbf{DCB0129:} gestión de riesgos clínicos por el fabricante; \textbf{DCB0160}, por la organización que lo implanta. \textbf{DSPT:} autoevaluación de seguridad/\allowbreak{}protección de datos. «NHS compliant» no constituye un certificado externo único. \href{https://digital.nhs.uk/services/digital-technology-assessment-criteria-dtac}{DTAC}, \href{https://digital.nhs.uk/data-and-information/information-standards/governance/latest-activity/standards-and-collections/dcb0129-clinical-risk-management-its-application-in-the-manufacture-of-health-it-systems/}{DCB0129}, \href{https://www.dsptoolkit.nhs.uk/}{DSPT}. \\
\textbf{TGV /\allowbreak{} EFVP: Québec} & \textbf{TGV} (\emph{Trousse globale de vérification}): certificación del producto exigida por el sistema sanitario provincial. \textbf{EFVP:} evaluación del impacto en privacidad de su implantación, responsabilidad de la organización compradora; son requisitos distintos. \href{https://www.quebec.ca/sante/professionnels/programme-transcription-intelligence-artificielle}{Gobierno de Québec}. \\
\textbf{FCC /\allowbreak{} ISED /\allowbreak{} conformidad radio UE-UK; ensayo SGS} & En Heidi Remote, credenciales de radio del \textbf{hardware} y un ensayo antibacteriano concreto de SGS. No certifican el escriba ni convierten el micrófono en producto sanitario. \href{https://www.heidihealth.com/en-ca/docs/remote-user-manual}{Manual y referencias regulatorias}. \\
\textbf{WEEE /\allowbreak{} RoHS /\allowbreak{} EMC: hardware} & \textbf{WEEE:} gestión de residuos electrónicos; aquí, adhesión a un sistema de recogida/\allowbreak{}reciclaje. \textbf{RoHS:} restricción de sustancias peligrosas. \textbf{EMC:} compatibilidad electromagnética. Las dos últimas figuran en la declaración del fabricante del micrófono Tandem. \href{https://environment.ec.europa.eu/topics/waste-and-recycling/rohs-directive\_en}{WEEE/\allowbreak{}RoHS}, \href{https://trust.tandemhealth.ai/doc/trust?rid=6a78f0c64c2eeb0a116d6fb1\&r=e0qkrepqo674el4jxq28px\&view=true}{declaración}. \\
\textbf{AAPC /\allowbreak{} ACR-VPAT} & \textbf{AAPC:} aquí se refiere a una evaluación de codificación anunciada por Ambience. \textbf{ACR/\allowbreak{}VPAT:} informe sobre accesibilidad de software, publicado por Oracle. No son certificaciones generales de seguridad clínica. \href{https://www.ambiencehealthcare.com/}{Ambience}, \href{https://docs.oracle.com/en/corporate/accessibility/templates/t2-16521.html}{Oracle}. \\
\textbf{ISO 22301} & Gestión de continuidad de negocio: preparación, respuesta y recuperación ante interrupciones. No acredita exactitud clínica. \href{https://www.iso.org/standard/75106.html}{ISO}. \\
\textbf{Dubai AI Seal} & DCAI evalúa y clasifica empresas proveedoras de IA; el sello tiene número y categoría. No es aprobación de un producto sanitario. \href{https://dub.ai/en/ai-seal/}{DCAI}. \\
\textbf{CSA STAR niveles 1–2 /\allowbreak{} GovRAMP /\allowbreak{} IRAP} & \textbf{STAR 1:} autoevaluación de seguridad cloud. \textbf{STAR 2:} certificación o atestación con evaluación independiente; es necesario identificar modalidad y alcance. \textbf{GovRAMP:} programa de seguridad cloud con distintos niveles de verificación; es necesario verificar el estado específico. \textbf{IRAP:} evaluación frente a requisitos australianos, no certificación o acreditación del sistema. \href{https://cloudsecurityalliance.org/star/}{CSA}, \href{https://govramp.org/security-program}{GovRAMP}, \href{https://www.cyber.gov.au/business-government/protecting-devices-systems/assessment-evaluation-programs/irap}{IRAP}. \\
\textbf{IM1: NHS} & Proceso de aseguramiento de interfaces con sistemas de atención primaria, actualmente EMIS Web y SystmOne. No evalúa los modelos de IA ni acredita integración con HCIS. \href{https://digital.nhs.uk/services/digital-services-for-integrated-care/im1-pairing-integration}{NHS}. \\
\textbf{CREST: pentests} & Acredita empresas de pruebas de ciberseguridad y certifica profesionales. Un pentest realizado por ellos evalúa un alcance y una fecha; no convierte al escriba en «certificado CREST». \href{https://www.crest-approved.org/about-us/what-we-do/}{CREST}, \href{https://www.crest-approved.org/skills-certifications-careers/crest-certifications/}{certificaciones profesionales}. \\
\end{longtable}

\textbf{Evidencia contractual, técnica y clínica complementaria:} declarar cumplimiento de RGPD, HIPAA o PIPEDA no demuestra una certificación externa; HHS no reconoce los sellos privados de «certificación HIPAA». DPA/\allowbreak{}BAA son contratos, PIA/\allowbreak{}EFVP son evaluaciones de privacidad; pentests, model cards y publicaciones son evidencias complementarias, con el alcance descrito en cada ficha. \href{https://www.hhs.gov/hipaa/for-professionals/faq/2003/are-we-required-to-certify-our-organizations-compliance-with-the-standards/index.html}{HHS}. Los estudios se separan y se señala la participación del proveedor. Una insignia «Verified» del portal no sustituye al certificado. El alcance de la revisión se limita a fuentes públicas, sin acceso a expedientes privados ni confirmación directa de los proveedores.

\endgroup

% !TEX root = ../main.tex
% GENERADO por scripts/sync_market_appendices.py; editar el Markdown de origen.
\clearpage
\section{Catálogo de proveedores y evidencia documental}
\label{app:proveedores}
\begingroup\footnotesize
El catálogo comprende 43 proveedores. Las fichas aplican criterios uniformes para distinguir declaraciones comerciales, documentación primaria examinada y limitaciones de la evidencia disponible.

\begin{longtable}{L{\dimexpr 0.08\linewidth-2\tabcolsep\relax}L{\dimexpr 0.92\linewidth-2\tabcolsep\relax}}
\toprule
\textbf{N.º} & \textbf{Proveedor /\allowbreak{} producto} \\
\midrule
\endfirsthead
\textbf{N.º} & \textbf{Proveedor /\allowbreak{} producto} \\
\midrule
\endhead
\midrule
\multicolumn{2}{r}{\footnotesize Continúa en la página siguiente}
\endfoot
\bottomrule
\endlastfoot
1 & Noa Notes: Doctoralia /\allowbreak{} Docplanner \\
2 & Invox Medical: Invox Genesis /\allowbreak{} Vócali \\
3 & Heidi: Scribe y Remote \\
4 & Solventum: Fluency Align \\
5 & Tandem Health \\
6 & Microsoft: Dragon Copilot \\
7 & Nabla \\
8 & Abridge \\
9 & Ambience Healthcare \\
10 & Suki \\
11 & DeepScribe \\
12 & Commure: Ambient AI \\
13 & Corti: plataforma/\allowbreak{}API y Assistant MD \\
14 & Oracle: Health Clinical AI Agent /\allowbreak{} Clinical Note \\
15 & Tortus \\
16 & Lyrebird \\
17 & Tali \\
18 & Scribeberry \\
19 & Quirónsalud: proyecto Scribe \\
20 & Freed \\
21 & 33n: CLEARnotes \\
22 & Accuro: Assistant \\
23 & Accurx: Scribe \\
24 & Anathem \\
25 & Anima: Annie /\allowbreak{} Ambient Scribe \\
26 & Aprobrium: Lexacom \\
27 & Beam Up: Beam Notes /\allowbreak{} Magic Notes \\
28 & Cerebric \\
29 & CLAI Ltd: ClinicLetter.ai /\allowbreak{} Mayden \\
30 & Dentistry Dashboard: AI Notes \\
31 & Dictate IT: AVT en Swift /\allowbreak{} DIT3 \\
32 & Docwise \\
33 & eConsult: Hi Scribe \\
34 & Epro: Bluewire Technologies \\
35 & G2 Speech: SpeechAmbient \\
36 & HealthOrbit AI \\
37 & Medhasyn Ltd: Clinical Companion \\
38 & Optum: EMIS Scribe \\
39 & Pungo: Joy /\allowbreak{} JoyNotes \\
40 & Scribetech: Augnito Omni \\
41 & T-Pro: Scribe /\allowbreak{} Copilot \\
42 & X-on Health: Surgery Intellect, powered by TORTUS \\
43 & AWS: HealthScribe \\
\end{longtable}
\subsection*{1. Noa Notes: Doctoralia /\allowbreak{} Docplanner}
\label{provider:1}

Documentación de referencia: \href{https://noa.ai/es-es/seguridad-y-privacidad}{Seguridad} y DPA. No se localizó un Trust Center específico de Notes con certificados descargables.

\begin{longtable}{L{\dimexpr 0.24\linewidth-2\tabcolsep\relax}L{\dimexpr 0.31\linewidth-2\tabcolsep\relax}L{\dimexpr 0.45\linewidth-2\tabcolsep\relax}}
\toprule
\textbf{Credencial} & \textbf{Emisor /\allowbreak{} documento} & \textbf{Alcance, vigencia y evidencia} \\
\midrule
\endfirsthead
\textbf{Credencial} & \textbf{Emisor /\allowbreak{} documento} & \textbf{Alcance, vigencia y evidencia} \\
\midrule
\endhead
\midrule
\multicolumn{3}{r}{\footnotesize Continúa en la página siguiente}
\endfoot
\bottomrule
\endlastfoot
ISO 27001 & Docplanner; \href{https://noa.ai/es-es/seguridad-y-privacidad}{Seguridad} & \textbf{Declarada} a nivel corporativo. Certificado, emisor, vigencia y alcance Notes pendientes de verificación. \\
\end{longtable}

\begin{itemize}
\item \textbf{Tratamiento de datos en España:} DPA público: AWS para alojamiento/\allowbreak{}transcripción y Azure para resúmenes, almacenamiento UE; también enumera servicios auxiliares en EE. UU./\allowbreak{}Australia. \textbf{Permite mejorar modelos con datos desidentificados si el paciente consiente}; esta condición excluye una prohibición absoluta de reutilización para entrenamiento. \href{https://notes.es.noa.ai/terms-and-conditions/noa-notes-es.pdf}{Contrato, §§3 y 6; anexo A y apéndice I}.
\item \textbf{Regulación sanitaria y calidad clínica:} el contrato declara que Notes \textbf{no es dispositivo médico} y exige revisión clínica. No se localizó evaluación del escriba revisada por pares; que Noa Evidence consulte literatura científica no valida Notes. \href{https://notes.es.noa.ai/terms-and-conditions/noa-notes-es.pdf}{Contrato}, \href{https://noa.ai/es-es/}{distinción de productos}.
\end{itemize}

\subsection*{2. Invox Medical: Invox Genesis /\allowbreak{} Vócali}
\label{provider:2}

Documentación de referencia: \href{https://www.invoxmedical.com/blog/invox-medical-refuerza-su-compromiso-con-la-seguridad-de-la-informacion-con-la-obtencion-de-la-certificacion-iso-27001-y-el-ens-nivel-alto}{Anuncio} y documentación del producto; no se localizó catálogo público de certificados completos.

\begin{longtable}{L{\dimexpr 0.24\linewidth-2\tabcolsep\relax}L{\dimexpr 0.31\linewidth-2\tabcolsep\relax}L{\dimexpr 0.45\linewidth-2\tabcolsep\relax}}
\toprule
\textbf{Credencial} & \textbf{Emisor /\allowbreak{} documento} & \textbf{Alcance, vigencia y evidencia} \\
\midrule
\endfirsthead
\textbf{Credencial} & \textbf{Emisor /\allowbreak{} documento} & \textbf{Alcance, vigencia y evidencia} \\
\midrule
\endhead
\midrule
\multicolumn{3}{r}{\footnotesize Continúa en la página siguiente}
\endfoot
\bottomrule
\endlastfoot
ISO 27001 & Emisor no identificado; \href{https://www.invoxmedical.com/blog/invox-medical-refuerza-su-compromiso-con-la-seguridad-de-la-informacion-con-la-obtencion-de-la-certificacion-iso-27001-y-el-ens-nivel-alto}{Anuncio} & \textbf{Declarada} por Invox/\allowbreak{}Vócali; certificado íntegro y vigencia no localizados. \\
ENS categoría ALTA & Emisor no identificado; \href{https://www.invoxmedical.com/blog/invox-medical-refuerza-su-compromiso-con-la-seguridad-de-la-informacion-con-la-obtencion-de-la-certificacion-iso-27001-y-el-ens-nivel-alto}{Anuncio} & \textbf{Declarada}; alcance Genesis pendiente. Legitec aparece como asesor, no certificador. \\
\end{longtable}

\begin{itemize}
\item \textbf{Tratamiento de datos:} anuncia migración de infraestructura a \textbf{AWS} y ausencia de almacenamiento de audios/\allowbreak{}información generada. País de procesamiento, DPA y lista completa de subencargados de Genesis pendientes de verificación. \href{https://www.invoxmedical.com/blog/migramos-a-aws-para-reforzar-la-seguridad-y-escalabilidad-de-nuestra-infraestructura-tecnologica}{Infraestructura}, \href{https://www.invoxmedical.com/}{producto}.
\item \textbf{Evidencia externa:} la \textbf{Comunidad de Madrid confirma Genesis en Psiquiatría del Hospital Fundación Alcorcón}, junio de 2025. Acredita adopción en un hospital público, sin demostrar despliegue general del SERMAS ni validación clínica revisada por pares. \href{https://www.comunidad.madrid/noticias/2025/06/10/hospital-publico-fundacion-alcorcon-incorpora-aplicacion-ia-optimizar-generacion-documentacion-clinica-partir-entrevistas}{Comunidad de Madrid}.
\end{itemize}

\subsection*{3. Heidi: Scribe y Remote}
\label{provider:3}

Documentación de referencia: \href{https://trust.heidihealth.com/}{Trust Center} y catálogo de recursos; combina PDFs públicos con documentos bajo solicitud.

\begin{longtable}{L{\dimexpr 0.24\linewidth-2\tabcolsep\relax}L{\dimexpr 0.31\linewidth-2\tabcolsep\relax}L{\dimexpr 0.45\linewidth-2\tabcolsep\relax}}
\toprule
\textbf{Credencial} & \textbf{Emisor /\allowbreak{} documento} & \textbf{Alcance, vigencia y evidencia} \\
\midrule
\endfirsthead
\textbf{Credencial} & \textbf{Emisor /\allowbreak{} documento} & \textbf{Alcance, vigencia y evidencia} \\
\midrule
\endhead
\midrule
\multicolumn{3}{r}{\footnotesize Continúa en la página siguiente}
\endfoot
\bottomrule
\endlastfoot
ISO 27001 & \textbf{A-LIGN}, ISMS-HE-042726; \href{https://trust.heidihealth.com/doc/trust?rid=6a4673e10d5928c7f9995cea\&r=qq3qqjsmu9pfv6v9q9xa7m\&view=true}{ISO} & \textbf{Certificado examinado}: ISMS de Scribe, Comms y Evidence; expira \textbf{10/\allowbreak{}05/\allowbreak{}2029}. \\
Cyber Essentials Plus & \textbf{AMTIVO}; \href{https://trust.heidihealth.com/doc/trust?rid=6a469d372506b7e736cfbd7d\&r=qq3qqjsmu9pfv6v9q9xa7m\&view=true}{CE+} & \textbf{Certificado examinado}: Heidi Health Ltd; recertificación \textbf{02/\allowbreak{}07/\allowbreak{}2027}. \\
TX-RAMP nivel 2 & \textbf{Texas DIR}, TX1569070; \href{https://trust.heidihealth.com/doc/trust?rid=66f7864e59f37adebcf7d679\&r=qq3qqjsmu9pfv6v9q9xa7m\&view=true}{TX-RAMP} & \textbf{Certificado examinado}; expira \textbf{26/\allowbreak{}09/\allowbreak{}2027}. \\
SOC 2 tipo II & \textbf{Prescient Assurance}; \href{https://trust.heidihealth.com/doc/trust?rid=6a1fad949f2b8d0d8c8937ca\&r=qq3qqjsmu9pfv6v9q9xa7m\&view=true}{carta SOC} & \textbf{Carta examinada}, sin salvedades; periodo \textbf{05/\allowbreak{}06/\allowbreak{}2024–04/\allowbreak{}06/\allowbreak{}2025}. Verificación del informe completo pendiente. \\
ISO 9001 /\allowbreak{} ISO 42001 & \href{https://www.heidihealth.com/it/safety}{Credenciales}; \href{https://trust.heidihealth.com/}{Trust Center} & \textbf{Declaradas}; emisor/\allowbreak{}alcance/\allowbreak{}vigencia pendientes de verificación. Certificado 42001 bajo solicitud. \\
C5 /\allowbreak{} UKCA & \href{https://trust.heidihealth.com/}{Trust Center} & Insignias publicadas; tipo C5, clase UKCA y alcance del software pendientes de verificación. \\
TGV: Québec & \href{https://www.quebec.ca/sante/professionnels/programme-transcription-intelligence-artificielle}{Québec} & \textbf{Registro oficial consultado}: TGV sí; EFVP de Santé Québec no. \\
\end{longtable}

\begin{itemize}
\item \textbf{Tratamiento de datos:} declara permanencia UE/\allowbreak{}EEE para la oferta europea, sin entrenamiento con consultas/\allowbreak{}notas; subencargados en Trust Centre. \textbf{Remote conserva audio cifrado localmente hasta sincronizar}, por lo que no procede generalizar «no almacena audio». \href{https://www.heidihealth.com/compliance/gdpr}{UE}, \href{https://support.heidihealth.com/en/articles/14648423-how-heidi-protects-your-data}{entrenamiento}, \href{https://trust.heidihealth.com/subprocessors}{subencargados}, \href{https://support.heidihealth.com/en/articles/14130260-getting-started-with-heidi-remote}{Remote}.
\item \textbf{Credenciales del dispositivo Remote:} manual declara FCC \textbf{2BULL-HV0}, ISED \textbf{35134-HV0}, conformidad radio UE/\allowbreak{}UK y ensayo antibacteriano SGS \textbf{GZF26-0000530-01}. Son credenciales del periférico; el fabricante indica que \textbf{no es dispositivo médico}. \href{https://www.heidihealth.com/en-ca/docs/remote-user-manual}{Manual}.
\item \textbf{Evidencia de evaluación:} ensayo aleatorizado con Heidi, \textbf{279 pacientes}, \emph{Journal of Arthroplasty} (2026): mejora eficiencia, pero registra \textbf{19 errores médicos frente a 9} con documentación convencional. Un único cirujano; no demuestra rendimiento general ni evalúa Remote. \href{https://doi.org/10.1016/j.arth.2026.03.069}{Artículo}.
\item \textbf{Evaluación en español:} preprint de agosto de 2026 con \textbf{385 consultas simuladas emparejadas en cinco países}, incluido un centro de Barcelona en español. Compara notas de Heidi y de clínicos con evaluación ciega; resultados agregados favorables. \textbf{Financiado íntegramente por Heidi, sin revisión por pares}; las 385 consultas no son todas españolas y no prueban resultados en pacientes reales. \href{https://www.medrxiv.org/content/10.64898/2026.08.18.26360701v1}{Estudio}.
\end{itemize}

\subsection*{4. Solventum: Fluency Align}
\label{provider:4}

Documentación de referencia: \href{https://assets.solventum.com/is/content/mmmspinco/fluency-align-key-benefits-and-security-data-sheet-his-cps-en-gbpdf}{Ficha UK} y \href{https://www.solventum.com/es-es/home/health-information-technology/solutions/fluency-align/}{FAQ España}; no se localizó expediente público de Align.

\begin{longtable}{L{\dimexpr 0.24\linewidth-2\tabcolsep\relax}L{\dimexpr 0.31\linewidth-2\tabcolsep\relax}L{\dimexpr 0.45\linewidth-2\tabcolsep\relax}}
\toprule
\textbf{Credencial} & \textbf{Emisor /\allowbreak{} documento} & \textbf{Alcance, vigencia y evidencia} \\
\midrule
\endfirsthead
\textbf{Credencial} & \textbf{Emisor /\allowbreak{} documento} & \textbf{Alcance, vigencia y evidencia} \\
\midrule
\endhead
\midrule
\multicolumn{3}{r}{\footnotesize Continúa en la página siguiente}
\endfoot
\bottomrule
\endlastfoot
ISO 27001 /\allowbreak{} Cyber Essentials & \href{https://assets.solventum.com/is/content/mmmspinco/fluency-align-key-benefits-and-security-data-sheet-his-cps-en-gbpdf}{Ficha UK} & \textbf{Declaradas}; certificados, emisor, vigencia y alcance Align con verificación documental pendiente. \\
ENS categoría ALTA & \href{https://www.solventum.com/es-es/home/health-information-technology/solutions/fluency-align/}{FAQ España} & Declara cumplimiento ENS ALTO; \textbf{certificado pendiente de verificación}. \\
ISO 42001 & \href{https://assets.solventum.com/is/content/mmmspinco/fluency-align-key-benefits-and-security-data-sheet-his-cps-en-gbpdf}{Ficha UK} & \textbf{Alineación}, no certificación demostrada. \\
\end{longtable}

\begin{itemize}
\item \textbf{Tratamiento de datos:} declara almacenamiento \textbf{y procesamiento en AWS Frankfurt, UE}; no cede datos para entrenar modelos externos. Esto no descarta entrenamiento interno. Retención, DPA y cadena completa de subencargados pendientes de verificar. \href{https://www.solventum.com/es-es/home/health-information-technology/solutions/fluency-align/}{FAQ España}.
\item \textbf{Regulación sanitaria y calidad clínica:} la ficha UK promete cumplimiento \textbf{DTAC/\allowbreak{}DCB0129 antes de firmar contrato}. No se localizó estudio revisado por pares específico de Align; el caso del Hospital del Sureste corresponde a \textbf{Fluency Direct}. \href{https://assets.solventum.com/is/content/mmmspinco/fluency-align-key-benefits-and-security-data-sheet-his-cps-en-gbpdf}{Ficha UK}, \href{https://www.solventum.com/es-es/home/health-information-technology/solutions/fluency-align/}{recursos y casos}.
\end{itemize}

\subsection*{5. Tandem Health}
\label{provider:5}

\textbf{Documentación externa examinada}, disponible en su \href{https://trust.tandemhealth.ai/}{Trust Center} y \href{https://trust.assuric.com/tandem/public/compliance}{catálogo Assuric}:

\begin{longtable}{L{\dimexpr 0.24\linewidth-2\tabcolsep\relax}L{\dimexpr 0.31\linewidth-2\tabcolsep\relax}L{\dimexpr 0.45\linewidth-2\tabcolsep\relax}}
\toprule
\textbf{Credencial} & \textbf{Emisor /\allowbreak{} documento} & \textbf{Alcance, vigencia y evidencia} \\
\midrule
\endfirsthead
\textbf{Credencial} & \textbf{Emisor /\allowbreak{} documento} & \textbf{Alcance, vigencia y evidencia} \\
\midrule
\endhead
\midrule
\multicolumn{3}{r}{\footnotesize Continúa en la página siguiente}
\endfoot
\bottomrule
\endlastfoot
\textbf{\href{https://file.assuric.com/orgs\%2Ftandem\%2FcontrolEvidence\%2FrTj5GDSXKa16FVqjjsV4\%2F25252991-\%20ENS\%20-\%20ALTA-TANDEM\%20HEALTH\%20AB\%20\%28SITE\%20VIRTUAL\%29-ESPANI\%CC\%80\_OL.pdf?alt=media\&token=77dff88d-f576-4c0a-a95e-77ed28223e1f}{ENS categoría ALTA}} & \textbf{Bureau Veritas}, ES154689-1 & Sistema de desarrollo/\allowbreak{}soporte del SaaS Tandem para atención clínica y administración; Azure Gävle. \textbf{10/\allowbreak{}09/\allowbreak{}2025–10/\allowbreak{}09/\allowbreak{}2027}. Copia pública: página 1/\allowbreak{}2. \\
\textbf{\href{https://file.assuric.com/orgs\%2Ftandem\%2FcontrolEvidence\%2F4YOEj8c6wuyZ7rpTKD4u\%2F05\%20ISMS\%20Certificate\%20v2.0\%20-\%20Tandem\%20Health\%20AB.pdf?alt=media\&token=8e3ca92c-f0db-475a-bc3b-06044a35716f}{ISO 27001:2022}} & \textbf{Insight Assurance}, IA-2024-12-17-02 & Gestión de seguridad de Tandem Health AB; versión 05/\allowbreak{}02/\allowbreak{}2026, expira \textbf{17/\allowbreak{}12/\allowbreak{}2027}. \\
\textbf{\href{https://file.assuric.com/orgs\%2Ftandem\%2FcontrolEvidence\%2Fs4KoRqo4iYEUdz0Jmu7p\%2FSCAR-9.1.1.0\%20-\%20Tandem\%20-\%20ISO\%2013485\%20Certificate.pdf?alt=media\&token=b87e7394-d280-44fd-a73f-cb75d27b1126}{ISO 13485:2016}} & \textbf{Scarlet}, SCAR-9.1.1.0 & Diseño/\allowbreak{}fabricación de software sanitario de documentación; \textbf{15/\allowbreak{}10/\allowbreak{}2025–14/\allowbreak{}10/\allowbreak{}2028}. \\
\textbf{\href{https://file.assuric.com/orgs\%2Ftandem\%2FcontrolEvidence\%2FpoCpC7GQYwIJRb65B40G\%2FSCAR-9.9.1.0.pdf?alt=media\&token=3196aaed-1602-4fc2-b196-2e6939bb6cc4}{ISO 42001:2023}} & \textbf{Scarlet}, SCAR-9.9.1.0 & Software sanitario con IA y usos internos de IA; expira \textbf{06/\allowbreak{}05/\allowbreak{}2029}. \\
\textbf{\href{https://file.assuric.com/orgs\%2Ftandem\%2FcontrolEvidence\%2FqkNFoPUk2Ck4cMsMGwgr\%2F0245002\_ENG\_14001\_Tandem\%20Health\%20AB.pdf?alt=media\&token=9c97057e-cff8-404a-86bc-43c491c3dd3c}{ISO 14001:2015}} & \textbf{Intertek}, 0245002 & Gestión ambiental del desarrollo/\allowbreak{}operación de software sanitario; expira \textbf{06/\allowbreak{}05/\allowbreak{}2029}. \\
\textbf{\href{https://file.assuric.com/orgs\%2Ftandem\%2FcontrolEvidence\%2FjGRsY5eN0u0cu7ynzA9o\%2FNL\%202864.1.2\%20-\%20Tandem\%20Health\%20AB\%20-\%20NEN\%207510.pdf?alt=media\&token=51a06ba9-5b80-4eb5-96ef-17818c7d3ee2}{NEN 7510-1:2024}} & \textbf{Brand Compliance}, NL 2864.1.2 & Seguridad de la plataforma de documentación clínica; alojamiento externalizado. Expira \textbf{15/\allowbreak{}04/\allowbreak{}2029}. \\
\textbf{\href{https://trust.tandemhealth.ai/doc/trust?rid=6a79840c39de23960acdd51f\&r=e0qkrepqo674el4jxq28px\&view=true}{C5 tipo II: dictamen}} & \textbf{HKKG}, 16/\allowbreak{}06/\allowbreak{}2026 & SaaS Tandem; periodo \textbf{17/\allowbreak{}02–18/\allowbreak{}05/\allowbreak{}2026}. \textbf{Con salvedades}: verificaciones de personal, formación, manejo de metadatos y estrategias de salida de proveedores. \\
\textbf{\href{https://file.assuric.com/orgs\%2Ftandem\%2FcontrolEvidence\%2F7f7koYGmQJa8TWFMTNqm\%2FSCAR-9.2.1.5\%20-\%20Tandem\%20-\%20Quality\%20EU\%20MDR\%20Certificate.pdf?alt=media\&token=aa8686b7-6a03-4f04-9d4a-f40e217d6bd1}{CE MDR clase IIa}} & \textbf{Scarlet NB B.V., 3022}, SCAR-9.2.1.5 & Incluye explícitamente \textbf{AI Medical Scribe}, Coder y Guide; revisión 18/\allowbreak{}06/\allowbreak{}2026, expira \textbf{15/\allowbreak{}02/\allowbreak{}2031}. Requiere revisión profesional de las salidas. \\
\end{longtable}

\begin{itemize}
\item \textbf{Otras credenciales:} \textbf{Cyber Essentials Plus} declarado, con enlace a registro; \textbf{Ethik-IA} anunciado, sin certificado localizado. \textbf{Astori/\allowbreak{}LVV confirma Scribe v1.0.0, clase A1}, certificado de seguridad \textbf{FI251028-214}, 28/\allowbreak{}10/\allowbreak{}2025–28/\allowbreak{}10/\allowbreak{}2028. \textbf{DSPT 2025–26: “Standards exceeded”}; DTAC/\allowbreak{}DCB0129 con documentación bajo solicitud. \href{https://trust.tandemhealth.ai/}{Trust Center}, \href{https://astori.lvv.fi/}{Astori: buscar Tandem}, \href{https://www.dsptoolkit.nhs.uk/OrganisationSearch/TAHEAB001}{DSPT oficial}.
\item \textbf{Documentación complementaria examinada:} \textbf{libre venta}, Agencia Sueca del Medicamento, 10/\allowbreak{}08/\allowbreak{}2026, incluye Scribe/\allowbreak{}Coder/\allowbreak{}Guide; \textbf{WEEE 2026}, adhesión a El-Kretsen para residuos electrónicos; micrófono \textbf{BM10}, declaración propia RoHS/\allowbreak{}EMC. La representación suiza es un mandato, no certificación del producto; también enlaza inscripción estonia. \href{https://trust.tandemhealth.ai/doc/trust?rid=6a86ba9438622132aac7533c\&r=e0qkrepqo674el4jxq28px\&view=true}{Libre venta}, \href{https://trust.tandemhealth.ai/doc/trust?rid=6a78f060820290d89042a93e\&r=e0qkrepqo674el4jxq28px\&view=true}{WEEE}, \href{https://trust.tandemhealth.ai/doc/trust?rid=6a78f0c64c2eeb0a116d6fb1\&r=e0qkrepqo674el4jxq28px\&view=true}{micrófono}, \href{https://www.chrep.me/certificates-ch-rep/ch-rep-eu6romde4}{mandato suizo}, \href{https://trust.tandemhealth.ai/resources}{catálogo}.
\item \textbf{Tratamiento de datos:} declara procesamiento/\allowbreak{}almacenamiento UE, audio eliminado tras transcribir y sin entrenamiento con datos de pacientes. Lista actual: \textbf{Azure, ElevenLabs, Snowflake, Intercom} y entidades del grupo, con ubicaciones UE declaradas. \href{https://tandemhealth.ai/product/data-security}{Seguridad}, \href{https://trust.tandemhealth.ai/subprocessors}{subencargados}.
\item \textbf{Evidencia de evaluación:} \emph{JMIR Medical Informatics} (2026), \textbf{236.153 notas y encuesta a 177 usuarios en Suecia}, observacional, con autores de Tandem: tiempo autodeclarado por nota \ensuremath{-}29\%; no es ensayo de seguridad clínica ni validación en español. El documento de orientación para la contratación es una guía comercial de contratación. \href{https://tandemhealth.ai/resources/research/evaluation-of-an-ai-medical-scribe-after-236153-notes-generated-across-care-levels-in-a-european-health-system}{Estudio y resumen}, \href{https://doi.org/10.2196/90052}{DOI}, \href{https://mcusercontent.com/92915049aef2f51e4bcdd8434/files/a252c926-b88a-0af8-4071-a19fed8b9bf5/Procurement\_of\_Ambient\_Voice\_Technology\_EN.pdf}{whitepaper}.
\end{itemize}

\subsection*{6. Microsoft: Dragon Copilot}
\label{provider:6}

Documentación de referencia: \href{https://learn.microsoft.com/en-us/industry/healthcare/dragon-copilot/about/faqs}{FAQ} y catálogo de cumplimiento de Microsoft; informes completos con verificación documental pendiente.

\begin{longtable}{L{\dimexpr 0.24\linewidth-2\tabcolsep\relax}L{\dimexpr 0.31\linewidth-2\tabcolsep\relax}L{\dimexpr 0.45\linewidth-2\tabcolsep\relax}}
\toprule
\textbf{Credencial} & \textbf{Emisor /\allowbreak{} documento} & \textbf{Alcance, vigencia y evidencia} \\
\midrule
\endfirsthead
\textbf{Credencial} & \textbf{Emisor /\allowbreak{} documento} & \textbf{Alcance, vigencia y evidencia} \\
\midrule
\endhead
\midrule
\multicolumn{3}{r}{\footnotesize Continúa en la página siguiente}
\endfoot
\bottomrule
\endlastfoot
ISO 42001 & \href{https://learn.microsoft.com/en-us/compliance/regulatory/offering-iso-42001}{Alcance 42001} & Microsoft incluye Dragon Copilot en el alcance publicado; verificación del certificado íntegro pendiente. \\
SOC 1 II /\allowbreak{} SOC 2 II /\allowbreak{} C5 II & \href{https://learn.microsoft.com/en-us/industry/healthcare/dragon-copilot/about/faqs}{FAQ} & \textbf{Declarados} para Dragon Copilot; auditor, periodos e informes pendientes de verificación. \\
NEN 7510 /\allowbreak{} 7512 /\allowbreak{} 7513 & \href{https://learn.microsoft.com/en-us/industry/healthcare/dragon-copilot/about/faqs}{FAQ} & Cumplimiento declarado; certificados individuales con verificación documental pendiente. \\
ISO 27001 /\allowbreak{} ISO 27701 & \href{https://learn.microsoft.com/en-us/industry/healthcare/dragon-copilot/about/faqs}{FAQ} & \textbf{Pendientes} en la FAQ de 27/\allowbreak{}07/\allowbreak{}2026; no se cuentan como obtenidas. \\
UKCA clase I & \href{https://learn.microsoft.com/en-us/industry/healthcare/dragon-copilot/about/eifu}{Instrucciones} & \textbf{Autodeclaración} del fabricante; no evaluación independiente de clase IIa. \\
\end{longtable}

\begin{itemize}
\item \textbf{Tratamiento de datos:} EU Data Boundary, con excepciones contractuales; documentación europea identifica Francia/\allowbreak{}Alemania, sin concretar España. Retención hasta \textbf{90 días}; \textbf{reutiliza datos anonimizados para mejorar modelos}, con posible intervención humana. Publica enlaces a DPA/\allowbreak{}subencargados. \href{https://learn.microsoft.com/en-us/industry/healthcare/dragon-copilot/whitepapers/privacy}{Privacidad}, \href{https://www.microsoft.com/licensing/terms/en-US/product/PrivacyandSecurityTerms/all}{condiciones}.
\item \textbf{Regulación sanitaria y evaluación clínica:} UKCA clase I \textbf{autocertificada}. El ensayo \emph{NEJM AI} con 238 médicos evaluó \textbf{DAX Copilot}, su predecesor: relevante, pero no valida toda la versión actual; DAX no redujo significativamente el tiempo de documentación (\emph{time-in-note}) frente al control. \href{https://learn.microsoft.com/en-us/industry/healthcare/dragon-copilot/about/eifu}{Instrucciones}, \href{https://doi.org/10.1056/AIoa2501000}{ensayo}.
\end{itemize}

\subsection*{7. Nabla}
\label{provider:7}

Documentación de referencia: \href{https://trust.nabla.com/?format=html}{Trust Center}; certificados, SOC y pentest requieren acceso.

\begin{longtable}{L{\dimexpr 0.24\linewidth-2\tabcolsep\relax}L{\dimexpr 0.31\linewidth-2\tabcolsep\relax}L{\dimexpr 0.45\linewidth-2\tabcolsep\relax}}
\toprule
\textbf{Credencial} & \textbf{Emisor /\allowbreak{} documento} & \textbf{Alcance, vigencia y evidencia} \\
\midrule
\endfirsthead
\textbf{Credencial} & \textbf{Emisor /\allowbreak{} documento} & \textbf{Alcance, vigencia y evidencia} \\
\midrule
\endhead
\midrule
\multicolumn{3}{r}{\footnotesize Continúa en la página siguiente}
\endfoot
\bottomrule
\endlastfoot
ISO 27001 & \href{https://trust.nabla.com/?format=html}{Trust Center} & \textbf{Declarada}; documento de septiembre de 2025 bajo solicitud. \\
SOC 2 tipo II & \href{https://trust.nabla.com/?format=html}{Trust Center} & \textbf{Declarado}; periodo noviembre 2024–octubre 2025. Informe/\allowbreak{}auditor con verificación documental pendiente. \\
TX-RAMP nivel 2 & \href{https://trust.nabla.com/?format=html}{Trust Center} & \textbf{Declarado}; documento enero de 2026 bajo solicitud. \\
HDS & \href{https://www.nabla.com/compliance-security}{oferta europea} & Atribuido a \textbf{infraestructura}; certificado propio de Nabla no demostrado. \\
\end{longtable}

\begin{itemize}
\item \textbf{Tratamiento de datos:} GCP Bélgica para clientes fuera de EE. UU.; oferta europea declara tratamiento UE. Audio no conservado por defecto; notas \textbf{14 días}, configurable. DPA identifica Google/\allowbreak{}Microsoft Ireland y permite reutilización anonimizada con autorización expresa; la matriz regional de subencargados requiere confirmación contractual. \href{https://trust.nabla.com/?format=html}{Trust Center}, \href{https://www.nabla.com/fr/security-documents/data-protection-agreement}{DPA}.
\item \textbf{Evidencia de evaluación:} ensayo aleatorizado \emph{NEJM AI} (2025), \textbf{238 médicos/\allowbreak{}14 especialidades}, Nabla/\allowbreak{}DAX/\allowbreak{}control: reducción modesta del tiempo de documentación con Nabla; un centro y consultas en inglés. Model card y adhesión CHAI no se cuentan como certificación clínica. \href{https://doi.org/10.1056/AIoa2501000}{Ensayo}, \href{https://trust.nabla.com/?format=html}{documentos de gobernanza}.
\end{itemize}

\subsection*{8. Abridge}
\label{provider:8}

Documentación de referencia: \href{https://trust.abridge.com/}{Trust Center}; SOC, pentest y documentación de IA requieren acceso.

\begin{longtable}{L{\dimexpr 0.24\linewidth-2\tabcolsep\relax}L{\dimexpr 0.31\linewidth-2\tabcolsep\relax}L{\dimexpr 0.45\linewidth-2\tabcolsep\relax}}
\toprule
\textbf{Credencial} & \textbf{Emisor /\allowbreak{} documento} & \textbf{Alcance, vigencia y evidencia} \\
\midrule
\endfirsthead
\textbf{Credencial} & \textbf{Emisor /\allowbreak{} documento} & \textbf{Alcance, vigencia y evidencia} \\
\midrule
\endhead
\midrule
\multicolumn{3}{r}{\footnotesize Continúa en la página siguiente}
\endfoot
\bottomrule
\endlastfoot
SOC 2 tipos I y II & \href{https://trust.abridge.com/}{Trust Center} & \textbf{Declarados} para sus productos, seguridad/\allowbreak{}confidencialidad; informes bajo solicitud. \\
TX-RAMP & \href{https://trust.abridge.com/}{Trust Center} & Mención en portal; \textbf{nivel, vigencia y documento pendientes de verificación}. \\
\end{longtable}

\begin{itemize}
\item \textbf{Tratamiento de datos:} soporte declara almacenamiento/\allowbreak{}procesamiento en \textbf{EE. UU.} La política general excluye los datos clínicos de clientes, regidos por contrato/\allowbreak{}BAA: retención, entrenamiento y subencargados del escriba necesitan esos documentos. \href{https://support.abridge.com/hc/en-us/articles/30235294201619-Data-Security-and-HIPAA-Compliance}{Soporte}, \href{https://www.abridge.com/privacy}{privacidad}.
\item \textbf{Evidencia de evaluación:} \emph{JAMA Network Open} (2025), antes/\allowbreak{}después, \textbf{263 profesionales de seis sistemas}, mejora autodeclarada de carga/\allowbreak{}burnout; sin aleatorización y con coautor del proveedor. También hay trabajo específico de aseguramiento de calidad en \emph{NEJM AI}. \href{https://jamanetwork.com/journals/jamanetworkopen/fullarticle/2839542}{JAMA}, \href{https://doi.org/10.1056/AIcs2400977}{NEJM AI}, \href{https://www.abridge.com/ai/publications}{catálogo}.
\end{itemize}

\subsection*{9. Ambience Healthcare}
\label{provider:9}

Documentación de referencia: \href{https://trust.ambiencehealthcare.com/}{Trust Center y recursos}; documentos bajo solicitud.

\begin{longtable}{L{\dimexpr 0.24\linewidth-2\tabcolsep\relax}L{\dimexpr 0.31\linewidth-2\tabcolsep\relax}L{\dimexpr 0.45\linewidth-2\tabcolsep\relax}}
\toprule
\textbf{Credencial} & \textbf{Emisor /\allowbreak{} documento} & \textbf{Alcance, vigencia y evidencia} \\
\midrule
\endfirsthead
\textbf{Credencial} & \textbf{Emisor /\allowbreak{} documento} & \textbf{Alcance, vigencia y evidencia} \\
\midrule
\endhead
\midrule
\multicolumn{3}{r}{\footnotesize Continúa en la página siguiente}
\endfoot
\bottomrule
\endlastfoot
SOC 2 tipos I y II & \href{https://trust.ambiencehealthcare.com/}{Trust Center} & \textbf{Declarados}; catálogo lista SOC 2 II 2026, pentest 2026 y subencargados, todos bajo solicitud. \\
Evaluación de codificación AAPC & \href{https://www.ambiencehealthcare.com/}{AAPC según proveedor} & Anuncia \textbf{95\%}; informe externo completo y metodología no localizados. \\
ISO 27001 /\allowbreak{} Cyber Essentials del revendedor & \href{https://www.applytosupply.digitalmarketplace.service.gov.uk/g-cloud/services/285898651648070}{ficha Hitachi} & Titular \textbf{Hitachi}; no se atribuyen automáticamente al escriba Ambience. \\
\end{longtable}

\begin{itemize}
\item \textbf{Tratamiento de datos:} esa oferta G-Cloud declara almacenamiento/\allowbreak{}procesamiento \textbf{UK}; no demuestra residencia UE. La política remite al contrato/\allowbreak{}BAA para datos clínicos y permite compartir información desidentificada según contrato; retención y lista completa de subencargados pendientes de verificación. \href{https://www.applytosupply.digitalmarketplace.service.gov.uk/g-cloud/services/285898651648070}{Ficha}, \href{https://www.ambiencehealthcare.com/privacy}{privacidad}.
\item \textbf{Evidencia de evaluación:} \emph{npj Health Systems} (2026), implantación en Cleveland Clinic a \textbf{más de 4.000 profesionales} en cuatro meses, con coautores del proveedor. Es evidencia de implementación, no ensayo de exactitud clínica. \href{https://www.nature.com/articles/s44401-026-00144-6}{Artículo}.
\end{itemize}

\subsection*{10. Suki}
\label{provider:10}

Documentación de referencia: \href{https://trust.suki.ai/}{Trust Center} y subencargados; informes bajo solicitud.

\begin{longtable}{L{\dimexpr 0.24\linewidth-2\tabcolsep\relax}L{\dimexpr 0.31\linewidth-2\tabcolsep\relax}L{\dimexpr 0.45\linewidth-2\tabcolsep\relax}}
\toprule
\textbf{Credencial} & \textbf{Emisor /\allowbreak{} documento} & \textbf{Alcance, vigencia y evidencia} \\
\midrule
\endfirsthead
\textbf{Credencial} & \textbf{Emisor /\allowbreak{} documento} & \textbf{Alcance, vigencia y evidencia} \\
\midrule
\endhead
\midrule
\multicolumn{3}{r}{\footnotesize Continúa en la página siguiente}
\endfoot
\bottomrule
\endlastfoot
SOC 2 tipo II & \href{https://trust.suki.ai/}{Trust Center} & \textbf{Declarado}; informe \textbf{26/\allowbreak{}05/\allowbreak{}2025–25/\allowbreak{}05/\allowbreak{}2026} bajo solicitud; auditor pendiente de verificación. \\
TX-RAMP & \href{https://trust.suki.ai/}{Trust Center} & \textbf{Declarado}; nivel y vigencia pendientes de verificación. No se encontró respaldo primario para ISO/\allowbreak{}HITRUST. \\
\end{longtable}

\begin{itemize}
\item \textbf{Tratamiento de datos:} Compose declara entrenamiento con \textbf{datos anonimizados} y grabaciones en nube cifrada; BAA público y subencargados publicados en el Trust Center. Residencia UE y plazo del audio pendientes de verificación. Las condiciones deben verificarse para la solución contratada. \href{https://www.suki.ai/suki-compose/}{Compose}, \href{https://app.suki.ai/business-associate-agreement}{BAA}, \href{https://trust.suki.ai/subprocessors}{subencargados}.
\item \textbf{Evidencia de evaluación:} \emph{Frontiers in Artificial Intelligence} (2025), comparación ciega de notas de \textbf{97 consultas/\allowbreak{}cinco especialidades}: alucinaciones en 31\% de notas Suki frente a 20\% de referencia médica. Datos de 2024, en inglés, autores vinculados comercialmente a Suki; no es tasa universal del producto actual. \href{https://www.frontiersin.org/journals/artificial-intelligence/articles/10.3389/frai.2025.1691499/full}{Artículo}.
\end{itemize}

\subsection*{11. DeepScribe}
\label{provider:11}

Documentación de referencia: \href{https://trust.deepscribe.ai/}{Trust Center y actualizaciones}; documentos bajo solicitud.

\begin{longtable}{L{\dimexpr 0.24\linewidth-2\tabcolsep\relax}L{\dimexpr 0.31\linewidth-2\tabcolsep\relax}L{\dimexpr 0.45\linewidth-2\tabcolsep\relax}}
\toprule
\textbf{Credencial} & \textbf{Emisor /\allowbreak{} documento} & \textbf{Alcance, vigencia y evidencia} \\
\midrule
\endfirsthead
\textbf{Credencial} & \textbf{Emisor /\allowbreak{} documento} & \textbf{Alcance, vigencia y evidencia} \\
\midrule
\endhead
\midrule
\multicolumn{3}{r}{\footnotesize Continúa en la página siguiente}
\endfoot
\bottomrule
\endlastfoot
SOC 2 tipo II & \href{https://trust.deepscribe.ai/}{Trust Center y actualizaciones} & \textbf{Declarado}; informe \textbf{01/\allowbreak{}01–31/\allowbreak{}12/\allowbreak{}2024}, carta de continuidad 03/\allowbreak{}10/\allowbreak{}2025, bajo solicitud. \\
Auditoría HIPAA & \href{https://trust.deepscribe.ai/}{Trust Center y actualizaciones} & Documento bajo solicitud; \textbf{no certificación HIPAA oficial}. Vanta aloja el portal, no es el auditor identificado. \\
\end{longtable}

\begin{itemize}
\item \textbf{Tratamiento de datos:} identifica \textbf{AWS SageMaker/\allowbreak{}Bedrock}; publica cifrado, desidentificación y controles de acceso. Retención según necesidad/\allowbreak{}contrato/\allowbreak{}BAA; residencia UE, cadena completa de subencargados y prohibición contractual de entrenamiento pendientes de verificación. \href{https://www.deepscribe.ai/resources/deepscribe-launches-new-and-innovative-trust-and-safety-suite-for-its-ambient-ai-technology}{Arquitectura y controles}, \href{https://www.deepscribe.ai/privacy}{privacidad}.
\item \textbf{Evidencia de evaluación:} \textbf{DeepScore}, arXiv (2024), describe evaluación de errores/\allowbreak{}entidades/\allowbreak{}aceptación; \textbf{preprint del proveedor}, sin publicación revisada por pares identificada en las fuentes consultadas. Su \emph{Trust and Safety Suite} es funcionalidad propia, no acreditación externa; no se localizó validación oncológica independiente específica. \href{https://arxiv.org/abs/2409.16307}{DeepScore}, \href{https://www.deepscribe.ai/resources/deepscribe-launches-new-and-innovative-trust-and-safety-suite-for-its-ambient-ai-technology}{suite}.
\end{itemize}

\subsection*{12. Commure: Ambient AI}
\label{provider:12}

Documentación de referencia: \href{https://www.commure.com/trust-center}{Trust Center}; no ofrece un informe que resuelva la ambigüedad de alcance SOC.

\begin{longtable}{L{\dimexpr 0.24\linewidth-2\tabcolsep\relax}L{\dimexpr 0.31\linewidth-2\tabcolsep\relax}L{\dimexpr 0.45\linewidth-2\tabcolsep\relax}}
\toprule
\textbf{Credencial} & \textbf{Emisor /\allowbreak{} documento} & \textbf{Alcance, vigencia y evidencia} \\
\midrule
\endfirsthead
\textbf{Credencial} & \textbf{Emisor /\allowbreak{} documento} & \textbf{Alcance, vigencia y evidencia} \\
\midrule
\endhead
\midrule
\multicolumn{3}{r}{\footnotesize Continúa en la página siguiente}
\endfoot
\bottomrule
\endlastfoot
SOC 2 & \href{https://discover.commure.com/hubfs/Commure\%20Ambient\%20AI\%202025\%20First\%20Look\%20-\%20Final-1.pdf}{KLAS, p. 4}; \href{https://www.commure.com/trust-center}{Trust Center} & Declaración en KLAS; portal ambiguo y referencia tipo II de \textbf{centros de datos}. Informe del Ambient pendiente de verificación. \\
HITRUST & \href{https://discover.commure.com/hubfs/Commure\%20Ambient\%20AI\%202025\%20First\%20Look\%20-\%20Final-1.pdf}{KLAS, p. 4} & Ficha 2025: \textbf{no}. Certificación posterior del producto no localizada. \\
CE MDR clase I /\allowbreak{} registro MHRA & \href{https://www.commure.com/press-releases/commure-ambient-ai-achieves-class-i-certification-and-mhra-registration}{Anuncio} & Anuncio 03/\allowbreak{}09/\allowbreak{}2026: \textbf{autodeclarados}; no certificado de organismo notificado. \\
ISO 14971 /\allowbreak{} IEC 62366-1 & \href{https://www.commure.com/press-releases/commure-ambient-ai-achieves-class-i-certification-and-mhra-registration}{Anuncio} & Normas aplicadas; certificados externos con verificación documental pendiente. \\
\end{longtable}

\begin{itemize}
\item \textbf{Tratamiento de datos y evaluación:} anuncia residencia \textbf{UK}, mientras el Trust Center describe infraestructura estadounidense; retención contractual y subencargados bajo solicitud. La modalidad regional requiere confirmación contractual. KLAS recoge experiencia de siete organizaciones, \textbf{sin revisión por pares ni validación clínica}; no se localizó una publicación científica del producto Ambient actual. \href{https://www.commure.com/press-releases/commure-ambient-ai-achieves-class-i-certification-and-mhra-registration}{Anuncio UK}, \href{https://www.commure.com/trust-center}{datos}, \href{https://discover.commure.com/hubfs/Commure\%20Ambient\%20AI\%202025\%20First\%20Look\%20-\%20Final-1.pdf}{KLAS}.
\end{itemize}

\subsection*{13. Corti: plataforma/\allowbreak{}API y Assistant MD}
\label{provider:13}

Documentación de referencia: \href{https://corti.ai/safety}{Seguridad} y base de conocimiento; \href{https://app.drata.com/trust/9cc6d701-0c38-11ee-865f-029d78a187d9}{portal Drata} con acceso técnico al portal, pero sin contenido susceptible de examen documental.

\begin{longtable}{L{\dimexpr 0.24\linewidth-2\tabcolsep\relax}L{\dimexpr 0.31\linewidth-2\tabcolsep\relax}L{\dimexpr 0.45\linewidth-2\tabcolsep\relax}}
\toprule
\textbf{Credencial} & \textbf{Emisor /\allowbreak{} documento} & \textbf{Alcance, vigencia y evidencia} \\
\midrule
\endfirsthead
\textbf{Credencial} & \textbf{Emisor /\allowbreak{} documento} & \textbf{Alcance, vigencia y evidencia} \\
\midrule
\endhead
\midrule
\multicolumn{3}{r}{\footnotesize Continúa en la página siguiente}
\endfoot
\bottomrule
\endlastfoot
ISO 27001 /\allowbreak{} 27017 /\allowbreak{} 27018 & \textbf{A-LIGN}, según \href{https://corti.ai/stories/corti-completes-iso-audits-with-zero-nonconformities}{ISO} & Auditoría anunciada \textbf{16/\allowbreak{}06/\allowbreak{}2026}; certificados y vigencias con verificación documental pendiente. \\
ISO 42001 /\allowbreak{} Cyber Essentials Plus & \href{https://corti.ai/stories/corti-achieves-iso-iec-42001-certification}{42001/\allowbreak{}CE+} & \textbf{Declaradas}; documentación íntegra pendiente de verificación. \\
SOC 2 II /\allowbreak{} C5 II /\allowbreak{} ISAE 3000 II & \href{https://corti.ai/stories/how-corti-supports-medical-device-manufacturers-regulatory-ready-ai}{auditorías} & \textbf{Declaradas}; última, privacidad. Auditor, periodos y dictámenes con verificación documental pendiente. \\
Clase I UK/\allowbreak{}UE: Assistant MD & \href{https://corti.ai/stories/corti-achieves-eu-medical-device-registration-strengthening-the-foundation-for-regulated-healthcare-ai}{Alcance regulatorio} & Registro declarado para \textbf{Assistant MD}; no se traslada a la API ni a productos desarrollados sobre la plataforma. \\
ISO 13485 /\allowbreak{} ISO 14971 & \href{https://corti.ai/stories/corti-achieves-eu-medical-device-registration-strengthening-the-foundation-for-regulated-healthcare-ai}{Alcance regulatorio} & Conformidad declarada; \textbf{certificación no demostrada}. \\
\end{longtable}

\begin{itemize}
\item \textbf{Tratamiento de datos:} ofrece región UE, DPA público y aviso de nuevos subencargados; contrato prohíbe entrenar modelos para otros clientes salvo \textbf{opt-in expreso}. No significa ausencia de todo ajuste específico del cliente. \href{https://corti.ai/safety}{Seguridad}, \href{https://corti.ai/legal/master-terms-of-service}{contrato}, \href{https://corti.ai/legal/data-processing-agreement}{DPA}.
\item \textbf{Evidencia de evaluación:} \textbf{FactsR} evalúa documentación, pero es preprint del proveedor; el trabajo de ictus no valida el escriba. \href{https://arxiv.org/abs/2505.10360}{FactsR}, \href{https://corti.ai/research}{investigación}.
\end{itemize}

\subsection*{14. Oracle: Health Clinical AI Agent /\allowbreak{} Clinical Note}
\label{provider:14}

Documentación de referencia: \href{https://www.oracle.com/uk/corporate/acquisitions/cerner/security/}{Certificaciones y auditorías OHAI}; no se localizó certificado específico de Clinical Note.

\begin{longtable}{L{\dimexpr 0.24\linewidth-2\tabcolsep\relax}L{\dimexpr 0.31\linewidth-2\tabcolsep\relax}L{\dimexpr 0.45\linewidth-2\tabcolsep\relax}}
\toprule
\textbf{Credencial} & \textbf{Emisor /\allowbreak{} documento} & \textbf{Alcance, vigencia y evidencia} \\
\midrule
\endfirsthead
\textbf{Credencial} & \textbf{Emisor /\allowbreak{} documento} & \textbf{Alcance, vigencia y evidencia} \\
\midrule
\endhead
\midrule
\multicolumn{3}{r}{\footnotesize Continúa en la página siguiente}
\endfoot
\bottomrule
\endlastfoot
ISO 27001 & \href{https://www.oracle.com/uk/corporate/acquisitions/cerner/security/}{Certificaciones y auditorías OHAI} & Programa corporativo OHAI; \textbf{alcance Clinical Note pendiente de verificación}. \\
SOC 1 /\allowbreak{} SOC 2 tipo II & \href{https://www.oracle.com/uk/corporate/acquisitions/cerner/security/}{Certificaciones y auditorías OHAI} & Determinados entornos de alojamiento; no atribuidos automáticamente al escriba. \\
ACR/\allowbreak{}VPAT & \href{https://docs.oracle.com/en/corporate/accessibility/templates/t2-16521.html}{ACR} & Declaración de \textbf{accesibilidad} de Oracle; no certificación clínica. \\
\end{longtable}

\begin{itemize}
\item \textbf{Tratamiento de datos:} OHAI describe alojamiento en el país del cliente salvo acuerdo distinto, con acceso operativo internacional posible. No acredita procesamiento íntegro UE de Clinical Note; región, retención, entrenamiento y subencargados específicos pendientes. La consola contempla archivar/\allowbreak{}descargar grabaciones. \href{https://www.oracle.com/uk/corporate/acquisitions/cerner/security/}{OHAI}, \href{https://docs.oracle.com/en/corporate/accessibility/templates/t2-16689.html}{consola}.
\item \textbf{Evidencia de calidad clínica:} los resultados localizados son pilotos/\allowbreak{}casos NHS; no se localizó evaluación del escriba revisada por pares ni documentación primaria de certificación sanitaria accesible. \href{https://www.oracle.com/uk/news/announcement/oracle-health-clinicalai-agent-helps-uk-doctors-spend-more-time-on-patient-care-2026-02-12/}{Pilotos}.
\end{itemize}

\subsection*{15. Tortus}
\label{provider:15}

Documentación de referencia: \href{https://trust.tortus.ai/}{Trust Center} y recursos públicos; otros documentos requieren acceso.

\begin{longtable}{L{\dimexpr 0.24\linewidth-2\tabcolsep\relax}L{\dimexpr 0.31\linewidth-2\tabcolsep\relax}L{\dimexpr 0.45\linewidth-2\tabcolsep\relax}}
\toprule
\textbf{Credencial} & \textbf{Emisor /\allowbreak{} documento} & \textbf{Alcance, vigencia y evidencia} \\
\midrule
\endfirsthead
\textbf{Credencial} & \textbf{Emisor /\allowbreak{} documento} & \textbf{Alcance, vigencia y evidencia} \\
\midrule
\endhead
\midrule
\multicolumn{3}{r}{\footnotesize Continúa en la página siguiente}
\endfoot
\bottomrule
\endlastfoot
ISO 27001 & \textbf{LRQA}, 10712929; \href{https://trust.tortus.ai/doc/trust?rid=68cc0d7dc8e88b2e67d98932\&r=p97jk29hxd122x8d1cq4\&view=true}{27001} & \textbf{Certificado examinado}: alcance escriba TORTUS; expira \textbf{24/\allowbreak{}08/\allowbreak{}2028}. \\
ISO 13485 & \textbf{Scarlet NB UK}, SCAR-16.4.1.0; \href{https://trust.tortus.ai/doc/trust?rid=699582984a4bafd95d114b9a\&r=p97jk29hxd122x8d1cq4\&view=true}{13485} & \textbf{Certificado examinado}: software de documentación; expira \textbf{16/\allowbreak{}02/\allowbreak{}2029}. \\
Cyber Essentials Plus & \textbf{CyberSmart}; \href{https://trust.tortus.ai/doc/trust?rid=68cc0e3a1650a1721bb52f61\&r=p97jk29hxd122x8d1cq4\&view=true}{CE+} & \textbf{Certificado examinado}; recertificación \textbf{24/\allowbreak{}07/\allowbreak{}2027}. \\
UKCA clase IIa & \textbf{Scarlet, 8536}, según \href{https://tortus.ai/years-in-the-making}{UKCA} & Anunciado \textbf{24/\allowbreak{}06/\allowbreak{}2026}; verificación documental de UKCA pendiente. El catálogo conserva un registro clase I anterior. \\
\end{longtable}

\begin{itemize}
\item \textbf{Tratamiento de datos:} lista \textbf{Google Cloud, OpenAI, Auth0 y Datadog}; infraestructura Londres, endpoint OpenAI UE con retención cero y sin entrenamiento, Datadog UE. No demuestra residencia íntegra UE; transcripciones para depuración pueden conservarse hasta resolver incidencias. \href{https://tortus.ai/privacy}{Política del producto}.
\item \textbf{Evidencia de evaluación:} \emph{npj Digital Medicine} (2025), \textbf{25 consultas simuladas, 18 experimentos, 50 evaluadores}, con autores de Tortus. El \textbf{75\% menos de alucinaciones mayores} corresponde a cambios de prompt (4\ensuremath{\rightarrow}1), no a una certificación independiente del sistema Shell actual. \href{https://www.nature.com/articles/s41746-025-01670-7}{Artículo}.
\end{itemize}

\subsection*{16. Lyrebird}
\label{provider:16}

Documentación de referencia: \href{https://trust.lyrebirdhealth.com/}{Trust Center} y recursos; certificado ISO examinado, otros documentos con acceso o extracción limitada.

\begin{longtable}{L{\dimexpr 0.24\linewidth-2\tabcolsep\relax}L{\dimexpr 0.31\linewidth-2\tabcolsep\relax}L{\dimexpr 0.45\linewidth-2\tabcolsep\relax}}
\toprule
\textbf{Credencial} & \textbf{Emisor /\allowbreak{} documento} & \textbf{Alcance, vigencia y evidencia} \\
\midrule
\endfirsthead
\textbf{Credencial} & \textbf{Emisor /\allowbreak{} documento} & \textbf{Alcance, vigencia y evidencia} \\
\midrule
\endhead
\midrule
\multicolumn{3}{r}{\footnotesize Continúa en la página siguiente}
\endfoot
\bottomrule
\endlastfoot
ISO 27001 & \textbf{Certifi International}, A307RA001ISMS; \href{https://trust.lyrebirdhealth.com/doc/trust?rid=6a4b29f494ea8ceec985e54e\&r=5fwb6zo2bnid995vtydj6\&view=true}{Certificado ISO} & \textbf{Certificado examinado}: Subsidio Pty Ltd /\allowbreak{} Lyrebird, plataforma clínica; \textbf{26/\allowbreak{}06/\allowbreak{}2026–26/\allowbreak{}06/\allowbreak{}2029}. \\
Cyber Essentials /\allowbreak{} Plus & \href{https://trust.lyrebirdhealth.com/}{Trust Center} & \textbf{Declaradas}; PDF Plus localizado, contenido pendiente de verificación. \\
MHRA /\allowbreak{} UKCA clase I & \href{https://www.lyrebirdhealth.com/uk/uk-compliance}{compliance UK} & Registro/\allowbreak{}conformidad declarados; \textbf{no evaluación externa de clase IIa}. \\
\end{longtable}

\begin{itemize}
\item \textbf{Tratamiento de datos:} documentación UK: almacenamiento \textbf{UK}, inferencia transitoria \textbf{UE}, sin entrenamiento; audio eliminado tras transcribir (sin conexión, permanece en el dispositivo hasta reconectar); notas/\allowbreak{}transcripciones \textbf{siete días por defecto}, configurable. Subencargados en Trust Centre. La página de producto \texttt{/\allowbreak{}uk/\allowbreak{}} menciona Australia: \textbf{documentación regional inconsistente}. \href{https://www.lyrebirdhealth.com/uk/uk-compliance}{Compliance}, \href{https://help.lyrebirdhealth.com/en/articles/9229907-what-processes-does-lyrebird-use-to-protect-data}{retención}, \href{https://www.lyrebirdhealth.com/uk/product/clinical-notes}{producto}.
\item \textbf{Evidencia de evaluación:} \emph{BMC Health Services Research} (2026), estudio mixto de \textbf{16 semanas/\allowbreak{}7.499 consultas}; reporta errores/\allowbreak{}alucinaciones. Queensland Health identifica Lyrebird como herramienta evaluada. \href{https://pubmed.ncbi.nlm.nih.gov/41507920/}{Artículo}, \href{https://www.clinicalexcellence.qld.gov.au/improvement-exchange/enhancing-clinical-care-with-ambient-listening}{identificación del producto}.
\end{itemize}

\subsection*{17. Tali}
\label{provider:17}

Documentación de referencia: \href{https://tali.ai/security-privacy}{Seguridad}; SOC, PIA y pentest bajo solicitud, sin repositorio público específico localizado.

\begin{longtable}{L{\dimexpr 0.24\linewidth-2\tabcolsep\relax}L{\dimexpr 0.31\linewidth-2\tabcolsep\relax}L{\dimexpr 0.45\linewidth-2\tabcolsep\relax}}
\toprule
\textbf{Credencial} & \textbf{Emisor /\allowbreak{} documento} & \textbf{Alcance, vigencia y evidencia} \\
\midrule
\endfirsthead
\textbf{Credencial} & \textbf{Emisor /\allowbreak{} documento} & \textbf{Alcance, vigencia y evidencia} \\
\midrule
\endhead
\midrule
\multicolumn{3}{r}{\footnotesize Continúa en la página siguiente}
\endfoot
\bottomrule
\endlastfoot
SOC 2 tipo II & \href{https://tali.ai/security-privacy}{Seguridad} & \textbf{Declarado}; informe, auditor y periodo con verificación documental pendiente. \\
TGV: Québec & \href{https://www.quebec.ca/sante/professionnels/programme-transcription-intelligence-artificielle}{registro} & \textbf{Registro oficial consultado}, 09/\allowbreak{}07/\allowbreak{}2026: TGV sí; EFVP de Santé Québec no. \\
\end{longtable}

\begin{itemize}
\item \textbf{Tratamiento de datos:} Canadá si se selecciona esa región; otros países, EE. UU.; cuenta puede procesarse en EE. UU. Política: audio hasta \textbf{48 h}, notas/\allowbreak{}transcripciones \textbf{al menos 30 días}, desidentificación para mejoras. \textbf{Contradice el mensaje comercial de no conservar información clínica}. Proveedores de IA no identificados nominalmente. \href{https://tali.ai/privacy-policy}{Política}, \href{https://tali.ai/security-privacy}{seguridad}.
\item \textbf{Evidencia de evaluación:} \emph{Urology} (2025), comparación de cinco escribas con consultas simuladas. Tali y Nabla reciben mejores valoraciones, pero todos generan errores menores y críticos. \href{https://doi.org/10.1016/j.urology.2024.11.061}{Artículo}.
\end{itemize}

\subsection*{18. Scribeberry}
\label{provider:18}

Documentación de referencia: \href{https://trust.scribeberry.com/}{Trust Center} y documentos de seguridad; SOC bajo solicitud.

\begin{longtable}{L{\dimexpr 0.24\linewidth-2\tabcolsep\relax}L{\dimexpr 0.31\linewidth-2\tabcolsep\relax}L{\dimexpr 0.45\linewidth-2\tabcolsep\relax}}
\toprule
\textbf{Credencial} & \textbf{Emisor /\allowbreak{} documento} & \textbf{Alcance, vigencia y evidencia} \\
\midrule
\endfirsthead
\textbf{Credencial} & \textbf{Emisor /\allowbreak{} documento} & \textbf{Alcance, vigencia y evidencia} \\
\midrule
\endhead
\midrule
\multicolumn{3}{r}{\footnotesize Continúa en la página siguiente}
\endfoot
\bottomrule
\endlastfoot
SOC 2 tipo II & \href{https://scribeberry.com/product/security}{SOC}; \href{https://trust.scribeberry.com/}{Trust Center} & \textbf{Declarado}; informe bajo solicitud. \\
Evaluación de riesgos HIPAA/\allowbreak{}PIPEDA & \textbf{Ingrid Ruys}, 14/\allowbreak{}01/\allowbreak{}2024; \href{https://scribeberry.com/docs/security-risk-assessment.pdf}{evaluación firmada} & \textbf{Documento externo examinado}; no sustituye al SOC ni certifica HIPAA oficialmente. \\
TGV /\allowbreak{} EFVP: Québec & \href{https://www.quebec.ca/sante/professionnels/programme-transcription-intelligence-artificielle}{Québec} & \textbf{Registro oficial consultado}, 24/\allowbreak{}04/\allowbreak{}2026: ambos \textbf{no}. \\
\end{longtable}

\begin{itemize}
\item \textbf{Tratamiento de datos:} declara residencia \textbf{Canadá/\allowbreak{}EE. UU.} según región; Azure/\allowbreak{}Google Cloud y acuerdos con Microsoft, Google, OpenAI y Anthropic. Sin retención de audio ni entrenamiento con datos clínicos; notas con retención configurable. No acreditada residencia UE. \href{https://scribeberry.com/product/security}{Seguridad y documentos}.
\item \textbf{Evidencia de evaluación:} \emph{Healthcare} (2025), piloto prospectivo de Scribeberry V2 en paliativos, \textbf{dos residentes/\allowbreak{}19 casos}, resultados heterogéneos. La muestra reducida limita la generalización a otras versiones y contextos. \href{https://www.mdpi.com/2227-9032/13/17/2118}{Artículo}.
\end{itemize}

\subsection*{19. Quirónsalud: proyecto Scribe}
\label{provider:19}

Documentación de referencia: \href{https://www.quironsalud.com/es/comunicacion/actualidad/quironsalud-redisena-modelo-asistencial-reafirmandose-human}{Anuncio}; no se localizó Trust Center específico del proyecto Scribe.

\begin{longtable}{L{\dimexpr 0.24\linewidth-2\tabcolsep\relax}L{\dimexpr 0.31\linewidth-2\tabcolsep\relax}L{\dimexpr 0.45\linewidth-2\tabcolsep\relax}}
\toprule
\textbf{Credencial} & \textbf{Emisor /\allowbreak{} documento} & \textbf{Alcance, vigencia y evidencia} \\
\midrule
\endfirsthead
\textbf{Credencial} & \textbf{Emisor /\allowbreak{} documento} & \textbf{Alcance, vigencia y evidencia} \\
\midrule
\endhead
\midrule
\multicolumn{3}{r}{\footnotesize Continúa en la página siguiente}
\endfoot
\bottomrule
\endlastfoot
Certificación del sistema Scribe & \href{https://www.quironsalud.com/es/comunicacion/actualidad/quironsalud-redisena-modelo-asistencial-reafirmandose-human}{Anuncio} & \textbf{No localizada}. Noticia asistencial sin fabricante técnico, certificado ni oferta para terceros. \\
\end{longtable}

\begin{itemize}
\item \textbf{Tratamiento de datos:} anuncia confidencialidad, almacenamiento seguro y validación médica; no especifica país, subencargados, retención o entrenamiento. \href{https://www.quironsalud.com/es/comunicacion/actualidad/quironsalud-redisena-modelo-asistencial-reafirmandose-human}{Misma fuente}.
\item \textbf{Evidencia de evaluación:} \emph{International Journal of Medical Informatics} (2026): estudio retrospectivo de \textbf{48 hospitales españoles, 1.032.558 consultas con Scribe}, febrero de 2025–enero de 2026. Reporta menor duración y mejoras documentales. Revisado por pares, con autores del grupo; sin aleatorización. El resumen no desglosa rendimiento por idioma ni demuestra exactitud clínica de todas las notas. Publicación electrónica julio de 2026, número de noviembre. \href{https://pubmed.ncbi.nlm.nih.gov/42456601/}{PubMed}, \href{https://doi.org/10.1016/j.ijmedinf.2026.106584}{artículo}.
\item \textbf{Evidencia ambulatoria (actualización de 11/\allowbreak{}09/\allowbreak{}2026):} un segundo estudio analiza más de \textbf{2,33 millones de usos}. Duración media: \textbf{15,01 min con Scribe frente a 14,65 sin él}. El \textbf{87,4–89,2\% es similitud semántica entre borrador y texto corregido}, no exactitud clínica independiente. Encuestas transversales de baja respuesta; autores empleados por Quirón. Es evidencia española de implantación, sin validar todos los acentos ni demostrar HCIS o venta a terceros. Los usos de ambos estudios no deben sumarse como si correspondieran a pacientes independientes. \href{https://doi.org/10.3389/fdgth.2026.1874919}{Artículo y métodos}, ficha de lectura (archivo de notas del proyecto).
\end{itemize}

\subsection*{20. Freed}
\label{provider:20}

Documentación de referencia: \href{https://www.getfreed.ai/security}{Seguridad}; \href{https://app.drata.com/trust/1d1964c2-b90c-41ac-8d7c-2ea7f5af001c}{Trust Center} con contenido inaccesible durante la consulta; verificación del expediente pendiente.

\begin{longtable}{L{\dimexpr 0.24\linewidth-2\tabcolsep\relax}L{\dimexpr 0.31\linewidth-2\tabcolsep\relax}L{\dimexpr 0.45\linewidth-2\tabcolsep\relax}}
\toprule
\textbf{Credencial} & \textbf{Emisor /\allowbreak{} documento} & \textbf{Alcance, vigencia y evidencia} \\
\midrule
\endfirsthead
\textbf{Credencial} & \textbf{Emisor /\allowbreak{} documento} & \textbf{Alcance, vigencia y evidencia} \\
\midrule
\endhead
\midrule
\multicolumn{3}{r}{\footnotesize Continúa en la página siguiente}
\endfoot
\bottomrule
\endlastfoot
SOC 2 tipos I y II & \href{https://www.getfreed.ai/security}{Seguridad} & \textbf{Declarados}; informe, auditor y periodo con verificación documental pendiente. ISO/\allowbreak{}HITRUST o acreditación médica no localizados. \\
\end{longtable}

\begin{itemize}
\item \textbf{Tratamiento de datos:} \textbf{Azure EE. UU.}; audio temporal hasta completar resumen/\allowbreak{}comprobaciones y posterior borrado. Notas borrables manualmente o automáticamente a 30 días. \textbf{Sí entrena con notas desidentificadas}. Residencia UE y lista pública completa de subencargados pendientes de verificación. \href{https://www.getfreed.ai/security}{Seguridad}, \href{https://help.getfreed.ai/en/articles/9879953-privacy-policy-scribe-ai-chat}{privacidad}.
\item \textbf{Evidencia de evaluación:} artículo en \emph{Plastic and Reconstructive Surgery – Global Open} (2025), con experiencia y ejemplos de errores de Freed. Se trata de una \textbf{publicación narrativa y de recomendaciones}, sin diseño de validación clínica. \href{https://pmc.ncbi.nlm.nih.gov/articles/PMC11737491/}{Artículo}.
\end{itemize}

\subsection*{Proveedores adicionales del registro NHS}

El \href{https://digital.nhs.uk/services/ambient-scribing/ambient-voice-technology-self-certified-supplier-registry}{registro oficial}, actualizado el \textbf{02/\allowbreak{}09/\allowbreak{}2026}, contiene \textbf{28 proveedores}. Corti, Heidi, Lyrebird, Microsoft Dragon, Tandem y Tortus forman parte del conjunto inicial; los \textbf{22 restantes} se añaden a continuación. Es \textbf{autocertificación con evidencia presentada}, no aprobación clínica general del NHS. Los expedientes aportados al Procurement Hub no son públicos; su pertenencia al registro no se usa para inferir certificados concretos.

\subsection*{21. 33n: CLEARnotes}
\label{provider:21}

Documentación de referencia: \href{https://clearnotes.co.uk/product/}{producto} y \href{https://clearnotes.co.uk/}{seguridad en la web}; no se localizó Trust Center con certificados completos.

\begin{longtable}{L{\dimexpr 0.24\linewidth-2\tabcolsep\relax}L{\dimexpr 0.31\linewidth-2\tabcolsep\relax}L{\dimexpr 0.45\linewidth-2\tabcolsep\relax}}
\toprule
\textbf{Credencial} & \textbf{Emisor /\allowbreak{} documento} & \textbf{Alcance, vigencia y evidencia} \\
\midrule
\endfirsthead
\textbf{Credencial} & \textbf{Emisor /\allowbreak{} documento} & \textbf{Alcance, vigencia y evidencia} \\
\midrule
\endhead
\midrule
\multicolumn{3}{r}{\footnotesize Continúa en la página siguiente}
\endfoot
\bottomrule
\endlastfoot
Cyber Essentials Plus & \href{https://clearnotes.co.uk/}{Declaración de 33n} & Declara certificación; emisor, número y vigencia pendientes de verificación. \\
DCB0129 /\allowbreak{} DCB0160 /\allowbreak{} DSPT & \href{https://clearnotes.co.uk/}{Seguridad} & Cumplimiento declarado; no certificados clínicos externos. \\
\end{longtable}

\textbf{Tratamiento de datos y evaluación:} publica resultados de pilotos, pero esa página no aporta un estudio revisado por pares ni mapa completo de regiones/\allowbreak{}subencargados. \href{https://clearnotes.co.uk/clinically-led-deployment-and-compliance-the-key-to-sustainable-ai-integration-in-the-nhs/}{Despliegue y evaluación}.

\subsection*{22. Accuro: Assistant}
\label{provider:22}

Documentación de referencia: \href{https://accuro.co.uk/accuro-assistant-the-ambient-scribe-built-for-nhs-clinicians/}{Assistant} y \href{https://accuro.co.uk/sectors/healthcare/}{oferta sanitaria}; no se localizó catálogo público de certificados. Es Accuro UK, distinto del EHR canadiense Accuro.

\begin{longtable}{L{\dimexpr 0.24\linewidth-2\tabcolsep\relax}L{\dimexpr 0.31\linewidth-2\tabcolsep\relax}L{\dimexpr 0.45\linewidth-2\tabcolsep\relax}}
\toprule
\textbf{Credencial} & \textbf{Emisor /\allowbreak{} documento} & \textbf{Alcance, vigencia y evidencia} \\
\midrule
\endfirsthead
\textbf{Credencial} & \textbf{Emisor /\allowbreak{} documento} & \textbf{Alcance, vigencia y evidencia} \\
\midrule
\endhead
\midrule
\multicolumn{3}{r}{\footnotesize Continúa en la página siguiente}
\endfoot
\bottomrule
\endlastfoot
ISO 27001 & \href{https://accuro.co.uk/accuro-assistant-the-ambient-scribe-built-for-nhs-clinicians/}{Assistant} & Declarada para el servicio; certificado, emisor y vigencia con verificación documental pendiente. \\
Cyber Essentials Plus & \href{https://accuro.co.uk/sectors/healthcare/}{Seguridad} & Declarada; documento no localizado. HIPAA se trata como cumplimiento legal. \\
\end{longtable}

\textbf{Tratamiento de datos:} ofrece revisión por secretarios médicos en UK; no equivale a servicio sin acceso humano. Ubicaciones del procesamiento automático y subencargados completos pendientes. \href{https://accuro.co.uk/sectors/healthcare/}{Flujo del servicio}.

\subsection*{23. Accurx: Scribe}
\label{provider:23}

Documentación de referencia: \href{https://trust.assuric.com/accurx/public/compliance}{Trust Centre: certificados y recursos}. Publica certificados; riesgos clínicos y otros expedientes tienen acceso restringido.

\begin{longtable}{L{\dimexpr 0.24\linewidth-2\tabcolsep\relax}L{\dimexpr 0.31\linewidth-2\tabcolsep\relax}L{\dimexpr 0.45\linewidth-2\tabcolsep\relax}}
\toprule
\textbf{Credencial} & \textbf{Emisor /\allowbreak{} documento} & \textbf{Alcance, vigencia y evidencia} \\
\midrule
\endfirsthead
\textbf{Credencial} & \textbf{Emisor /\allowbreak{} documento} & \textbf{Alcance, vigencia y evidencia} \\
\midrule
\endhead
\midrule
\multicolumn{3}{r}{\footnotesize Continúa en la página siguiente}
\endfoot
\bottomrule
\endlastfoot
ISO 27001:2022 & \textbf{Amtivo /\allowbreak{} British Assessment Bureau}, 246818; \href{https://file.assuric.com/orgs\%2Faccurx\%2FcontrolEvidence\%2F1Jug9FI4p8olD3SUfIcN\%2FISO\_27001\_Certificate.pdf?alt=media\&token=cd8659cf-72a1-44ce-8d27-0604bdc6f4d8}{PDF} & \textbf{Documento examinado}: incluye explícitamente tecnología de voz ambiental; emitido 26/\allowbreak{}06/\allowbreak{}2026, expira \textbf{15/\allowbreak{}08/\allowbreak{}2027}. \\
Cyber Essentials /\allowbreak{} Plus & \href{https://trust.assuric.com/accurx/public/compliance}{Catálogo y PDF Plus 2025–26} & PDFs localizados; contenido y vigencia de Plus pendientes de verificación por limitaciones de legibilidad. \\
Clase I /\allowbreak{} IIa del escriba & \href{https://support.accurx.com/en/articles/768578-privacy-information-governance-and-security-an-overview}{Declaración de Accurx} & Anuncia ambas clases; certificado/\allowbreak{}región exactos pendientes de verificación para la oferta Accurx. \\
DTAC /\allowbreak{} DCB0129 /\allowbreak{} DSPT & \href{https://trust.assuric.com/accurx/public/compliance}{Trust Centre} & Documentación de cumplimiento; DTAC de plataforma v8, 18/\allowbreak{}08/\allowbreak{}2026. No aprobación clínica general. \\
\end{longtable}

\textbf{Tratamiento de datos y proveedores tecnológicos:} identifica \textbf{Tandem Health como subencargado del Scribe} y Azure; el audio se elimina tras transcribir. Las credenciales de Tandem deben comprobarse para la cadena contratada, no sumarse como certificados propios de Accurx. \href{https://trust.assuric.com/accurx/public}{Subencargados y descripción}.

\subsection*{24. Anathem}
\label{provider:24}

Documentación de referencia: \href{https://anathem.ai/medical-device}{página de aseguramiento y producto sanitario}; sin repositorio público de certificados localizado.

\begin{longtable}{L{\dimexpr 0.24\linewidth-2\tabcolsep\relax}L{\dimexpr 0.31\linewidth-2\tabcolsep\relax}L{\dimexpr 0.45\linewidth-2\tabcolsep\relax}}
\toprule
\textbf{Credencial} & \textbf{Emisor /\allowbreak{} documento} & \textbf{Alcance, vigencia y evidencia} \\
\midrule
\endfirsthead
\textbf{Credencial} & \textbf{Emisor /\allowbreak{} documento} & \textbf{Alcance, vigencia y evidencia} \\
\midrule
\endhead
\midrule
\multicolumn{3}{r}{\footnotesize Continúa en la página siguiente}
\endfoot
\bottomrule
\endlastfoot
ISO 27001 /\allowbreak{} Cyber Essentials Plus & \href{https://anathem.ai/medical-device}{Aseguramiento} & Declaradas; emisor, certificado y vigencia con verificación documental pendiente. \\
MHRA clase I & \href{https://anathem.ai/medical-device}{Página con enlace al registro} & Declara registro de ANATHEM LTD, sistema de dictado GMDN 36216; \textbf{IIa en desarrollo}. \\
DTAC /\allowbreak{} DSPT; pentest CREST & \href{https://anathem.ai/medical-device}{Aseguramiento} & Cumplimiento/\allowbreak{}pruebas declarados; sin verificación del informe externo. \\
\end{longtable}

\textbf{Tratamiento de datos y gobernanza de IA:} declara Azure UK, eliminación a diez días configurable, sin entrenamiento con datos sensibles. Describe pruebas con transcripciones sintéticas y trazabilidad; la documentación citada no aporta una evaluación independiente del mecanismo antialucinación. \href{https://anathem.ai/medical-device}{Misma documentación}.

\subsection*{25. Anima: Annie /\allowbreak{} Ambient Scribe}
\label{provider:25}

Documentación de referencia: \href{https://www.animahealth.com/security-compliance}{Security \& Compliance} y \href{https://www.animahealth.com/products/scribe}{Scribe}. El dominio animaapp.com corresponde a otro producto.

\begin{longtable}{L{\dimexpr 0.24\linewidth-2\tabcolsep\relax}L{\dimexpr 0.31\linewidth-2\tabcolsep\relax}L{\dimexpr 0.45\linewidth-2\tabcolsep\relax}}
\toprule
\textbf{Credencial} & \textbf{Emisor /\allowbreak{} documento} & \textbf{Alcance, vigencia y evidencia} \\
\midrule
\endfirsthead
\textbf{Credencial} & \textbf{Emisor /\allowbreak{} documento} & \textbf{Alcance, vigencia y evidencia} \\
\midrule
\endhead
\midrule
\multicolumn{3}{r}{\footnotesize Continúa en la página siguiente}
\endfoot
\bottomrule
\endlastfoot
Cyber Essentials /\allowbreak{} Plus & \href{https://www.animahealth.com/security-compliance}{Seguridad} & Declara Plus, pero una FAQ de la misma página lo mantiene «en curso»; vigencia \textbf{no resuelta}. \\
MHRA clase I & \href{https://www.animahealth.com/products/scribe}{Annie} & Declarada para el escriba; registro/\allowbreak{}certificado con verificación documental pendiente. \\
DTAC /\allowbreak{} DCB0129 /\allowbreak{} DSPT; IM1 & \href{https://www.animahealth.com/security-compliance}{Seguridad} & Expedientes bajo solicitud; IM1 es aseguramiento de interoperabilidad NHS, HCIS sin documentación específica identificada. \\
\end{longtable}

\textbf{Tratamiento de datos:} AWS Londres; DPA y DPIA específica del Scribe bajo solicitud. No se deduce la cadena de IA completa a partir del alojamiento de la plataforma. \href{https://www.animahealth.com/security-compliance}{Seguridad}.

\subsection*{26. Aprobrium: Lexacom}
\label{provider:26}

Documentación de referencia: \href{https://www.lexacom.co.uk/compliance-with-nhs-ai-guidance-and-assurance-standards/}{gobernanza} y \href{https://www.lexacom.co.uk/about-lexacom/}{certificaciones corporativas}.

\begin{longtable}{L{\dimexpr 0.24\linewidth-2\tabcolsep\relax}L{\dimexpr 0.31\linewidth-2\tabcolsep\relax}L{\dimexpr 0.45\linewidth-2\tabcolsep\relax}}
\toprule
\textbf{Credencial} & \textbf{Emisor /\allowbreak{} documento} & \textbf{Alcance, vigencia y evidencia} \\
\midrule
\endfirsthead
\textbf{Credencial} & \textbf{Emisor /\allowbreak{} documento} & \textbf{Alcance, vigencia y evidencia} \\
\midrule
\endhead
\midrule
\multicolumn{3}{r}{\footnotesize Continúa en la página siguiente}
\endfoot
\bottomrule
\endlastfoot
ISO 27001 /\allowbreak{} ISO 9001 & \href{https://www.lexacom.co.uk/about-lexacom/}{Página con imágenes de certificados} & Declara actualización agosto de 2024; contenido de las imágenes y vigencia con verificación documental pendiente. \\
Cyber Essentials Plus & \href{https://www.lexacom.co.uk/about-lexacom/}{Certificaciones} & Declara actualización diciembre de 2024; renovación vigente pendiente de verificación. \\
MHRA clase I & \href{https://www.lexacom.co.uk/compliance-with-nhs-ai-guidance-and-assurance-standards/}{Gobernanza} & Clase I bajo UK MDR; no certificado de evaluación externa IIa. \\
DTAC /\allowbreak{} DSPT /\allowbreak{} DCB0129; pentests CREST & \href{https://www.lexacom.co.uk/compliance-with-nhs-ai-guidance-and-assurance-standards/}{Gobernanza} & Cumplimiento y pruebas declarados; informes con verificación documental pendiente. \\
\end{longtable}

\textbf{Tratamiento de datos y gobernanza de IA:} Patient Shield elimina identificadores antes del procesamiento de IA, según proveedor; no implica por sí solo anonimización irreversible. \href{https://www.lexacom.co.uk/compliance-with-nhs-ai-guidance-and-assurance-standards/}{Gobernanza}.

\subsection*{27. Beam Up: Beam Notes /\allowbreak{} Magic Notes}
\label{provider:27}

Documentación de referencia: \href{https://trust.beam.org/}{Trust Portal}; políticas, informes y documentos de IA bajo solicitud.

\begin{longtable}{L{\dimexpr 0.24\linewidth-2\tabcolsep\relax}L{\dimexpr 0.31\linewidth-2\tabcolsep\relax}L{\dimexpr 0.45\linewidth-2\tabcolsep\relax}}
\toprule
\textbf{Credencial} & \textbf{Emisor /\allowbreak{} documento} & \textbf{Alcance, vigencia y evidencia} \\
\midrule
\endfirsthead
\textbf{Credencial} & \textbf{Emisor /\allowbreak{} documento} & \textbf{Alcance, vigencia y evidencia} \\
\midrule
\endhead
\midrule
\multicolumn{3}{r}{\footnotesize Continúa en la página siguiente}
\endfoot
\bottomrule
\endlastfoot
ISO 27001 /\allowbreak{} Cyber Essentials Plus & \href{https://trust.beam.org/}{Trust Portal} & Declaradas; titularidad corporativa, certificados y vigencias con verificación documental pendiente. \\
SOC 2 & \href{https://trust.beam.org/}{Trust Portal} & Mención e informe en catálogo; tipo, auditor y periodo pendientes de verificación. \\
CSA STAR nivel 1 /\allowbreak{} GovRAMP /\allowbreak{} IRAP & \href{https://trust.beam.org/}{Trust Portal} & Insignias: STAR 1 es autoevaluación; estatus y evaluaciones GovRAMP/\allowbreak{}IRAP pendientes de verificación. \\
DSPT & \href{https://www.dsptoolkit.nhs.uk/OrganisationSearch/U2A6H}{Registro NHS} & \textbf{Registro consultado}: 2025–26, estándares superados, 17/\allowbreak{}02/\allowbreak{}2026. \\
\end{longtable}

\textbf{Tratamiento de datos y evaluación:} oferta UK declara tratamiento UK/\allowbreak{}EEE; Nesta publicó una evaluación de aceptación social, \textbf{no validación de exactitud clínica}. \href{https://assets.applytosupply.digitalmarketplace.service.gov.uk/g-cloud-14/documents/717297/985901463553519-service-definition-document-2025-12-08-1626.pdf}{Ficha de servicio}, \href{https://www.nesta.org.uk/report/ai-sral-magicnotes/}{evaluación Nesta}.

\subsection*{28. Cerebric}
\label{provider:28}

Documentación de referencia: \href{https://trust.cerebric.io/public}{Trust Centre} y \href{https://trust.cerebric.io/public/compliance}{catálogo de cumplimiento}.

\begin{longtable}{L{\dimexpr 0.24\linewidth-2\tabcolsep\relax}L{\dimexpr 0.31\linewidth-2\tabcolsep\relax}L{\dimexpr 0.45\linewidth-2\tabcolsep\relax}}
\toprule
\textbf{Credencial} & \textbf{Emisor /\allowbreak{} documento} & \textbf{Alcance, vigencia y evidencia} \\
\midrule
\endfirsthead
\textbf{Credencial} & \textbf{Emisor /\allowbreak{} documento} & \textbf{Alcance, vigencia y evidencia} \\
\midrule
\endhead
\midrule
\multicolumn{3}{r}{\footnotesize Continúa en la página siguiente}
\endfoot
\bottomrule
\endlastfoot
Cyber Essentials /\allowbreak{} Plus & \href{https://trust.cerebric.io/public/compliance}{Catálogo con PDFs} & Declaradas; PDF Plus localizado, contenido no extraíble. Emisor/\allowbreak{}vigencia pendientes de verificación. \\
UKCA /\allowbreak{} MHRA clase I & \href{https://cerebric.io/product}{Producto} & Registro declarado para la plataforma; no evaluación externa IIa. \\
DTAC /\allowbreak{} DCB0129 /\allowbreak{} DSPT & \href{https://trust.cerebric.io/public/compliance}{Trust Centre} & Catálogo de expedientes; declara DSPT con estándares superados. \\
\end{longtable}

\textbf{Tratamiento de datos y gobernanza de IA:} identifica AWS y Neon; genera informes con referencias a fuentes. Regiones, retención y condiciones de entrenamiento completas pendientes de verificación. \href{https://trust.cerebric.io/public}{Subencargados}, \href{https://cerebric.io/product}{producto}.

\subsection*{29. CLAI Ltd: ClinicLetter.ai /\allowbreak{} Mayden}
\label{provider:29}

Documentación de referencia: \href{https://clinicletter.ai/}{página de producto y seguridad}; no se localizó Trust Center con informes descargables del escriba.

\begin{longtable}{L{\dimexpr 0.24\linewidth-2\tabcolsep\relax}L{\dimexpr 0.31\linewidth-2\tabcolsep\relax}L{\dimexpr 0.45\linewidth-2\tabcolsep\relax}}
\toprule
\textbf{Credencial} & \textbf{Emisor /\allowbreak{} documento} & \textbf{Alcance, vigencia y evidencia} \\
\midrule
\endfirsthead
\textbf{Credencial} & \textbf{Emisor /\allowbreak{} documento} & \textbf{Alcance, vigencia y evidencia} \\
\midrule
\endhead
\midrule
\multicolumn{3}{r}{\footnotesize Continúa en la página siguiente}
\endfoot
\bottomrule
\endlastfoot
Cyber Essentials Plus & \href{https://clinicletter.ai/}{Seguridad} & Declarada; certificado, emisor y vigencia pendientes de verificación. \\
MHRA clase I & \href{https://clinicletter.ai/}{Producto} & Registro declarado; no evaluación externa IIa. \\
DTAC /\allowbreak{} DSPT & \href{https://clinicletter.ai/}{Seguridad} & Cumplimiento declarado; no acreditación clínica externa. \\
\end{longtable}

\textbf{Tratamiento de datos:} anuncia residencia UK y borrado de sesiones completadas; cadena de subencargados/\allowbreak{}retención residual pendiente de verificación. \href{https://clinicletter.ai/}{Producto}.

\subsection*{30. Dentistry Dashboard: AI Notes}
\label{provider:30}

Documentación de referencia: \href{https://dentistrydashboard.com/trust-pack}{Trust Pack UK}; certificados e informes completos bajo solicitud.

\begin{longtable}{L{\dimexpr 0.24\linewidth-2\tabcolsep\relax}L{\dimexpr 0.31\linewidth-2\tabcolsep\relax}L{\dimexpr 0.45\linewidth-2\tabcolsep\relax}}
\toprule
\textbf{Credencial} & \textbf{Emisor /\allowbreak{} documento} & \textbf{Alcance, vigencia y evidencia} \\
\midrule
\endfirsthead
\textbf{Credencial} & \textbf{Emisor /\allowbreak{} documento} & \textbf{Alcance, vigencia y evidencia} \\
\midrule
\endhead
\midrule
\multicolumn{3}{r}{\footnotesize Continúa en la página siguiente}
\endfoot
\bottomrule
\endlastfoot
Cyber Essentials & \textbf{IASME}, según \href{https://dentistrydashboard.com/trust-pack}{Trust Pack} & Declara 11/\allowbreak{}09/\allowbreak{}2025, toda la organización; certificado con verificación documental pendiente. \textbf{No declara Plus propio}. \\
Pentest independiente & \textbf{Mondas Consulting}, según \href{https://dentistrydashboard.com/trust-pack}{Trust Pack} & Retest 10/\allowbreak{}04/\allowbreak{}2026 sin hallazgos pendientes en alcance; informe con verificación documental pendiente. \\
MHRA /\allowbreak{} UK MDR clase I & \href{https://dentistrydashboard.com/trust-pack}{Trust Pack} & Outlook Aesthetics Ltd, AI Notes, 20/\allowbreak{}09/\allowbreak{}2025; registro declarado. \\
ISO 27001 de infraestructura; DSPT/\allowbreak{}DTAC/\allowbreak{}DCB0129 & \href{https://dentistrydashboard.com/trust-pack}{Trust Pack} & \textbf{No acredita ISO propia}: corresponde a subencargados. Documentación NHS bajo solicitud. \\
\end{longtable}

\textbf{Tratamiento de datos:} AWS Londres/\allowbreak{}Azure UK South; audio transitorio, retención de notas configurable. Lista completa de subencargados bajo solicitud. \href{https://dentistrydashboard.com/trust-pack}{Trust Pack}.

\subsection*{31. Dictate IT: AVT en Swift /\allowbreak{} DIT3}
\label{provider:31}

Documentación de referencia: \href{https://dictateit.com/security-and-data-protection/}{biblioteca de seguridad} con certificados públicos.

\begin{longtable}{L{\dimexpr 0.24\linewidth-2\tabcolsep\relax}L{\dimexpr 0.31\linewidth-2\tabcolsep\relax}L{\dimexpr 0.45\linewidth-2\tabcolsep\relax}}
\toprule
\textbf{Credencial} & \textbf{Emisor /\allowbreak{} documento} & \textbf{Alcance, vigencia y evidencia} \\
\midrule
\endfirsthead
\textbf{Credencial} & \textbf{Emisor /\allowbreak{} documento} & \textbf{Alcance, vigencia y evidencia} \\
\midrule
\endhead
\midrule
\multicolumn{3}{r}{\footnotesize Continúa en la página siguiente}
\endfoot
\bottomrule
\endlastfoot
ISO 27001:2022 & \textbf{Amtivo /\allowbreak{} British Assessment Bureau}, 186273; \href{https://dictateit.com/wp-content/uploads/2025/04/Dictate.IT-ISO-27001-Certificate.pdf}{PDF} & \textbf{Documento examinado}: dictado, transcripción y reconocimiento de voz sanitario; expira \textbf{16/\allowbreak{}04/\allowbreak{}2027}. \\
ISO 9001:2015 & \textbf{Amtivo /\allowbreak{} British Assessment Bureau}, 185884; \href{https://dictateit.com/wp-content/uploads/2025/04/Dictate-IT-ISO-9001-Certificate.pdf}{PDF} & \textbf{Documento examinado}: mismo ámbito; expira \textbf{14/\allowbreak{}04/\allowbreak{}2027}. \\
Cyber Essentials Plus & \href{https://dictateit.com/wp-content/uploads/2024/09/Dictate-IT-Cyber-Essentials-Plus-Certificate.pdf}{PDF publicado} & \textbf{Documento examinado, vencido 12/\allowbreak{}07/\allowbreak{}2025}; excluye sistemas heredados en red separada. Renovación actual pendiente de verificación. \\
MHRA clase I; DTAC/\allowbreak{}DCB0129/\allowbreak{}DSPT & \href{https://dictateit.com/wp-content/uploads/2025/09/Dictate-IT-NHS-AVT-standards-compliance-summary.pdf}{Declaración AVT} & Declarados para AVT en Swift/\allowbreak{}DIT3; no certificado clínico externo. \\
\end{longtable}

\textbf{Tratamiento de datos:} declara almacenamiento/\allowbreak{}procesamiento UK, sin proveedores externos de IA; pentest anual NCC Group. Es declaración de julio de 2025, no auditoría actual de toda la arquitectura. \href{https://dictateit.com/wp-content/uploads/2025/09/Dictate-IT-NHS-AVT-standards-compliance-summary.pdf}{Documento AVT}.

\subsection*{32. Docwise}
\label{provider:32}

Documentación de referencia: \href{https://www.docwise.co.uk/compliance}{Compliance}: enlaces a registros; DTAC/\allowbreak{}DCB0129 bajo solicitud. El análisis corresponde al escriba de docwise.co.uk; excluye empresas homónimas de gestión documental.

\begin{longtable}{L{\dimexpr 0.24\linewidth-2\tabcolsep\relax}L{\dimexpr 0.31\linewidth-2\tabcolsep\relax}L{\dimexpr 0.45\linewidth-2\tabcolsep\relax}}
\toprule
\textbf{Credencial} & \textbf{Emisor /\allowbreak{} documento} & \textbf{Alcance, vigencia y evidencia} \\
\midrule
\endfirsthead
\textbf{Credencial} & \textbf{Emisor /\allowbreak{} documento} & \textbf{Alcance, vigencia y evidencia} \\
\midrule
\endhead
\midrule
\multicolumn{3}{r}{\footnotesize Continúa en la página siguiente}
\endfoot
\bottomrule
\endlastfoot
Cyber Essentials & \href{https://www.docwise.co.uk/compliance}{Compliance, enlace a registro} & Declarada; certificado, emisor y vigencia con verificación documental pendiente. \\
UKCA /\allowbreak{} MHRA clase I & \href{https://www.docwise.co.uk/compliance}{Compliance} & Marcado/\allowbreak{}registro declarados; no evaluación externa IIa. \\
DTAC /\allowbreak{} DCB0129 /\allowbreak{} DSPT & \href{https://www.docwise.co.uk/compliance}{Compliance} & Declara DSPT «standards met»; documentos clínicos bajo solicitud. \\
\end{longtable}

\textbf{Tratamiento de datos y evaluación:} declara servidores UK; lista completa de subencargados pendiente de verificación. La web presenta resultados de piloto sin referencia bibliográfica que permita verificar la revisión por pares. \href{https://www.docwise.co.uk/compliance}{Seguridad}, \href{https://www.docwise.co.uk/}{piloto}.

\subsection*{33. eConsult: Hi Scribe}
\label{provider:33}

Documentación de referencia: \href{https://econsult.net/primary-care/security}{seguridad corporativa} y \href{https://econsult.net/hi-scribe-privacy-policy}{política específica Hi Scribe}.

\begin{longtable}{L{\dimexpr 0.24\linewidth-2\tabcolsep\relax}L{\dimexpr 0.31\linewidth-2\tabcolsep\relax}L{\dimexpr 0.45\linewidth-2\tabcolsep\relax}}
\toprule
\textbf{Credencial} & \textbf{Emisor /\allowbreak{} documento} & \textbf{Alcance, vigencia y evidencia} \\
\midrule
\endfirsthead
\textbf{Credencial} & \textbf{Emisor /\allowbreak{} documento} & \textbf{Alcance, vigencia y evidencia} \\
\midrule
\endhead
\midrule
\multicolumn{3}{r}{\footnotesize Continúa en la página siguiente}
\endfoot
\bottomrule
\endlastfoot
ISO 27001 & \textbf{Alcumus ISOQAR}, según \href{https://info.econsult.health/hubfs/eConsult\%20Health\%20-\%20GDPR\%20and\%20Data\%20Protection\%20and\%20Information\%20Security\%20Compliance\%20Statement\%20-\%202024.pdf}{declaración 2024} & Certificado corporativo 20623 citado; vigencia actual y alcance explícito Hi Scribe pendientes de verificación. \\
DSPT & \href{https://econsult.net/hi-scribe-privacy-policy}{Política Hi Scribe} & Declara estándares superados; no certificación clínica. \\
\end{longtable}

\textbf{Tratamiento de datos:} política marzo de 2026: Google Speech-to-Text y AWS Comprehend en UK; Vertex AI en UE; CloudWatch/\allowbreak{}Datadog para metadatos. Audio no almacenado, texto en sesión; feedback voluntario hasta dos semanas. \href{https://econsult.net/hi-scribe-privacy-policy}{Mapa y retención}.

\subsection*{34. Epro: Bluewire Technologies}
\label{provider:34}

Documentación de referencia: \href{https://epro.com/blog/epro-has-been-awarded-with-the-iso27001-accreditation-the-internationally-recognised-standard-for-information-security-management-system/}{anuncio ISO}, \href{https://epro.com/}{producto} y política; no se localizó Trust Center de certificados completos.

\begin{longtable}{L{\dimexpr 0.24\linewidth-2\tabcolsep\relax}L{\dimexpr 0.31\linewidth-2\tabcolsep\relax}L{\dimexpr 0.45\linewidth-2\tabcolsep\relax}}
\toprule
\textbf{Credencial} & \textbf{Emisor /\allowbreak{} documento} & \textbf{Alcance, vigencia y evidencia} \\
\midrule
\endfirsthead
\textbf{Credencial} & \textbf{Emisor /\allowbreak{} documento} & \textbf{Alcance, vigencia y evidencia} \\
\midrule
\endhead
\midrule
\multicolumn{3}{r}{\footnotesize Continúa en la página siguiente}
\endfoot
\bottomrule
\endlastfoot
ISO 27001 & \href{https://epro.com/blog/epro-has-been-awarded-with-the-iso27001-accreditation-the-internationally-recognised-standard-for-information-security-management-system/}{Anuncio, 24/\allowbreak{}10/\allowbreak{}2025} & Declarada; emisor, número, alcance y vigencia pendientes de verificación. \\
ISO 9001 & \href{https://www.linkedin.com/company/epro\_health/}{Declaración corporativa} & Declarada por Epro; certificado con verificación documental pendiente. \\
MHRA clase I /\allowbreak{} DCB0129 & \href{https://epro.com/}{Producto} & Conformidad declarada; no evaluación externa IIa. \\
\end{longtable}

\textbf{Tratamiento de datos:} alojamiento cloud anunciado; la política web no basta para establecer países, retención y subencargados del AVT. \href{https://epro.com/product/ai-speech/policy}{Política}.

\subsection*{35. G2 Speech: SpeechAmbient}
\label{provider:35}

Documentación de referencia: \href{https://www.g2speech.com/security}{Security \& Compliance}; no se localizó repositorio público de certificados completos.

\begin{longtable}{L{\dimexpr 0.24\linewidth-2\tabcolsep\relax}L{\dimexpr 0.31\linewidth-2\tabcolsep\relax}L{\dimexpr 0.45\linewidth-2\tabcolsep\relax}}
\toprule
\textbf{Credencial} & \textbf{Emisor /\allowbreak{} documento} & \textbf{Alcance, vigencia y evidencia} \\
\midrule
\endfirsthead
\textbf{Credencial} & \textbf{Emisor /\allowbreak{} documento} & \textbf{Alcance, vigencia y evidencia} \\
\midrule
\endhead
\midrule
\multicolumn{3}{r}{\footnotesize Continúa en la página siguiente}
\endfoot
\bottomrule
\endlastfoot
ISO 27001 /\allowbreak{} ISO 27701 & \href{https://www.g2speech.com/security}{Seguridad} & Declaradas; certificado, emisor y vigencia con verificación documental pendiente. \\
NEN 7510 /\allowbreak{} Cyber Essentials Plus & \href{https://www.g2speech.com/security}{Seguridad} & Declaradas; alcance concreto de SpeechAmbient pendiente. \\
MHRA clase I & \href{https://www.g2speech.com/security}{Seguridad} & Registro declarado; no evaluación externa IIa. \\
DCB0129 /\allowbreak{} DSPT /\allowbreak{} DTAC & \href{https://www.g2speech.com/security}{Seguridad} & Cumplimiento; no certificaciones clínicas independientes. \\
\end{longtable}

\textbf{Tratamiento de datos:} declara alojamiento UK o UE, sin conservar audio y acceso de soporte solicitado por el cliente. Región concreta y lista completa de subencargados pendientes. \href{https://www.g2speech.com/security}{Seguridad}.

\subsection*{36. HealthOrbit AI}
\label{provider:36}

Documentación de referencia: \href{https://healthorbit.ai/trust-center/}{Trust Center}: expediente completo, DPA, DPIA y riesgos clínicos bajo solicitud.

\begin{longtable}{L{\dimexpr 0.24\linewidth-2\tabcolsep\relax}L{\dimexpr 0.31\linewidth-2\tabcolsep\relax}L{\dimexpr 0.45\linewidth-2\tabcolsep\relax}}
\toprule
\textbf{Credencial} & \textbf{Emisor /\allowbreak{} documento} & \textbf{Alcance, vigencia y evidencia} \\
\midrule
\endfirsthead
\textbf{Credencial} & \textbf{Emisor /\allowbreak{} documento} & \textbf{Alcance, vigencia y evidencia} \\
\midrule
\endhead
\midrule
\multicolumn{3}{r}{\footnotesize Continúa en la página siguiente}
\endfoot
\bottomrule
\endlastfoot
ISO 27001 /\allowbreak{} SOC 2 tipo II & \href{https://healthorbit.ai/trust-center/}{Trust Center} & Declaradas; emisor, alcance y vigencia/\allowbreak{}periodo con verificación documental pendiente. \\
Cyber Essentials /\allowbreak{} Plus & \href{https://healthorbit.ai/trust-center/}{Trust Center} & Insignias publicadas; certificados con verificación documental pendiente. \\
Dubai AI Seal & \textbf{DCAI}, según \href{https://healthorbit.ai/news/healthorbit-ai-officially-recognised-with-dubai-ai-seal/}{anuncio} & Anunciado 07/\allowbreak{}11/\allowbreak{}2025; número/\allowbreak{}categoría oficial pendientes de verificación. No certificación médica. \\
UKCA /\allowbreak{} MOHAP & \href{https://healthorbit.ai/trust-center/}{Trust Center} & Logotipos sin documento, clase o autorización identificados; no se cuentan como aprobaciones verificadas. MOHAP es el Ministerio de Salud y Prevención de EAU, no el nombre de una certificación. \\
DTAC /\allowbreak{} DCB0129 /\allowbreak{} DSPT & \href{https://healthorbit.ai/trust-center/}{Trust Center} & Cumplimiento declarado; FAQ aún describe DTAC completo en preparación. \\
\end{longtable}

\textbf{Tratamiento de datos:} identifica AWS y Google Workspace, pero la lista completa requiere solicitud. No verificada residencia íntegra UE. \href{https://healthorbit.ai/trust-center/}{Subencargados}.

\subsection*{37. Medhasyn Ltd: Clinical Companion}
\label{provider:37}

Documentación de referencia: \href{https://www.medhasyn.com/}{sitio oficial} y \href{https://digital.nhs.uk/services/ambient-scribing/ambient-voice-technology-self-certified-supplier-registry}{registro NHS}.

\begin{longtable}{L{\dimexpr 0.24\linewidth-2\tabcolsep\relax}L{\dimexpr 0.31\linewidth-2\tabcolsep\relax}L{\dimexpr 0.45\linewidth-2\tabcolsep\relax}}
\toprule
\textbf{Credencial} & \textbf{Emisor /\allowbreak{} documento} & \textbf{Alcance, vigencia y evidencia} \\
\midrule
\endfirsthead
\textbf{Credencial} & \textbf{Emisor /\allowbreak{} documento} & \textbf{Alcance, vigencia y evidencia} \\
\midrule
\endhead
\midrule
\multicolumn{3}{r}{\footnotesize Continúa en la página siguiente}
\endfoot
\bottomrule
\endlastfoot
Certificación externa del escriba & \href{https://www.medhasyn.com/}{Web} & \textbf{No localizada}: información pública limitada, sin certificados, informes ni Trust Center identificable. \\
Registro AVT NHS & \href{https://digital.nhs.uk/services/ambient-scribing/ambient-voice-technology-self-certified-supplier-registry}{NHS} & Pertenencia verificada; expediente no público, no permite inferir ISO/\allowbreak{}CE/\allowbreak{}otros. \\
\end{longtable}

\textbf{Tratamiento de datos:} alude a flujos locales cuando se necesitan; no demuestra una configuración contratada, regiones, subencargados o retención. \href{https://www.medhasyn.com/}{Web}.

\subsection*{38. Optum: EMIS Scribe}
\label{provider:38}

Documentación de referencia: \href{https://www.emishealth.com/news-insights/march-2026-newsletter}{producto} y \href{https://www.emishealth.com/terms-conditions}{condiciones/\allowbreak{}subencargados}; no se localizó expediente público de certificación específico del Scribe.

\begin{longtable}{L{\dimexpr 0.24\linewidth-2\tabcolsep\relax}L{\dimexpr 0.31\linewidth-2\tabcolsep\relax}L{\dimexpr 0.45\linewidth-2\tabcolsep\relax}}
\toprule
\textbf{Credencial} & \textbf{Emisor /\allowbreak{} documento} & \textbf{Alcance, vigencia y evidencia} \\
\midrule
\endfirsthead
\textbf{Credencial} & \textbf{Emisor /\allowbreak{} documento} & \textbf{Alcance, vigencia y evidencia} \\
\midrule
\endhead
\midrule
\multicolumn{3}{r}{\footnotesize Continúa en la página siguiente}
\endfoot
\bottomrule
\endlastfoot
Certificaciones propias del Scribe & \href{https://www.emishealth.com/terms-conditions}{Condiciones} & \textbf{No verificadas}. No se trasladan las credenciales de otros productos de Optum ni las del proveedor T-Pro. \\
Registro AVT NHS & \href{https://digital.nhs.uk/services/ambient-scribing/ambient-voice-technology-self-certified-supplier-registry}{NHS} & Pertenencia verificada; no certificación independiente del producto. \\
\end{longtable}

\textbf{Tratamiento de datos y proveedores tecnológicos:} \textbf{T-Pro aporta la IA de EMIS Scribe}; la tabla contractual indica que ese servicio no procesa fuera de UK. No localizados plazos de retención ni mapa técnico completo del escriba. \href{https://www.emishealth.com/terms-conditions}{Subencargados}.

\subsection*{39. Pungo: Joy /\allowbreak{} JoyNotes}
\label{provider:39}

Documentación de referencia: \href{https://trust.assuric.com/the-joy-app/public/compliance}{Trust Centre}: ISO público, expedientes DCB/\allowbreak{}DTAC/\allowbreak{}DPA con acceso restringido.

\begin{longtable}{L{\dimexpr 0.24\linewidth-2\tabcolsep\relax}L{\dimexpr 0.31\linewidth-2\tabcolsep\relax}L{\dimexpr 0.45\linewidth-2\tabcolsep\relax}}
\toprule
\textbf{Credencial} & \textbf{Emisor /\allowbreak{} documento} & \textbf{Alcance, vigencia y evidencia} \\
\midrule
\endfirsthead
\textbf{Credencial} & \textbf{Emisor /\allowbreak{} documento} & \textbf{Alcance, vigencia y evidencia} \\
\midrule
\endhead
\midrule
\multicolumn{3}{r}{\footnotesize Continúa en la página siguiente}
\endfoot
\bottomrule
\endlastfoot
ISO 27001:2022 & \textbf{Citation ISO Certification}, 363532021; \href{https://file.assuric.com/orgs\%2Fthe-joy-app\%2FcontrolEvidence\%2F22pghJ4mEH7v7DPByRHz\%2FISO\%2027001\%20\_\%202022\%20Certificate.pdf?alt=media\&token=bbae1827-a31f-48e6-ac52-fd720c6b4fc3}{PDF} & \textbf{Documento examinado}: software de prescripción social; expira \textbf{05/\allowbreak{}01/\allowbreak{}2027}. JoyNotes no nombrado expresamente. \\
Cyber Essentials /\allowbreak{} Plus & \href{https://trust.assuric.com/the-joy-app/public/compliance}{Catálogo} & CE público localizado; contenido no legible. Plus figuraba en una versión indexada anterior; vigencia actual pendiente de verificación. \\
MHRA clase I: JoyNotes & \href{https://trust.assuric.com/the-joy-app/public/compliance}{Catálogo con enlace MHRA} & Registro declarado; declaración de conformidad con verificación documental pendiente. \\
DTAC /\allowbreak{} DCB0129 /\allowbreak{} DSPT & \href{https://trust.assuric.com/the-joy-app/public/compliance}{Trust Centre} & Recursos específicos JoyNotes; no certificación clínica externa. \\
\end{longtable}

\textbf{Tratamiento de datos:} DPA y DPIA identificados en catálogo, con verificación íntegra pendiente; regiones, subencargados y entrenamiento del escriba pendientes. \href{https://trust.assuric.com/the-joy-app/public/compliance}{Recursos}.

\subsection*{40. Scribetech: Augnito Omni}
\label{provider:40}

Documentación de referencia: \href{https://scribetech.co.uk/press-releases/nhs-ready-ai-medical-scribe-augnito-omni-ai-among-first-to-fully-meet-new-nhs-england-ambient-voice-tech-guidelines}{aseguramiento Omni} y documentación corporativa; sin Trust Center de certificados completos localizado.

\begin{longtable}{L{\dimexpr 0.24\linewidth-2\tabcolsep\relax}L{\dimexpr 0.31\linewidth-2\tabcolsep\relax}L{\dimexpr 0.45\linewidth-2\tabcolsep\relax}}
\toprule
\textbf{Credencial} & \textbf{Emisor /\allowbreak{} documento} & \textbf{Alcance, vigencia y evidencia} \\
\midrule
\endfirsthead
\textbf{Credencial} & \textbf{Emisor /\allowbreak{} documento} & \textbf{Alcance, vigencia y evidencia} \\
\midrule
\endhead
\midrule
\multicolumn{3}{r}{\footnotesize Continúa en la página siguiente}
\endfoot
\bottomrule
\endlastfoot
ISO 27001 & \textbf{BSI}, según \href{https://scribetech.co.uk/wp-content/uploads/2022/11/ProductComparison.pdf}{documentación Scribetech} & Declarada; certificado/\allowbreak{}vigencia actuales y alcance exacto Omni con verificación documental pendiente. \\
ISO 9001 /\allowbreak{} ISO 14001 /\allowbreak{} Cyber Essentials Plus & \href{https://scribetech.co.uk/press-releases/nhs-ready-ai-medical-scribe-augnito-omni-ai-among-first-to-fully-meet-new-nhs-england-ambient-voice-tech-guidelines}{Anuncio Omni, 26/\allowbreak{}07/\allowbreak{}2025} & Declaradas para la plataforma; certificados y vigencias con verificación documental pendiente. \\
SOC 2 & \href{https://scribetech.co.uk/blog/2023/06/26/the-realities-of-successfully-integrating-natural-speech-recognition-into-mental-health-services}{Artículo corporativo 2023} & Menciona conformidad; tipo, auditor y cobertura Omni no demostrados. \\
UKCA /\allowbreak{} MHRA clase I; DCB0129/\allowbreak{}DCB0160/\allowbreak{}DSPT & \href{https://scribetech.co.uk/press-releases/nhs-ready-ai-medical-scribe-augnito-omni-ai-among-first-to-fully-meet-new-nhs-england-ambient-voice-tech-guidelines}{Anuncio Omni} & Registro/\allowbreak{}cumplimiento declarados; \textbf{IIa era objetivo}, no credencial obtenida en esa fuente. \\
\end{longtable}

\textbf{Tratamiento de datos:} declara procesamiento UK y retención opcional dentro del país; entrenamiento previo con datos clínicos UK no aclara reutilización contractual de nuevas consultas. \href{https://scribetech.co.uk/press-releases/nhs-ready-ai-medical-scribe-augnito-omni-ai-among-first-to-fully-meet-new-nhs-england-ambient-voice-tech-guidelines}{Omni}.

\subsection*{41. T-Pro: Scribe /\allowbreak{} Copilot}
\label{provider:41}

El \href{https://trust.tpro.io/}{Trust Center} presentó limitaciones de acceso al contenido. Fuentes alternativas: \href{https://www.applytosupply.digitalmarketplace.service.gov.uk/g-cloud/services/310198615907338}{ficha G-Cloud} y \href{https://info.tpro.io/hubfs/T-Pro\%20Scribe\%20with\%20Compliance\%20with\%20NHS\%20England\%20AVT\%20Requirements.pdf?hsLang=en}{declaración AVT}.

\begin{longtable}{L{\dimexpr 0.24\linewidth-2\tabcolsep\relax}L{\dimexpr 0.31\linewidth-2\tabcolsep\relax}L{\dimexpr 0.45\linewidth-2\tabcolsep\relax}}
\toprule
\textbf{Credencial} & \textbf{Emisor /\allowbreak{} documento} & \textbf{Alcance, vigencia y evidencia} \\
\midrule
\endfirsthead
\textbf{Credencial} & \textbf{Emisor /\allowbreak{} documento} & \textbf{Alcance, vigencia y evidencia} \\
\midrule
\endhead
\midrule
\multicolumn{3}{r}{\footnotesize Continúa en la página siguiente}
\endfoot
\bottomrule
\endlastfoot
ISO 27001 & \textbf{SGS UK}, según \href{https://www.applytosupply.digitalmarketplace.service.gov.uk/g-cloud/services/310198615907338}{ficha presentada por T-Pro} & Fecha citada 07/\allowbreak{}11/\allowbreak{}2022; dictado/\allowbreak{}transcripción/\allowbreak{}correspondencia. Verificación del certificado vigente pendiente. \\
Cyber Essentials Plus & \href{https://info.tpro.io/hubfs/T-Pro\%20Scribe\%20with\%20Compliance\%20with\%20NHS\%20England\%20AVT\%20Requirements.pdf?hsLang=en}{Declaración AVT, 18/\allowbreak{}06/\allowbreak{}2025} & Renovación anual declarada; documento/\allowbreak{}auditor pendientes de verificación. \\
ISO 13485 /\allowbreak{} ISO 14971 & \href{https://www.linkedin.com/posts/tpro\_healthcareai-ambientvoicetechnology-clinicaldocumentation-activity-7467515379280457728-9-we}{Comunicación de T-Pro} & Referencias a estándares; \textbf{certificación externa no demostrada} en las fuentes consultadas. \\
MHRA clase I; DTAC/\allowbreak{}DCB0129/\allowbreak{}DSPT & \href{https://info.tpro.io/hubfs/T-Pro\%20Scribe\%20with\%20Compliance\%20with\%20NHS\%20England\%20AVT\%20Requirements.pdf?hsLang=en}{Declaración AVT} & En junio de 2025 el registro médico estaba \textbf{en curso}; estado posterior pendiente de verificación. Cumplimiento NHS declarado. \\
\end{longtable}

\textbf{Tratamiento de datos:} DPIA y flujos por cliente; el documento no permite establecer residencia UE y cadena completa para España. La ISO 9001 de material de \textbf{InventAsia} no se atribuye automáticamente a T-Pro Scribe. \href{https://info.tpro.io/hubfs/T-Pro\%20Scribe\%20with\%20Compliance\%20with\%20NHS\%20England\%20AVT\%20Requirements.pdf?hsLang=en}{AVT}, \href{https://info.tpro.io/hubfs/PDFs/T-Pro\%20Compliance\%20-\%20Invent\%20Asia.pdf?hsLang=en}{alcance InventAsia}.

\subsection*{42. X-on Health: Surgery Intellect, powered by TORTUS}
\label{provider:42}

Documentación de referencia: \href{https://www.x-on.co.uk/trust-centre/}{Trust Centre} y \href{https://surgeryconnect.academy/trust-centre/}{Academy: documentos}. Certificados públicos; parte del expediente clínico requiere autenticación.

\begin{longtable}{L{\dimexpr 0.24\linewidth-2\tabcolsep\relax}L{\dimexpr 0.31\linewidth-2\tabcolsep\relax}L{\dimexpr 0.45\linewidth-2\tabcolsep\relax}}
\toprule
\textbf{Credencial} & \textbf{Emisor /\allowbreak{} documento} & \textbf{Alcance, vigencia y evidencia} \\
\midrule
\endfirsthead
\textbf{Credencial} & \textbf{Emisor /\allowbreak{} documento} & \textbf{Alcance, vigencia y evidencia} \\
\midrule
\endhead
\midrule
\multicolumn{3}{r}{\footnotesize Continúa en la página siguiente}
\endfoot
\bottomrule
\endlastfoot
ISO 9001:2015 & \textbf{Citation}, 136442021; \href{https://surgeryconnect.academy/wp-content/uploads/2026/05/ISO-9001\_2015-Certificate.pdf}{PDF} & \textbf{Documento examinado}: navegación asistencial y soluciones digitales sanitarias; expira \textbf{02/\allowbreak{}08/\allowbreak{}2027}. \\
ISO 22301:2019 & \textbf{Citation}, 427522023; \href{https://surgeryconnect.academy/wp-content/uploads/2026/05/ISO-22301\_2019-Certificate.pdf}{PDF} & \textbf{Documento examinado}, mismo ámbito; expira \textbf{13/\allowbreak{}09/\allowbreak{}2027}. \\
ISO 14001:2015 & \textbf{Citation}, 429932023; \href{https://surgeryconnect.academy/wp-content/uploads/2026/05/ISO-14001\_2015-Certificate.pdf}{PDF} & \textbf{Documento examinado}, mismo ámbito; expira \textbf{17/\allowbreak{}10/\allowbreak{}2027}. \\
ISO 27001 /\allowbreak{} ISO 42001 & \href{https://surgeryconnect.academy/wp-content/uploads/2026/07/ISO27001-Certificate.pdf}{27001}, \href{https://surgeryconnect.academy/wp-content/uploads/2026/08/ISO42001-Certificate.pdf}{42001} & Certificados públicos localizados; sin extracción legible. Emisor, alcance y vigencia pendientes de verificación. \\
Cyber Essentials /\allowbreak{} Plus & \href{https://surgeryconnect.academy/trust-centre/}{Academy} & Declara Plus; enlaza documento titulado CE octubre 2025–octubre 2026, con verificación íntegra pendiente. \\
UKCA clase IIa: tecnología TORTUS & \href{https://www.x-on.co.uk/surgery-intellect/}{Surgery Intellect} & Declaración explícita de la tecnología incorporada; sin verificación de un certificado independiente de X-on. \\
DTAC /\allowbreak{} DCB0129 /\allowbreak{} DSPT & \href{https://www.x-on.co.uk/trust-centre/}{Trust Centre} & Expedientes por producto; cumplimiento, no certificación clínica externa. \\
\end{longtable}

\textbf{Tratamiento de datos:} política corporativa declara UK con excepciones contractuales; para el escriba es necesario el expediente \textbf{X-on + TORTUS}, no solo el de telefonía. \href{https://www.x-on.co.uk/privacy-notice/}{Privacidad}, \href{https://surgeryconnect.academy/trust-centre/}{DPIA y cadena del producto}.

\subsection*{43. AWS: HealthScribe}
\label{provider:43}

\textbf{Revisión específica: 17/\allowbreak{}09/\allowbreak{}2026.} API gestionada para construir aplicaciones clínicas: transcripción y borradores de notas, por lotes o streaming. \textbf{Solo inglés de EE. UU. y región US East (N. Virginia)} según la guía; sin conector HCIS identificado. \href{https://aws.amazon.com/es/healthscribe/}{Producto}, \href{https://docs.aws.amazon.com/transcribe/latest/dg/health-scribe.html}{documentación técnica}.

\textbf{Interpretación del alcance:} AWS presenta HealthScribe dentro de Amazon Transcribe. Sus catálogos incluyen Transcribe y extienden la cobertura a funciones de disponibilidad general, salvo exclusión, con evaluación en la siguiente oportunidad de auditoría. La aplicación a HealthScribe se infiere de esa regla cuando no se nombra expresamente; \textbf{no demuestra certificación de la aplicación de un partner} ni revisión de cada versión posterior. \href{https://aws.amazon.com/compliance/services-in-scope/ENS-High/}{Regla y alcance ENS}.

\begin{longtable}{L{\dimexpr 0.24\linewidth-2\tabcolsep\relax}L{\dimexpr 0.31\linewidth-2\tabcolsep\relax}L{\dimexpr 0.45\linewidth-2\tabcolsep\relax}}
\toprule
\textbf{Credencial} & \textbf{Emisor /\allowbreak{} documento} & \textbf{Alcance, vigencia y evidencia} \\
\midrule
\endfirsthead
\textbf{Credencial} & \textbf{Emisor /\allowbreak{} documento} & \textbf{Alcance, vigencia y evidencia} \\
\midrule
\endhead
\midrule
\multicolumn{3}{r}{\footnotesize Continúa en la página siguiente}
\endfoot
\bottomrule
\endlastfoot
\textbf{ISO 27001:2022} & \textbf{EY CertifyPoint}, 2013-009; \href{https://d1.awsstatic.com/onedam/marketing-channels/website/aws/en\_US/certification/compliance/iso\_27001\_global\_certification.pdf}{certificado público, 7 páginas} & \textbf{Documento examinado}: emitido 25/\allowbreak{}11/\allowbreak{}2025, reexpedido 24/\allowbreak{}08/\allowbreak{}2026; expira 30/\allowbreak{}11/\allowbreak{}2028. Su p. 2 remite al catálogo de servicios, que incluye Transcribe; p. 4 incluye Virginia. \\
\textbf{ENS ALTA} & \href{https://aws.amazon.com/compliance/services-in-scope/ENS-High/}{Catálogo AWS de servicios en alcance} & Transcribe figura incluido. Cobertura de HealthScribe inferida según la regla anterior; \textbf{certificado completo, emisor, vigencia y anexos pendientes de examen} mediante AWS Artifact. No se clasifica como ENS ausente ni como expediente íntegramente verificado. \\
\textbf{SOC 1 /\allowbreak{} 2 /\allowbreak{} 3} & \textbf{Ernst \& Young LLP} para el \href{https://d1.awsstatic.com/onedam/marketing-channels/website/aws/en\_US/whitepapers/compliance/AWS\_SOC3\_Report.pdf}{SOC 3 público}; \href{https://aws.amazon.com/compliance/services-in-scope/SOC/}{catálogo SOC} & \textbf{SOC 3 examinado}, periodo 01/\allowbreak{}04/\allowbreak{}2025–31/\allowbreak{}03/\allowbreak{}2026, firmado 29/\allowbreak{}05/\allowbreak{}2026; Transcribe en p. 7. El auditor excluye expresamente evaluar el rendimiento/\allowbreak{}integridad de la IA. SOC 1/\allowbreak{}2: informes restringidos, no examinados. \\
ISO 27017 /\allowbreak{} 27018 /\allowbreak{} 27701; 9001 /\allowbreak{} 22301 /\allowbreak{} 20000-1; CSA STAR & \href{https://aws.amazon.com/compliance/iso-certified/}{Catálogo ISO de AWS} & Certificaciones declaradas y Transcribe incluido; documentos individuales y alcance detallado pendientes. No se atribuye a todos el emisor o vigencia del certificado 27001. \\
\textbf{ISO 42001:2023} & \textbf{Schellman Compliance LLC}, según \href{https://aws.amazon.com/blogs/machine-learning/aws-achieves-iso-iec-420012023-artificial-intelligence-management-system-accredited-certification/}{anuncio AWS}; \href{https://aws.amazon.com/compliance/iso-42001-faqs/}{acceso al certificado} & Anuncio de 2024 incluye Transcribe. Certificado vigente y anexos no examinados; gestión de IA, no validación clínica de HealthScribe. \\
C5 /\allowbreak{} HITRUST CSF & \href{https://aws.amazon.com/compliance/services-in-scope/C5/}{Alcance C5}, \href{https://aws.amazon.com/compliance/services-in-scope/HITRUST-CSF/}{alcance HITRUST} & Transcribe incluido en ambos catálogos. Informes/\allowbreak{}certificados completos, periodos y modalidad HITRUST pendientes; sin atribuir tipo C5 específico. \\
HIPAA: elegibilidad & \href{https://aws.amazon.com/compliance/hipaa-eligible-services-reference/}{Listado AWS} & Mención explícita a Transcribe \textbf{incluido HealthScribe}. Requiere BAA y configuración adecuada; \textbf{no es certificación externa HIPAA}. \\
\end{longtable}

\textbf{Datos y privacidad:} AWS declara no entrenar con entradas/\allowbreak{}salidas ni retenerlas tras entregar resultados; las copias en S3 quedan bajo gestión del cliente. DPA y lista de subencargados públicos, sin mapa específico completo de HealthScribe identificado. La disponibilidad en Virginia \textbf{no acredita residencia UE}; almacenar la salida en otro lugar tampoco cambia dónde se procesa. Los términos contemplan inferencia entre regiones para servicios basados en Bedrock, incluido HealthScribe: requiere aclaración contractual del flujo efectivo. \href{https://aws.amazon.com/healthscribe/faqs/}{FAQ}, \href{https://d1.awsstatic.com/legal/aws-dpa/aws-dpa.pdf}{DPA}, \href{https://aws.amazon.com/compliance/sub-processors/}{subencargados}, \href{https://aws.amazon.com/service-terms/}{términos §§1.24 y 50.3}.

\textbf{Calidad y seguridad clínica:} la AI Service Card, aplicable a la versión de \textbf{28/\allowbreak{}11/\allowbreak{}2023}, publica evaluación propia: F1 de reconocimiento de palabras \ensuremath{\geq}84\% en 28 grupos y F1 de resumen 72\% en conversaciones guionizadas. No son porcentajes de notas clínicamente correctas ni resultados en español. Incluye vínculos de las notas a la transcripción y sus limitaciones. No se localizó validación clínica específica revisada por pares en la búsqueda dirigida. \href{https://docs.aws.amazon.com/ai/responsible-ai/aws-healthscribe/overview.html}{Service Card}.

\textbf{Partners:} AWS identifica \textbf{3M HIS/\allowbreak{}Solventum, Netsmart, ScribeEMR, TeleTracking y Pariveda}. Netsmart documenta \textbf{Bells Virtual Scribe con HealthScribe}; las otras referencias combinan productos, colaboraciones y desarrollo a medida. La relación con 3M no demuestra que Fluency Align España utilice esta API. \href{https://aws.amazon.com/healthscribe/customers/}{Clientes AWS}, \href{https://aws.amazon.com/blogs/industries/netsmart-transforms-behavioral-healthcare-with-new-feature-of-aws-healthscribe/}{implementación Netsmart}. El detalle comercial y las páginas revisadas constan en la nota de revisión AWS (archivo de notas del proyecto).

\textbf{Distinción necesaria:} \href{https://docs.aws.amazon.com/transcribe/latest/dg/transcribe-medical.html}{Amazon Transcribe Medical} ofrece transcripción médica, también limitada a inglés de EE. UU.; no sustituye la generación de notas de HealthScribe. Las garantías de una aplicación terminada requieren evaluar además al integrador. \textbf{Sin inclusión en la selección corta actual:} español, procesamiento UE e integración HCIS no acreditados.
\endgroup

% !TEX root = ../main.tex
% GENERADO por scripts/sync_market_appendices.py; editar el Markdown de origen.
\clearpage
\begin{landscape}
\section{Interoperabilidad con Dedalus HCIS}
\label{app:hcis}
\begingroup\footnotesize

Revisión pública inicial: \textbf{9 de septiembre de 2026}; AWS HealthScribe añadido el \textbf{17/\allowbreak{}09/\allowbreak{}2026}. El análisis comprende los \textbf{43 proveedores} del \hyperref[app:proveedores]{documento de certificaciones}, incluidos los 22 adicionales del registro NHS. Las otras 42 filas conservan su fecha de revisión.

\textbf{En este informe, HCIS designa Dedalus Care HCIS /\allowbreak{} HP-HCIS}, no cualquier historia clínica electrónica ni otros productos de Dedalus. \href{https://www.dedalus.com/spain/es/our-offer/products/hcis-can-be-implemented-using-a-best-of-breed-approach/}{Producto HCIS}.

\textbf{Sí} = compatibilidad específica con HCIS publicada por el proveedor; se especifican el carácter comercial de la declaración y el producto cubierto. \textbf{No} = integración disponible con HCIS sin evidencia identificada en las fuentes consultadas, o indisponibilidad declarada por el proveedor. Esta clasificación no implica imposibilidad técnica. Las APIs, las interfaces HL7/\allowbreak{}FHIR, la integración con otros EHR y la transferencia manual de texto no acreditan un conector HCIS.

\textbf{Resultado principal:} Invox anuncia dictado dentro de HCIS; \textbf{no queda confirmada esa integración para Genesis, su escriba ambiental}. Tandem tiene páginas dedicadas a HCIS, pero ambas indican que la integración aún no está disponible. La revisión pública no ha confirmado un escriba ambiental de esta lista con incorporación operativa de notas en HCIS. Las fuentes y límites están en las filas 2 y 5.

Las páginas enlazadas son las consultadas para esta revisión; se indica cuándo el acceso se limitó al contenido indexado. El alcance se limita a la documentación pública y excluye pruebas funcionales de conectores y expedientes de acceso privado. Las vías descritas corresponden a lo publicado por cada proveedor.

\begin{longtable}{L{\dimexpr 0.18\linewidth-2\tabcolsep\relax}L{\dimexpr 0.08\linewidth-2\tabcolsep\relax}L{\dimexpr 0.48\linewidth-2\tabcolsep\relax}L{\dimexpr 0.26\linewidth-2\tabcolsep\relax}}
\toprule
\textbf{Proveedor /\allowbreak{} producto} & \textbf{HCIS documentado} & \textbf{Modalidad de integración y alcance de la evidencia} & \textbf{Fuentes documentales} \\
\midrule
\endfirsthead
\textbf{Proveedor /\allowbreak{} producto} & \textbf{HCIS documentado} & \textbf{Modalidad de integración y alcance de la evidencia} & \textbf{Fuentes documentales} \\
\midrule
\endhead
\midrule
\multicolumn{4}{r}{\footnotesize Continúa en la página siguiente}
\endfoot
\bottomrule
\endlastfoot
1. Noa Notes: Doctoralia & \textbf{No} & Oferta hospitalaria con integración en HCE, sin identificar HCIS ni documentar un conector específico. & \href{https://noa.ai/es-es/ia-para-hospitales}{Hospitales}, \href{https://noa.ai/es-mx/noa-notes}{Noa Notes}. \\
2. Invox Medical & \textbf{Sí: dictado} & \textbf{Dictation} anuncia transcripción directa y cumplimentación de campos en HCIS. \textbf{Genesis: no confirmado}. La página comercial usa hospitales de ejemplo «XXX/\allowbreak{}YYY»; no aporta un despliegue identificable. & \href{https://dictation.invoxmedical.com/software-medico/dedalus-hcis/}{HCIS: contenido indexado; apertura directa bloqueada}, \href{https://dictation.invoxmedical.com/software-medico/}{catálogo compatible}, \href{https://www.invoxmedical.com/productos/generador-informes-medicos-inteligencia-artificial}{Genesis}. \\
3. Heidi: Scribe /\allowbreak{} Remote & \textbf{No} & Embedded incorpora la interfaz; Connect conecta datos del EHR. Capacidades y sistemas dependen del conector; HCIS no identificado. & \href{https://www.heidihealth.com/en-us/integrations}{Catálogo}, \href{https://support.heidihealth.com/en/articles/14075160-heidi-integrations-overview}{guía de integraciones}. \\
4. Solventum: Fluency Align & \textbf{No} & Anuncia trabajo dentro del flujo de la HCE y combina documentación ambiental con dictado. No se localizó compatibilidad HCIS específica de Align. & \href{https://www.solventum.com/es-es/home/health-information-technology/solutions/fluency-align/}{Align España}, \href{https://www.solventum.com/en-us/home/health-information-technology/solutions/fluency-align/}{Align EE. UU.}. \\
5. Tandem & \textbf{No} & Las dos fichas HCIS dicen \textbf{«No disponible»} y \textbf{«No, todavía no»}. Ofrece solicitar desarrollo/\allowbreak{}configuración; el texto sobre transferencia de notas es condicional. & \href{https://tandemhealth.ai/es/producto/integraciones/orbis-hcis}{Orbis /\allowbreak{} HCIS}, \href{https://tandemhealth.ai/es/producto/integraciones/hp-his-hp-hcis}{HP-HIS /\allowbreak{} HP-HCIS}. \\
6. Microsoft: Dragon Copilot & \textbf{No} & APIs REST para socios: contexto de paciente/\allowbreak{}episodio, audio y devolución de notas. Permiten construir una integración; no documentan HCIS. & \href{https://learn.microsoft.com/en-us/industry/healthcare/dragon-copilot/sdk/partner-apis/}{API para socios}, \href{https://www.microsoft.com/en-us/microsoft-cloud/blog/healthcare/2025/03/03/microsoft-partners-are-shaping-the-future-of-healthcare-with-microsoft-dragon-copilot/}{ecosistema de socios}. \\
7. Nabla & \textbf{No} & Nabla Connect y vías de integración en la interfaz del EHR; no se localizó una implementación identificada para HCIS. & \href{https://docs.nabla.com/guides/ehr-integration}{Guía técnica}, \href{https://www.nabla.com/press-release/nabla-launches-nabla-connect}{Nabla Connect}. \\
8. Abridge & \textbf{No} & Publica integración con Epic, Oracle Health y athenahealth; HCIS sin documentación específica identificada. & \href{https://www.abridge.com/press-release/patient-centered-clinician-intelligence-platform-keynote}{Plataforma e integraciones anunciadas}, \href{https://www.abridge.com/}{producto}. \\
9. Ambience Healthcare & \textbf{No} & Publica anuncios de integración con Epic, Oracle/\allowbreak{}Cerner y athenahealth; no se identificó HCIS. & \href{https://www.ambiencehealthcare.com/blog}{Catálogo de anuncios de integración}, \href{https://www.ambiencehealthcare.com/}{producto}. \\
10. Suki & \textbf{No} & Conectores bidireccionales para varios EHR y API/\allowbreak{}SDK para integrar funciones de documentación. HCIS no figura en las fuentes consultadas. & \href{https://www.suki.ai/ehr-integrations/}{Integraciones EHR}, \href{https://developer.suki.ai/api-overview/overview}{API}. \\
11. DeepScribe & \textbf{No} & Integraciones específicas, incluidas Epic, iKnowMed y OncoEMR en oncología. Sus condiciones distinguen servicio independiente de integración contratada. & \href{https://www.deepscribe.ai/ehr-integrations}{Integraciones}, \href{https://www.deepscribe.ai/resources/ambient-ai-and-oncology-value-chain-reduce-complexity}{oncología}, \href{https://www.deepscribe.ai/terms}{condiciones}. \\
12. Commure: Ambient AI & \textbf{No} & API, extensión Chrome y conectores EHR; describe lectura/\allowbreak{}escritura para sistemas identificados, sin HCIS. & \href{https://docs.ambient.commure.com/get\_started/overview}{Documentación técnica}, \href{https://www.commure.com/ambient-ai}{Ambient AI}. \\
13. Corti & \textbf{No} & API para incorporar transcripción/\allowbreak{}documentación a aplicaciones. Su aplicación a HCIS requeriría un desarrollo específico; las fuentes consultadas no documentan una integración disponible. & \href{https://docs.corti.ai/get\_started/welcome}{API}, \href{https://corti.ai/enterprise}{Enterprise}. \\
14. Oracle: Clinical AI Agent & \textbf{No} & Documentación integrada en \textbf{Oracle Health EHR}; no se localizó un conector con Dedalus HCIS. & \href{https://www.oracle.com/health/clinical-suite/clinical-ai-agent/}{Clinical AI Agent}, \href{https://www.oracle.com/a/ocom/docs/oracle-health-ehr-product-brochure.pdf}{ficha Oracle Health EHR}. \\
15. Tortus & \textbf{No} & Publica integración en EHR del NHS y caso con Epic; también se incorpora a Surgery Intellect. No se identificó HCIS. & \href{https://tortus.ai/}{Producto y casos}, \href{https://www.x-on.co.uk/surgery-intellect/}{Surgery Intellect}. \\
16. Lyrebird & \textbf{No} & Partner API, FHIR y webhooks; sincronización de citas y devolución de notas según integración. HCIS no identificado. & \href{https://www.lyrebirdhealth.com/us/integrations}{Integraciones}, \href{https://help.lyrebirdhealth.com/en/collections/10480349-other-integrations}{guías de otros EHR}. \\
17. Tali & \textbf{No} & Conectores documentados para EHR principalmente canadienses y uso superpuesto a otras aplicaciones. No hay prueba específica de HCIS. & \href{https://help.tali.ai/en/articles/how-to-set-up-your-tali-ehr-integration}{Configuración EHR}, \href{https://tali.ai/faq}{FAQ}. \\
18. Scribeberry & \textbf{No} & API/\allowbreak{}SDK, widget y extensión para EMR. La compatibilidad genérica con aplicaciones web no acredita funcionamiento en HCIS. & \href{https://console.scribeberry.com/docs}{Documentación para desarrolladores}, \href{https://help.scribeberry.com/emr-integrations/5r5f7zLEpjW2DEW9bnaAaW}{integraciones EMR}. \\
19. Quirónsalud: Scribe & \textbf{No} & Proyecto asistencial interno; el anuncio no identifica HCIS ni una interfaz comercial para terceros. & \href{https://www.quironsalud.com/es/comunicacion/actualidad/quironsalud-redisena-modelo-asistencial-reafirmandose-human}{Anuncio de Quirónsalud}. \\
20. Freed & \textbf{No} & Documenta instalación de Vim para eClinicalWorks y athenahealth; no se encontró HCIS. & \href{https://www.getfreed.ai/help/vim-install-guide}{Instalación Vim}, \href{https://help.getfreed.ai/en/collections/9672396-integrations}{catálogo de integraciones}. \\
21. 33n: CLEARnotes & \textbf{No} & Describe incorporación al flujo clínico, sin documentación pública de un conector HCIS. & \href{https://clearnotes.co.uk/product/}{Producto}, \href{https://clearnotes.co.uk/}{sitio oficial}. \\
22. Accuro: Assistant & \textbf{No} & Ofrece entrega de documentos y opciones de integración EPR, sin identificar HCIS. Es Accuro UK, no el EHR canadiense homónimo. & \href{https://accuro.co.uk/sectors/healthcare/}{Healthcare}, \href{https://accuro.co.uk/accuro-assistant-the-ambient-scribe-built-for-nhs-clinicians/}{Assistant}. \\
23. Accurx: Scribe & \textbf{No} & Scribe usa Tandem. La plataforma Accurx publica compatibilidad con sistemas de atención primaria UK; no demuestra HCIS ni idéntico alcance para cada función de Scribe. & \href{https://www.accurx.com/bundles}{Oferta e integraciones de plataforma}, \href{https://trust.assuric.com/accurx/public}{Trust Center y dependencia}. \\
24. Anathem & \textbf{No} & Anuncia compatibilidad HL7/\allowbreak{}FHIR y conversaciones para integrar EHR. No equivale a un conector HCIS disponible. & \href{https://anathem.ai/medical-device}{Medical device e interoperabilidad}. \\
25. Anima: Annie & \textbf{No} & Describe lectura de contexto y envío de notas/\allowbreak{}códigos a EMIS y SystmOne. HCIS sin documentación específica identificada. & \href{https://www.animahealth.com/products/scribe}{Scribe}, \href{https://www.animahealth.com/security-compliance}{seguridad e IM1}. \\
26. Aprobrium: Lexacom & \textbf{No} & Integración IM1; el registro NHS identifica Lexacom con Optum/\allowbreak{}EMIS. Esto no acredita el módulo ambiental en HCIS. & \href{https://www.lexacom.co.uk/compliance-with-nhs-ai-guidance-and-assurance-standards/}{Gobernanza}, \href{https://digital.nhs.uk/services/digital-services-for-integrated-care/im1-pairing-integration}{registro IM1}. \\
27. Beam Up: Beam Notes /\allowbreak{} Magic Notes & \textbf{No} & En el caso público de Newcastle, Magic Notes no se integra con otros sistemas del ayuntamiento. No se identificó otro caso HCIS. & \href{https://www.gov.uk/algorithmic-transparency-records/newcastle-city-council-magic-notes}{Ficha oficial de Newcastle}, \href{https://trust.beam.org/}{Trust Portal}. \\
28. Cerebric & \textbf{No} & Extensión Chrome/\allowbreak{}Edge con guía de importación desde Odoo; oferta modular e integración a medida. Sin compatibilidad HCIS documentada en las fuentes consultadas. & \href{https://cerebric.io/extension-guide}{Guía de extensión}, \href{https://cerebric.io/product}{producto}. \\
29. CLAI: ClinicLetter.ai & \textbf{No} & Integración mediante MPage en Oracle Health Millennium: contexto del paciente y devolución de documentación. No es HCIS. & \href{https://clinicletter.ai/}{Producto e integración}. \\
30. Dentistry Dashboard & \textbf{No} & Afirma funcionar con cualquier software de gestión dental, sin prueba específica de HCIS ni conector hospitalario documentado. & \href{https://dentistrydashboard.com/}{Producto}, \href{https://dentistrydashboard.com/trust-pack}{Trust Pack}. \\
31. Dictate IT: Swift /\allowbreak{} DIT3 & \textbf{No} & DIT3: HL7 mediante NextGen Connect, FHIR e interfaces a medida. Swift: APIs de socios EMIS/\allowbreak{}TPP. El documento AVT no nombra HCIS. & \href{https://dictateit.com/wp-content/uploads/2025/09/Dictate-IT-NHS-AVT-standards-compliance-summary.pdf}{Documento técnico AVT}, \href{https://www.emishealth.com/partners/dictation-software/dictate-swift}{ficha Swift en EMIS}. \\
32. Docwise & \textbf{No} & Anuncia compatibilidad con EHR de forma general; no se localizó catálogo técnico o conector HCIS. & \href{https://www.docwise.co.uk/}{Producto}, \href{https://www.docwise.co.uk/compliance}{Compliance}. \\
33. eConsult: Hi Scribe & \textbf{No} & eConsult Toolbar con EMIS Web/\allowbreak{}SystmOne: recibe contexto básico y archiva el resumen aprobado por el profesional. HCIS sin documentación específica identificada. & \href{https://econsult.net/hi-scribe-privacy-policy}{Política técnica de Hi Scribe}, \href{https://digital.nhs.uk/services/digital-services-for-integrated-care/im1-pairing-integration}{registro IM1}. \\
34. Epro: Bluewire & \textbf{No} & Ofrece desarrollo de APIs e integración a medida; Digital Outpatient publica HL7/\allowbreak{}MESH y formatos heredados. Alcance HCIS del escriba no identificado. & \href{https://epro.com/services/}{Servicios de integración}, \href{https://epro.com/digital-outpatient/}{Digital Outpatient}. \\
35. G2 Speech: SpeechAmbient & \textbf{No} & No se localizó una integración HCIS específica en la ficha del producto y los anuncios revisados. & \href{https://www.g2speech.com/speechambient}{SpeechAmbient}, \href{https://www.g2speech.com/news-uk}{noticias}. \\
36. HealthOrbit AI & \textbf{No} & Describe FHIR, HL7 y componentes de integración a medida; HCIS no identificado. & \href{https://healthorbit.ai/trust-center/}{Trust Center, FAQ de integración}. \\
37. Medhasyn: Clinical Companion & \textbf{No} & Información técnica pública mínima; no documenta conectores HCIS. La pertenencia al registro NHS no acredita integración. & \href{https://www.medhasyn.com/}{Sitio oficial}, \href{https://digital.nhs.uk/services/ambient-scribing/ambient-voice-technology-self-certified-supplier-registry}{registro NHS}. \\
38. Optum: EMIS Scribe & \textbf{No} & Integrado en EMIS-X/\allowbreak{}EMIS Web; T-Pro aporta la IA. No se publica conector HCIS del producto EMIS Scribe. & \href{https://www.emishealth.com/news-insights/march-2026-newsletter}{Producto, marzo de 2026}, \href{https://www.emishealth.com/terms-conditions}{condiciones y T-Pro}. \\
39. Pungo: JoyNotes & \textbf{No} & Publica envío de notas a EMIS/\allowbreak{}SystmOne; HCIS sin documentación específica identificada. & \href{https://www.thejoyapp.com/products/joynotes}{JoyNotes}, \href{https://support.thejoyapp.com/en/articles/10471051-how-to-send-notes-made-on-joy-to-a-client-s-clinical-record-from-the-client-record}{guía de envío al registro clínico}. \\
40. Scribetech: Augnito Omni & \textbf{No} & Oferta de Voice Services mediante API/\allowbreak{}SDK junto a Omni y Spectra; no se localizó un conector HCIS de Omni. & \href{https://scribetech.co.uk/}{Productos y Voice Services}, \href{https://scribetech.co.uk/press-releases/nhs-ready-ai-medical-scribe-augnito-omni-ai-among-first-to-fully-meet-new-nhs-england-ambient-voice-tech-guidelines}{descripción Omni para NHS}. \\
41. T-Pro: Scribe /\allowbreak{} Copilot & \textbf{No} & Documento AVT describe APIs, HL7/\allowbreak{}FHIR y varios EHR (incluidos Cerner, Epic y EMIS), sin identificar HCIS. & \href{https://info.tpro.io/hubfs/T-Pro\%20Scribe\%20with\%20Compliance\%20with\%20NHS\%20England\%20AVT\%20Requirements.pdf?hsLang=en}{Interoperabilidad en documento AVT}. \\
42. X-on: Surgery Intellect & \textbf{No} & Combina Surgery Connect con Tortus y envío al registro mediante IM1 para atención primaria UK. Sin compatibilidad HCIS identificada en las fuentes consultadas. & \href{https://www.x-on.co.uk/surgery-intellect/}{Producto e integración}, \href{https://www.x-on.co.uk/trust-centre/}{Trust Centre}. \\
43. AWS: HealthScribe & \textbf{No} & API de transcripción y notas para construir aplicaciones; no conector HCIS identificado. Inglés de EE. UU. y región Virginia. La integración EHR genérica de partners no acredita HCIS. Transcribe Medical solo aporta transcripción. & \href{https://docs.aws.amazon.com/transcribe/latest/dg/health-scribe.html}{HealthScribe}, \href{https://docs.aws.amazon.com/transcribe/latest/dg/transcribe-medical.html}{Transcribe Medical}, \href{https://aws.amazon.com/healthscribe/customers/}{clientes/\allowbreak{}partners}, \href{https://www.ntst.com/carefabric/careguidance-solutions/ai-documentation-assistant}{Bells de Netsmart}. \\
\end{longtable}

\textbf{Referencia adicional, fuera de los 43 proveedores:} Nunsys y Dedalus anuncian que \textbf{Vekai se integra dentro de Dedalus Care HCIS}. Es una referencia específica y relevante, pero el anuncio no permite confirmar que incluya captura ambiental de consultas ni un despliegue completo de ese módulo. \href{https://www.nunsys.com/nunsys-group-y-dedalus-se-alian-para-impulsar-desde-la-comunitat-un-nuevo-modelo-de-historia-clinica-digital/}{Anuncio de Nunsys}.

La verificación de las vías de integración propuestas requiere especificar \textbf{versión/\allowbreak{}módulo de HCIS, lectura de paciente y episodio, y escritura de la nota aprobada}. La demostración en la instalación del hospital debe distinguir la integración automatizada, el dictado en campos, la transferencia manual y las funcionalidades pendientes de desarrollo.
\endgroup
\end{landscape}

% !TEX root = ../main.tex
% GENERADO por scripts/sync_market_appendices.py; editar el Markdown de origen.
\clearpage
\begin{landscape}
\section{Condiciones económicas y transparencia de la facturación}
\label{app:precios}
\begingroup\footnotesize

Revisión inicial: \textbf{9 de septiembre de 2026}; AWS HealthScribe añadido el \textbf{17/\allowbreak{}09/\allowbreak{}2026}, sin revisar de nuevo los demás precios. El análisis comprende los 43 proveedores incluidos en \hyperref[app:proveedores]{certificaciones} e \hyperref[app:hcis]{integraciones HCIS}. \hyperref[sec:candidatos]{Evaluación comparativa}.

\textbf{Alcance de la comparación:} precio publicado del producto, unidad facturable, compromiso, límites y extras. Se prioriza España; las ofertas de otros países y las tarifas históricas se identifican expresamente. No se convierten monedas ni se equipara una API con un producto terminado. «No localizado» no implica gratuidad. \textbf{Ningún importe de esta tabla demuestra el coste de una implantación ambiental con HCIS.}

\textbf{Transparencia de la oferta publicada (criterio de esta revisión):} \textbf{A}, importe, unidad, periodo y principales límites identificables; \textbf{B}, información parcial, importes ambiguos o tarifas antiguas; \textbf{C}, sin tarifa ni mecanismo suficientemente concreto para calcular; \textbf{N/\allowbreak{}A}, sin oferta comercial identificada. La categoría A describe la transparencia de la tarifa; no implica un coste reducido, ausencia de cargos adicionales o disponibilidad de un presupuesto hospitalario completo. Los fallos de acceso se indican y no se interpretan como ocultación deliberada.

Precios sin normalizar impuestos: la inclusión o exclusión del IVA solo se indica cuando la fuente lo especifica. \textbf{\$ sin código} conserva el símbolo de la fuente; no debe convertirse a euros sin confirmar la moneda. Los equivalentes mensuales de contratos anuales no implican un derecho de cancelación mensual.

\begin{longtable}{L{\dimexpr 0.15\linewidth-2\tabcolsep\relax}L{\dimexpr 0.2\linewidth-2\tabcolsep\relax}L{\dimexpr 0.36\linewidth-2\tabcolsep\relax}L{\dimexpr 0.1\linewidth-2\tabcolsep\relax}L{\dimexpr 0.19\linewidth-2\tabcolsep\relax}}
\toprule
\textbf{Proveedor} & \textbf{Precio público documentado} & \textbf{Facturación, límites y condiciones pendientes} & \textbf{Transparencia} & \textbf{Fuentes documentales} \\
\midrule
\endfirsthead
\textbf{Proveedor} & \textbf{Precio público documentado} & \textbf{Facturación, límites y condiciones pendientes} & \textbf{Transparencia} & \textbf{Fuentes documentales} \\
\midrule
\endhead
\midrule
\multicolumn{5}{r}{\footnotesize Continúa en la página siguiente}
\endfoot
\bottomrule
\endlastfoot
1. Noa Notes & \textbf{España: 30 €/\allowbreak{}mes + IVA}. & Notes independiente, notas ilimitadas. Doctoralia Pro y organizaciones: oferta personalizada. Número de usuarios, permanencia y coste de integración pendientes de confirmación. Los \textbf{950 MXN/\allowbreak{}mes} aparecen en condiciones promocionales de 2025; no son una tarifa española vigente. & \textbf{B} & \href{https://noa.ai/es-es/precio}{Precio España}, \href{https://noa.ai/es-mx/precio}{precio México actual}, \href{https://noa.ai/es-mx/terminos-y-condiciones-noa-notes-oferta-2-meses}{condiciones de 2025}. \\
2. Invox Genesis & \textbf{No localizado.} & Propuesta comercial previa demostración. La documentación pública consultada no especifica precio por usuario o minuto, mínimos, implantación ni mantenimiento de Genesis; las tarifas de Dictation corresponden a otro producto. & \textbf{C} & \href{https://www.invoxmedical.com/productos/generador-informes-medicos-inteligencia-artificial}{Genesis y contacto}, \href{https://www.invoxmedical.com/soluciones/inteligencia-artificial-para-hospitales}{soluciones hospitalarias}. \\
3. Heidi: Scribe /\allowbreak{} Remote & Plan gratuito; \textbf{tarifa de pago en España pendiente de verificación}. Referencia del blog estadounidense: Clinician \textbf{\$110/\allowbreak{}mes con anualidad}, Practice \textbf{\$180/\allowbreak{}mes con anualidad}, Remote \textbf{\$299,99}. & Planes/\allowbreak{}funciones públicos, pero los importes no resultaron accesibles en la versión consultada de la página española. Las cifras del blog corresponden a otra región y no constituyen una oferta para España. Remote se compra aparte; envío e impuestos pendientes de confirmación. El acuerdo empresarial remite a pedido y declara confidenciales sus precios. & \textbf{B} & \href{https://www.heidihealth.com/es-es/tarifas}{Tarifas España}, \href{https://www.heidihealth.com/en-us/blog/sully-ai-alternative}{referencia del propio proveedor}, \href{https://www.heidihealth.com/legal/remote-terms-of-sale}{Remote: venta}, \href{https://www.heidihealth.com/es-es/legal/acuerdo-marco-de-servicios}{acuerdo empresarial}. \\
4. Solventum: Fluency Align & \textbf{No localizado.} & Demostración, prueba y propuesta comercial; no tarifa pública del paquete Align, volumen incluido o implantación. Fluency Direct está incluido en la experiencia anunciada, pero no permite deducir precio de Align. & \textbf{C} & \href{https://www.solventum.com/es-es/home/health-information-technology/solutions/fluency-align/}{Align España}, \href{https://www.solventum.com/en-us/home/health-information-technology/solutions/fluency-align/}{paquete Align}. \\
5. Tandem & Free; Pro \textbf{125 €/\allowbreak{}usuario/\allowbreak{}mes, facturado anualmente}: \textbf{1.500 €/\allowbreak{}usuario/\allowbreak{}año}, importe anual calculado. & Free: notas ilimitadas, dos plantillas propias y acciones Pro limitadas. Pro: notas/\allowbreak{}acciones ilimitadas. Enterprise por usuario/\allowbreak{}mes, importe personalizado, incluye integración a medida en su descripción. Precio mensual sin compromiso anual y tratamiento del IVA pendientes de verificación. La integración con HCIS permanece pendiente. & \textbf{A: Pro; B: Enterprise} & \href{https://tandemhealth.ai/es/producto/precios}{Precios España}, \href{https://tandemhealth.ai/es/producto/integraciones/hp-his-hp-hcis}{estado HCIS}. \\
6. Microsoft: Dragon Copilot & Flex: \textbf{0,25 USD por sesión asistida por IA}, además de licencia base cuyo importe no se localizó. & Desde 04/\allowbreak{}05/\allowbreak{}2026: 25 unidades \ensuremath{\times} 0,01 USD. Modalidad alternativa por usuario con funcionalidades incluidas. En Flex una sesión puede contener varias acciones; un episodio hospitalario puede requerir varias sesiones. Requiere Azure vinculado; no equivale a 0,25 USD por paciente. & \textbf{B} & \href{https://www.microsoft.com/licensing/guidance/Dragon-Copilot}{Guía oficial de licencias y consumo}. \\
7. Nabla & \textbf{No localizado en fuentes actuales revisadas.} & Condiciones remiten a web/\allowbreak{}portal de suscripción; importe, cuota incluida y cargos adicionales pendientes de verificación. No se utilizan tarifas de comparadores como confirmación del precio vigente. La ruta \texttt{/\allowbreak{}pricing} no devolvió página accesible. & \textbf{C: evidencia disponible} & \href{https://www.nabla.com/}{Producto}, \href{https://www.nabla.com/docs/terms-baa}{condiciones y tarifas}, \href{https://www.nabla.com/pricing}{ruta intentada}. \\
8. Abridge & \textbf{No localizado.} & Contacto comercial; el calculador de ahorro no publica una tarifa del escriba. Unidad, mínimos y servicios de integración pendientes de propuesta. & \textbf{C} & \href{https://www.abridge.com/terms}{Contacto y condiciones}, \href{https://www.abridge.com/calculator}{calculador}. \\
9. Ambience & \textbf{No localizado.} & Anuncia honorarios ligados a resultados, pero no importes ni fórmula pública suficiente para presupuestar. Implantación, módulos y condiciones del compromiso pendientes. & \textbf{B: modelo parcial} & \href{https://www.ambiencehealthcare.com/}{Producto y honorarios ligados a resultados}, \href{https://www.ambiencehealthcare.com/resources/klas-ambient-ai-scribe-roi-validations}{resultados KLAS}. \\
10. Suki & \textbf{No localizado como tarifa general vigente.} & Pedido y servicios profesionales separados en contrato. Caso McLeod: pago por encuentro con tope, sin cuantías publicadas; no es necesariamente la oferta estándar. & \textbf{B: modelo parcial} & \href{https://www.suki.ai/terms-of-service/}{Condiciones, §§2.8 y 3}, \href{https://www.suki.ai/blog/when-ai-actually-pays-off-what-three-health-systems-found-when-they-let-klas-do-the-math/}{modelo contratado por McLeod}. \\
11. DeepScribe & \textbf{No localizado.} & Contratación específica; integración depende del pedido. Las cifras genéricas de su blog sobre el mercado no son una tarifa propia. & \textbf{C} & \href{https://www.deepscribe.ai/terms}{Condiciones}, \href{https://www.deepscribe.ai/resources/best-ai-medical-scribes}{guía comercial}, \href{https://www.deepscribe.ai/pricing}{ruta de precios intentada}. \\
12. Commure Scribe & \textbf{\$89/\allowbreak{}mes} o \textbf{\$59/\allowbreak{}mes facturado anualmente}; promoción publicada de \textbf{\$45 el primer mes}. & Scribe individual ilimitado. La oferta hospitalaria Ambient AI requiere una propuesta distinta. La página comercial abrevia los precios; el blog propio aclara mensual/\allowbreak{}anual. La promoción y el alcance de integración requieren confirmación contractual. & \textbf{A: individual} & \href{https://scribe.commure.com/}{Scribe}, \href{https://www.commure.com/blog-scribe/sunoh-ai-review}{periodicidad aclarada por Commure}, \href{https://www.commure.com/ambient-ai}{oferta hospitalaria}. \\
13. Corti: API & Voz a texto: \textbf{\$0,0065/\allowbreak{}min}. Generación de notas: \textbf{\$4/\allowbreak{}millón de tokens de entrada + \$16/\allowbreak{}millón de salida}. & Se suman audio y generación. Créditos iniciales de \$50; créditos comprados caducan al año. Acceleration Pack: \$1.000/\allowbreak{}mes con créditos y cinco horas de soporte de desarrollo. Aplicación, conector HCIS y operación son costes adicionales a estimar. & \textbf{A: API} & \href{https://corti.ai/pricing}{Precios y reglas de créditos}. \\
14. Oracle: Clinical AI Agent & \textbf{No localizado.} & Contacto comercial; la página presenta conjuntamente capacidades disponibles y previstas. No hay importe ni base suficiente para calcular el escriba con licencias EHR, implantación y soporte. & \textbf{C} & \href{https://www.oracle.com/health/clinical-suite/clinical-ai-agent/}{Producto, contacto y condiciones de disponibilidad}. \\
15. Tortus & \textbf{Sin importe público en las fuentes consultadas.} & Hospitales: licencia anual por clínico o contrato por volumen; también describe acuerdos ligados a resultados. SDK: por hora de audio, descuentos por volumen y 100 horas de prueba. Modalidades de facturación identificadas; importes pendientes de propuesta. & \textbf{B} & \href{https://tortus.ai/hospitals}{Hospitales, sección comercial}, \href{https://tortus.ai/partners}{SDK y cobro por audio}. \\
16. Lyrebird & UK: \textbf{£59/\allowbreak{}mes} o \textbf{£480/\allowbreak{}año}. EE. UU.: \textbf{89 USD/\allowbreak{}mes} o \textbf{828 USD/\allowbreak{}año}. & Lite gratuito: 50 acciones/\allowbreak{}mes. La ayuda indica que otros países, por tanto España, reciben precios EE. UU.; organizaciones por oferta. Equivalentes anuales: £40 o 69 USD/\allowbreak{}mes. Los planes australianos de clínica no son equivalentes al escriba individual. & \textbf{A: individual} & \href{https://help.lyrebirdhealth.com/en/articles/10066768-pricing-how-do-i-pay}{Guía regional actualizada 07/\allowbreak{}08/\allowbreak{}2026}, \href{https://help.lyrebirdhealth.com/en/articles/14332033-subscription-faqs}{suscripciones}. \\
17. Tali & \textbf{Discrepancias entre tarifas publicadas:} Pro muestra \textbf{150 CAD/\allowbreak{}mes} y otros bloques \textbf{100/\allowbreak{}68 CAD/\allowbreak{}mes}. & Las discrepancias impiden establecer una tarifa vigente. Publica anual/\allowbreak{}mensual, Free con cinco escribas/\allowbreak{}mes y Pro ilimitado; hay versiones duplicadas y enlaces de staging. Plan, moneda y periodo de facturación pendientes de confirmación contractual. & \textbf{B: ambigua} & \href{https://tali.ai/pricing}{Precios}, \href{https://help.tali.ai/en/articles/8562936-which-plan-should-i-choose}{guía de planes}. \\
18. Scribeberry & Pro \textbf{\$99/\allowbreak{}mes}; Enterprise \textbf{desde \$79/\allowbreak{}usuario/\allowbreak{}mes, mínimo cinco}; moneda no identificada en la tarjeta. & Free: tres días ilimitados, luego 20 usos/\allowbreak{}mes. Pro ilimitado. API separada: entrada \$0,008/\allowbreak{}1.000 tokens; salida \$0,032/\allowbreak{}1.000; audio archivo \$0,204/\allowbreak{}h o directo \$0,85/\allowbreak{}h. La agregación de costes de suscripción y API depende de la modalidad contratada. & \textbf{B: moneda/\allowbreak{}alcance; API detallada} & \href{https://scribeberry.com/pricing}{Planes}, \href{https://console.scribeberry.com/docs/pricing}{API y facturación}. \\
19. Quirónsalud: Scribe & \textbf{N/\allowbreak{}A.} & Proyecto interno; no se localizó oferta de venta o licencia para terceros. La ausencia de oferta comercial impide asignar un precio comparable. & \textbf{N/\allowbreak{}A} & \href{https://www.quironsalud.com/es/comunicacion/actualidad/quironsalud-redisena-modelo-asistencial-reafirmandose-human}{Anuncio del proyecto}. \\
20. Freed & Starter \textbf{\$39/\allowbreak{}mes}; Core \textbf{\$79/\allowbreak{}mes}; Premier \textbf{\$119/\allowbreak{}mes} o \textbf{\$1.248/\allowbreak{}año}. & Un clínico; Starter 40 notas/\allowbreak{}mes, Core ilimitado, integración EHR en Premier para sistemas compatibles. Prueba siete días. Grupos desde dos clínicos, oferta según tamaño. La documentación de ayuda aclara los bloques de precios duplicados de la página comercial. & \textbf{A} & \href{https://help.getfreed.ai/en/articles/9242518-freed-pricing-discounts-and-referral-program}{Guía de cobro, límites y cancelación}, \href{https://www.getfreed.ai/pricing}{precios}. \\
21. 33n: CLEARnotes & \textbf{No localizado.} & Remite a contacto para planes y precio; no unidad facturable pública identificada. & \textbf{C} & \href{https://clearnotes.co.uk/contact-us/}{Contacto sobre precios}, \href{https://clearnotes.co.uk/product/}{producto}. \\
22. Accuro Assistant & Healthcare: \textbf{importe no localizado}. Referencia \textbf{social}, no tarifa hospitalaria: \textbf{£0,06/\allowbreak{}min}. & Healthcare anuncia pago por uso sin suscripción. La oferta social sí separa audio y secretaría (£2/\allowbreak{}tarea); esos importes no se extrapolan al escriba NHS. & \textbf{B} & \href{https://accuro.co.uk/accuro-assistant-the-ambient-scribe-built-for-nhs-clinicians/}{Assistant NHS}, \href{https://accuro.co.uk/accuro-ambient-social-scribe-ai-uk/}{oferta social y tarifas}. \\
23. Accurx Scribe & \textbf{£0,50 /\allowbreak{} £0,47 /\allowbreak{} £0,43 /\allowbreak{} £0,40 por paciente adscrito y año}, según tamaño del cupo. & Tramos: 1–14.999; 15.000–89.999; 90.000–799.999; \ensuremath{\geq}800.000 pacientes. Escribas ilimitados, códigos, integración EPR y formación incluidos. Se cobra por población registrada, no por consulta grabada. Oferta de primaria UK. & \textbf{A} & \href{https://www.accurx.com/scribe-pricing}{Tarifa y tramos}. \\
24. Anathem & \textbf{£60/\allowbreak{}mes} o \textbf{£540/\allowbreak{}año}, sin IVA. & Clínico individual de salud mental; transcripciones/\allowbreak{}notas/\allowbreak{}documentos ilimitados, prueba 14 días. Organizaciones y plantillas específicas por propuesta. Anual equivalente a £45/\allowbreak{}mes. & \textbf{A: individual} & \href{https://anathem.ai/pricing}{Precios}. \\
25. Anima: Annie & \textbf{No localizado.} & Producto integrado en plataforma de primaria; importe y unidad de cobro de Scribe no publicados en la ficha consultada. & \textbf{C} & \href{https://www.animahealth.com/products/scribe}{Clinical Scribe}. \\
26. Lexacom & \textbf{Importe ambiental no localizado.} & Echo publica licencia anual o perpetua, precio según usuarios; es dictado. No demuestra precio o inclusión del escriba ambiental. & \textbf{C: ambiental} & \href{https://www.lexacom.co.uk/}{Producto}, \href{https://www.lexacom.co.uk/medical-speech-recognition-software/}{precio/\allowbreak{}modelo de Echo}. \\
27. Beam: Magic Notes & G-Cloud, junio 2026: Core \textbf{£30.000/\allowbreak{}año, hasta 6.000 h}; hasta \textbf{£225.000/\allowbreak{}año, 150.000 h}, con tramos intermedios. & Usuarios ilimitados; pago anticipado por bolsa anual. Flex usa créditos/\allowbreak{}módulos con otra escala. Publica implantación, formación y soporte; Talk puede añadir £2.500. Costes por exceso de uso y condiciones para España pendientes de confirmación. El cálculo excluye el descuento del piloto, vencido en mayo. & \textbf{A: catálogo UK} & \href{https://assets.applytosupply.digitalmarketplace.service.gov.uk/g-cloud-14/documents/717297/985901463553519-pricing-document-2026-06-04-1149.pdf}{Tarifario completo junio 2026}. \\
28. Cerebric & \textbf{No localizado.} & Oferta modular por contacto; sin tarifa del escriba. Su comparador publica precios de competidores, que no son precios de Cerebric. & \textbf{C} & \href{https://cerebric.io/product}{Producto}, \href{https://cerebric.io/research/choosing-software-uk-neurodevelopmental-clinic}{comparador propio}. \\
29. ClinicLetter.ai & \textbf{No localizado.} & Sitio con extracción limitada; no se encontró tarifa o unidad facturable en contenido público indexado. No permite concluir que el precio se oculte deliberadamente. & \textbf{No verificable} & \href{https://clinicletter.ai/}{Sitio consultado}. \\
30. Dentistry Dashboard & UK/\allowbreak{}UE Dentist Zone: \textbf{£40,99/\allowbreak{}dentista/\allowbreak{}mes}, una licencia. & AI Notes ilimitado, prueba 14 días con tarjeta. Practice Hub £99/\allowbreak{}local/\allowbreak{}mes; paquete Hub + Dentist Zone £130/\allowbreak{}mes. Licencias de asociados adicionales; tratamiento fiscal pendiente de confirmación. Oferta dental, no hospital general. & \textbf{A: licencia dental} & \href{https://dentistrydashboard.com/pricing}{Precios UK/\allowbreak{}UE}. \\
31. Dictate IT: AVT & \textbf{No localizado.} & Prueba o demostración; no tarifa ambiental ni volumen incluido públicos en la ficha revisada. Las tarifas de T-Pro Dictate corresponden a otra empresa. & \textbf{C} & \href{https://dictateit.com/solutions/ambient-listening/}{Ambient Listening}, \href{https://dictateit.com/}{producto}. \\
32. Docwise & \textbf{No localizado.} & Anuncia precio por volumen para grupos, sin escala o unidad facturable. & \textbf{C} & \href{https://www.docwise.co.uk/}{Producto y oferta por volumen}. \\
33. eConsult: Hi Scribe & \textbf{No localizado.} & Acceso/\allowbreak{}prueba mediante equipo comercial; no cuota de Hi Scribe o condiciones de inclusión en eConsult verificadas. & \textbf{C} & \href{https://help.econsult.health/en/articles/11331697-hi-scribe-smarter-notes-instantly}{Guía de acceso y prueba}, \href{https://econsult.net/hi-scribe-privacy-policy}{producto específico}. \\
34. Epro: AVT & \textbf{Importe ambiental no localizado.} & G-Cloud muestra £30/\allowbreak{}licencia para \textbf{AI Speech and Dictation}, sin demostrar tarifa del escriba ambiental actual. Implantación y módulos pendientes de oferta. & \textbf{C: ambiental} & \href{https://epro.com/}{Oferta Epro}, \href{https://www.applytosupply.digitalmarketplace.service.gov.uk/g-cloud/services/619277374438387}{ficha de dictado, alcance distinto}. \\
35. G2 Speech: SpeechAmbient & \textbf{No localizado.} & Demostración y contacto comercial; sin unidad, cuota o mínimos públicos identificados. & \textbf{C} & \href{https://www.g2speech.com/speechambient}{Ficha y contacto}. \\
36. HealthOrbit AI & Pro \textbf{£40/\allowbreak{}clínico/\allowbreak{}mes}, o \textbf{£33/\allowbreak{}clínico/\allowbreak{}mes con anualidad}. & Minutos ilimitados. Enterprise personalizado; intérprete, Ask AI y otros módulos son de pago sin precio en esa tabla. Total anual pendiente de confirmación: «dos meses gratis» no coincide exactamente con 33 \ensuremath{\times} 12 por redondeo. & \textbf{B: extras y redondeo} & \href{https://healthorbit.ai/pricing-plans/}{Precios y módulos}. \\
37. Medhasyn: Clinical Companion & \textbf{No localizado.} & Información pública limitada, sin planes, tarifas ni unidad facturable identificados. & \textbf{C} & \href{https://www.medhasyn.com/}{Sitio oficial}. \\
38. Optum: EMIS Scribe & \textbf{No localizado.} & Importe remitido a propuesta; la tarifa de T-Pro no es directamente atribuible al producto EMIS Scribe. & \textbf{C} & \href{https://www.emishealth.com/news-insights/march-2026-newsletter}{Producto}, \href{https://www.emishealth.com/terms-conditions}{condiciones}. \\
39. Pungo: JoyNotes & Publicita \textbf{acceso anticipado gratuito}; precio posterior no localizado. & Plazo de acceso, límites y continuidad de la gratuidad pendientes de verificación. Organizaciones mediante contacto. & \textbf{B: gratuidad promocional} & \href{https://www.thejoyapp.com/products/joynotes}{JoyNotes y acceso anticipado}. \\
40. Scribetech: Augnito Omni & \textbf{No localizado.} & Anuncia consultas, minutos e informes ilimitados; no publica en las fuentes consultadas cuota, mínimos o coste de integración. & \textbf{B: límites conocidos} & \href{https://scribetech.co.uk/augnito-omni-nhs-ai-scribe}{Oferta Omni}, \href{https://scribetech.co.uk/press-releases/scribetech-uk-to-launch-augnito-omni-ai-in-2025-0}{declaración de cobro ilimitado}. \\
41. T-Pro: Scribe /\allowbreak{} Copilot & \textbf{Referencia G-Cloud mayo 2024:} ambiental Standalone \textbf{£2.700/\allowbreak{}usuario/\allowbreak{}año}; Enterprise \textbf{£3.900/\allowbreak{}usuario/\allowbreak{}año}, sin impuestos. & Equivalentes £225/\allowbreak{}£325 al mes. PDF separa implantación, formación e interfaces como servicios adicionales; alojamiento/\allowbreak{}mantenimiento/\allowbreak{}soporte incluidos. \textbf{Vigencia comercial 2026 pendiente}; los precios de dictado del mismo PDF son menores y no equivalentes. & \textbf{B: tarifa antigua detallada} & \href{https://assets.applytosupply.digitalmarketplace.service.gov.uk/g-cloud-14/documents/704267/310198615907338-pricing-document-2024-05-07-1019.pdf}{PDF, pp. 3–4}, \href{https://info.tpro.io/partner\_with\_tpro}{modelo actual de distribución}. \\
42. X-on: Surgery Intellect & \textbf{No localizado.} & Oferta combinada con Surgery Connect/\allowbreak{}Tortus; desglose de costes del escriba, telefonía y servicios adicionales pendiente de verificación. & \textbf{C} & \href{https://www.x-on.co.uk/surgery-intellect/}{Producto}. \\
43. AWS: HealthScribe & \textbf{0,001667 USD/\allowbreak{}segundo}, aproximadamente \textbf{0,10 USD/\allowbreak{}minuto}. & Pago por audio procesado, incrementos de un segundo y mínimo de 15 segundos/\allowbreak{}petición. Incluye transcripción y resumen. Prueba: 300 min/\allowbreak{}mes durante los dos primeros meses. Aproximadamente 1,50 USD/\allowbreak{}consulta de 15 min, solo API; aplicación, almacenamiento, conector HCIS y operación adicionales. & \textbf{A: API} & \href{https://aws.amazon.com/healthscribe/pricing/}{Tarifa, alcance y prueba}, \href{https://aws.amazon.com/es/healthscribe/}{producto para desarrolladores}. \\
\end{longtable}

\textbf{Estimaciones ilustrativas a partir de tarifas públicas:} 10 usuarios Tandem Pro con anualidad = \textbf{15.000 €/\allowbreak{}año} antes de resolver impuestos e integración. Accurx para un cupo de 10.000 pacientes = \textbf{5.000 £/\allowbreak{}año}, independientemente del número de consultas dentro de esa oferta. Dragon Flex con 4.000 sesiones = \textbf{1.000 USD de componente variable}, más licencias y demás conceptos aplicables. Fuentes: filas 5, 23 y 6.

La comparación de propuestas para HCIS requiere un volumen de actividad común y el desglose de \textbf{licencias + consumo + conector/\allowbreak{}implantación + mantenimiento del conector + soporte/\allowbreak{}formación + hardware + impuestos}. Deben especificarse el hecho que activa la facturación (usuario nominal/\allowbreak{}activo, consulta, regeneración, minuto o token), mínimos, excesos, renovación y coste de salida. Con los precios públicos actuales \textbf{no se puede determinar la opción de menor coste total para la implantación del hospital}.
\endgroup
\end{landscape}

% !TEX root = ../main.tex
\clearpage
\section{Adquisición comercial y desarrollo propio: consideraciones técnicas}
\label{app:construir}
El análisis de una alternativa de desarrollo propio comprende componentes de procesamiento, recursos de evaluación y requisitos de integración. La documentación disponible incluye los modelos multimodales de pesos abiertos Gemma 4 y Qwen3-Omni, así como los recursos de evaluación PriMock57 y TEME-v1~\cite{J4955P47,P9P5U4KE,G2L7QBIV,CMRQCQTY}. Estos elementos pueden formar parte de una solución, pero su disponibilidad no equivale a disponer de un escriba clínico validado, certificado e integrado con HCIS.

Y-KNOT ilustra la normalización de documentos, la definición de eventos de activación y la devolución de borradores para aprobación dentro de la historia clínica electrónica. Su entrada se compone de registros existentes y el código no es público; estas características limitan su consideración como alternativa abierta a la captura ambiental~\cite{QQCRZ7VX}.

La comparación técnica debe incluir licencias y versiones, procesamiento local o remoto, capacidad para audio en español, trazabilidad de errores, conectores, autenticación, mantenimiento, actualización y coste total. Tanto una arquitectura de reconocimiento automático del habla y generación de texto (ASR+LLM) como un modelo multimodal requieren evaluar la nota final y el flujo completo. La selección de información y la revisión clínica siguen siendo funciones necesarias, con independencia de su distribución entre componentes.

\pendiente{Este análisis preliminar requiere una selección sistemática de aplicaciones abiertas y una estimación de infraestructura, operación y recursos de desarrollo.}

\printbibliography[heading=subbibliography,title={Referencias técnicas},keyword=tecnica]

% !TEX root = ../main.tex
\clearpage
\section{Opciones metodológicas para la evaluación}
\label{app:evaluacion-piloto}
\textbf{Propuesta para discusión.} Este apéndice concreta alternativas para la secuencia de la sección~\ref{sec:piloto}. La elección del diseño dependerá del objetivo institucional, los servicios disponibles y los recursos; no constituye un protocolo aprobado.

\subsection{Calidad de las notas sobre casos comunes}
Se propone evaluar los productos con las mismas consultas y el mismo contexto clínico, separando los casos de ajuste de plantillas de los reservados para evaluación. Las diferencias entre carga de archivos, reproducción de audio y captura en directo deben documentarse. Los trabajos de Biro y Ha ofrecen procedimientos con consultas estandarizadas y revisión clínica que pueden adaptarse; sus muestras exploratorias no determinan el tamaño necesario para este hospital~\cite{NNWFT2JX,PRTISMYF}.

La referencia combinaría audio, contexto disponible y valoración independiente. La nota original aporta información, pero también puede contener omisiones. Se propone contar con dos revisores, ocultar la identidad del proveedor cuando sea viable y resolver los desacuerdos. La valoración distinguiría errores clínicos, omisiones, gravedad y preferencias de estilo; la similitud automática entre textos no sustituye el juicio clínico~\cite{NNWFT2JX,4IHZ7QFQ}.

\subsection{Diseño para evaluar el efecto asistencial}
Un banco de consultas permite comparar borradores, pero no estimar por sí solo el ahorro de trabajo real. Para este segundo objetivo se plantean tres alternativas:
\begin{small}
\begin{longtable}{L{.25\linewidth}L{.31\linewidth}L{.34\linewidth}}
\toprule
Diseño & Utilidad & Limitación principal \\
\midrule\endhead
Antes/después con los mismos profesionales & Explorar factibilidad y cambios iniciales. & Confunde el efecto del producto con aprendizaje y cambios de actividad. \\
Grupo contemporáneo con documentación habitual; asignación aleatoria si es viable & Estimar el efecto frente al flujo habitual. & Requiere muestra suficiente y equilibrio entre servicios y carga inicial. \\
Incorporación escalonada o periodos alternos & Acompañar un despliegue progresivo. & Requiere controlar calendario, aprendizaje y posibles efectos persistentes entre periodos. \\
\bottomrule
\end{longtable}
\end{small}
Estas alternativas son propuestas de diseño, no una clasificación validada de protocolos para escribas. El ensayo de UCLA aporta un ejemplo de comparación aleatorizada con dos productos y control~\cite{PR93N8FN}. Para evitar confundir producto y especialidad, convendría comparar dentro de servicios o equilibrar la asignación.

\clearpage
\subsection{Resultados y decisiones pendientes}
Como objetivo principal inicial se plantea el \textbf{tiempo total hasta completar la documentación}, condicionado a una calidad clínica aceptable. Si la prioridad fuera bienestar profesional, deberían ajustarse el resultado principal y el seguimiento. Las medidas complementarias incluirían:
\begin{itemize}
\item \textbf{Calidad:} errores relevantes, omisiones y errores residuales tras revisión, con denominadores y gravedad definidos.
\item \textbf{Trabajo real:} revisión, edición dentro y fuera del escriba, incorporación a HCIS y trabajo fuera de horario. Los registros de una aplicación pueden omitir correcciones realizadas en otra~\cite{C7S8N9NA}.
\item \textbf{Adopción y experiencia:} uso sobre consultas elegibles, rechazos, abandono, carga percibida y experiencia del paciente. Las expectativas previas y la experiencia durante el uso deben analizarse por separado~\cite{PR93N8FN,54NFKG36}.
\item \textbf{Operación y coste:} fallos, latencia, recuperación y coste por consulta con documentación completada.
\end{itemize}

Antes del inicio deberán definirse condiciones de tratamiento del conjunto de prueba, formación, versiones, criterios de suspensión y cambios que exijan repetir la evaluación. La muestra responderá al objetivo y a la precisión deseada, considerando la agrupación de consultas por profesional. La guía \fuente{https://discovery.ucl.ac.uk/10152188/3/Supplementary_appendix_1_clean.pdf}{DECIDE-AI} aporta criterios de reporte sobre evaluación inicial y factores humanos; su adaptación al escriba documental sería metodológica, pues su ámbito original es el apoyo a decisiones clínicas. La autorización para preparar datos de evaluación no equivale a autorizar su transferencia a proveedores~\cite{rgpd2016}.

\end{document}
